% Merged research report. Standalone: no external graphics, no BibTeX run.
\documentclass[11pt]{article}
\usepackage[T1]{fontenc}
\usepackage[utf8]{inputenc}
\usepackage{lmodern}
\usepackage[margin=1in,headheight=15pt]{geometry}
\usepackage{amsmath,amssymb,amsthm,mathtools}
\usepackage{microtype}
\usepackage{booktabs,array,longtable,tabularx}
\usepackage{xcolor}
\usepackage{enumitem}
\usepackage{fancyhdr}
\usepackage{needspace}
\usepackage{xurl}
\usepackage{hyperref}
\usepackage{bookmark}
\definecolor{ink}{HTML}{233A45}
\definecolor{olive}{HTML}{596547}
\hypersetup{colorlinks=true,linkcolor=ink,citecolor=olive,urlcolor=ink,
  breaklinks=true,
  pdftitle={Holonomic Rigidity for Entire Hahn Functions},
  pdfauthor={Merged research report prepared for Vladimir Reshetnikov},
  pdfsubject={Arbitrary-rank Hahn fields, strong summability, linear
              differential and q-difference equations, first-order
              algebraic differential equations}}
\urlstyle{same}
\pagestyle{fancy}
\fancyhf{}
\fancyhead[L]{\small Holonomic rigidity for entire Hahn functions}
\fancyhead[R]{\small A sharp order-unit criterion}
\fancyfoot[C]{\thepage}
\renewcommand{\headrulewidth}{0.3pt}
\setlength{\parskip}{3.5pt}
\setlength{\parindent}{1.15em}
\setlength{\emergencystretch}{2em}
\widowpenalty=10000\clubpenalty=10000
\setcounter{tocdepth}{1}
\numberwithin{equation}{section}
\allowdisplaybreaks[2]
\newtheorem{maintheorem}{Theorem}
\renewcommand{\themaintheorem}{\Alph{maintheorem}}
\newtheorem{theorem}{Theorem}[section]
\newtheorem{lemma}[theorem]{Lemma}
\newtheorem{proposition}[theorem]{Proposition}
\newtheorem{corollary}[theorem]{Corollary}
\theoremstyle{definition}
\newtheorem{definition}[theorem]{Definition}
\newtheorem{convention}[theorem]{Convention}
\newtheorem{example}[theorem]{Example}
\newtheorem{question}[theorem]{Question}
\theoremstyle{remark}
\newtheorem{remark}[theorem]{Remark}
\newcommand{\C}{\mathbb C}
\newcommand{\R}{\mathbb R}
\newcommand{\Q}{\mathbb Q}
\newcommand{\N}{\mathbb N}
\newcommand{\Z}{\mathbb Z}
\newcommand{\No}{\mathbf{No}}
\newcommand{\K}{K_\Gamma}
\newcommand{\E}{\mathcal E_\Gamma}
\newcommand{\supp}{\operatorname{supp}}
\newcommand{\lc}{\operatorname{lc}}
\newcommand{\res}{\operatorname{res}}
\newcommand{\cf}{\operatorname{cf}}
\newcommand{\Dom}{\operatorname{Dom}_{\mathrm{str}}}
\newcommand{\fall}[2]{(#1)_{\underline{#2}}}
\newcommand{\abs}[1]{\lvert#1\rvert}
\newcommand{\repo}{\texttt{VladimirReshetnikov/Surreal}}
% Nonlinear part (source 10): any characteristic-zero coefficient field k.
\newcommand{\KK}{\mathbb K}
\newcommand{\EK}{\mathcal E_\Gamma(k)}
\newcommand{\Sfam}{\mathcal S}
\newcommand{\wt}{\operatorname{wt}}
\newcommand{\inpoly}{\operatorname{in}}
\setlist[itemize]{leftmargin=1.5em,itemsep=2pt,topsep=4pt}
\setlist[enumerate]{leftmargin=1.8em,itemsep=3pt,topsep=4pt}

\begin{document}
\hypersetup{pageanchor=false}
\begin{titlepage}
\centering
\vspace*{0in}
{\color{olive}\large\sffamily SURREAL AND SURCOMPLEX RESEARCH}\par
\vspace{0.18in}
{\color{ink}\huge\bfseries Holonomic Rigidity\\[6pt] for Entire Hahn Functions}\par
\vspace{0.2in}
{\Large A sharp order-unit criterion for dilations,\\[4pt]
with uniform exclusion bounds,\\[4pt]
and first-order nonlinear rigidity}\par
\vspace{0.2in}
{\large Merged research report\quad\textperiodcentered\quad 22 September 2026}\par
\vspace{0.12in}
\begin{minipage}{0.94\textwidth}
\begin{abstract}
Fix a nonzero set-sized ordered abelian group $\Gamma$ and the complex Hahn
field $\K=\C((t^\Gamma))$. \emph{Entire} means strong Hahn summability of the
evaluation family at every point of this one field. We prove that every
strongly entire D-finite power series is a polynomial, in one or in finitely
many variables, and that a fixed differential operator already excludes every
nonterminating formal solution on an explicit exterior valuation region,
with a common degree bound.

For dilations there is an exact alternative. If $q\in\K^\times$ has infinite
multiplicative order, a nonpolynomial strongly entire solution of a nonzero
linear $q$-difference equation with polynomial coefficients exists if and only
if $v(q)\ne0$ and $\abs{v(q)}$ is an order unit of $\Gamma$. A nonzero
valuation is not enough: its integer multiples must be cofinal among all
scales. The proof covers units whose residues are roots of unity, and
coefficients and parameters with arbitrary well-ordered Hahn supports. A
partial theta series proves sharpness, and its exact strong evaluation domain
is computed for an arbitrary Hahn parameter. Quantitatively, the recurrence
supplies an exclusion bound that is attained on its own boundary, and a
coarsened exclusion bound valid for every noncofinal dilation. A mixed
differential--dilation extension is proved for nonzero noncofinal dilation
valuations.

The common mechanism is an elementary finite-step escape chain: bounded
valuation costs in a coefficient recurrence force a subsequence of evaluated
leading exponents that never increases. This contradicts strong summability
without replacing an arbitrary-rank valuation by a real-valued norm. In an
explicit surreal workspace with countable cofinality and no order unit, the
series $\sum_{n\ge0}\omega^{-\omega^n}z^n$ is strongly entire yet satisfies no
nonzero linear differential equation and no nonzero linear dilation equation
for any nontorsion parameter of that workspace.

A third source adds a nonlinear part over $k((t^\Gamma))$, for any field $k$
of characteristic zero: every strongly entire solution of a nonzero
first-order algebraic differential equation $P(z,f,D_zf)=0$ is a polynomial,
with an exclusion bound that must depend on the solution. Higher order stays
open outside a nondegenerate class.

$\Gamma$ is \emph{not} assumed divisible. Divisibility is hypothesised only
where it is used, in the three threshold constructions listed in
Section~\ref{hol:sub:divisibility}, and each of those statements prints the
hypothesis itself.
\end{abstract}
\end{minipage}
\vfill
\begin{minipage}{0.94\textwidth}
\small
\textbf{Status.} An AI-assisted research report; priority and independent
refereeing are not certified. Lean covers the positive-support word lemma,
escape mechanism, bounded-cost exclusion and formal exponential domain.
The main differential and dilation classifications, the new torsion
argument and the whole nonlinear part remain unformalized. The variable
derivative is not the Berarducci--Mantova derivation on surreal coefficients.
\end{minipage}
\end{titlepage}
\hypersetup{pageanchor=true}

\tableofcontents
\clearpage

\section{The question and the principal results}\label{hol:sec:intro}

\subsection{Why the question is scale sensitive}

A familiar way to specify a special function is to give a finite linear
differential equation. Another is to give a finite relation between its values
at $z,qz,q^2z,\ldots$. Over the complex numbers both descriptions produce
nonpolynomial entire functions. Over a surreal Hahn workspace, ``entire'' asks
for evaluation not only at ordinary complex constants but at every infinite
and infinitesimal scale available in that workspace. In a Hahn field an
infinite expression is legal only when its supports satisfy an order condition
and a finiteness condition, and an equation can force either condition to fail
at a sufficiently large argument. The question of this report is what a finite
linear equation can accommodate under that stronger global requirement.
Sections~\ref{hol:nl:sec:setting}--\ref{hol:nl:sec:consequences} ask the same
of a finite \emph{nonlinear} algebraic differential equation, and answer it in
order one.

The answer depends decisively on which operator is used, and for dilation
equations it depends further on whether the valuation of the particular
parameter reaches every scale of the value group.

\subsection{Conventions}\label{hol:sub:conventions}

We work with a fixed ordered abelian group $\Gamma$ and
\[
   \K=\C((t^\Gamma)).
\]
An element of $\K$ is a formal sum $a=\sum_{\gamma}a_\gamma t^\gamma$ with
$a_\gamma\in\C$ whose support $\supp(a)=\{\gamma:a_\gamma\ne0\}$ is a
well-ordered set. For $a\ne0$ put
\[
 v(a)=\min\supp(a),\qquad \lc(a)=a_{v(a)}\in\C^\times,
 \qquad \res(a)=\lc(a)\text{ when }v(a)=0,
\]
and put $v(0)=+\infty$ only where convenient; this symbol is never subtracted.
Products and sums are the usual Hahn operations, with
\begin{equation}\label{hol:eq:valuationlaws}
 v(ab)=v(a)+v(b),\qquad
 v(a+b)\ge\min\{v(a),v(b)\},
\end{equation}
with equality in the second law when the two valuations differ. The
coefficient field $\C$ is trivially valued, so every nonzero ordinary integer
and every factorial has valuation zero. Throughout, $\N=\{0,1,2,\ldots\}$, and
$\abs\lambda$ means $\max(\lambda,-\lambda)$ in the ordered group.

\begin{convention}[Standing hypothesis on the value group]
\label{hol:conv:group}
$\Gamma$ is a nonzero set-sized ordered abelian group. Divisibility is
\emph{not} part of this standing hypothesis. Every statement below that needs
divisibility says so in its own statement; those statements are collected in
Section~\ref{hol:sub:divisibility}.
\end{convention}

\begin{definition}[Strong summability, evaluation domain, entireness]
\label{hol:def:entire}
A set-indexed family $(b_i)_{i\in I}$ in $\K$ is \emph{strongly Hahn summable}
if $\bigcup_i\supp(b_i)$ is well ordered and each $\gamma\in\Gamma$ lies in the
supports of only finitely many members. Its sum is then defined
coefficientwise, by a finite sum in $\C$ at each exponent. For
$f=\sum_{n\ge0}a_nz^n\in\K[[z]]$ let $\Dom(f)$ be the set of $x\in\K$ for
which $(a_nx^n)_{n\ge0}$ is strongly summable, and let
\[
 \E=\{f:\Dom(f)=\K\}.
\]
In $d$ variables the family to be tested is $(a_\alpha x^\alpha)_{\alpha\in\N^d}$.
\end{definition}

The word \emph{entire} always means Definition~\ref{hol:def:entire} relative to
one fixed workspace. It does not mean that a complex coefficient function is
entire, it does not mean convergence in the fine topology inherited from the
full surreal class, and it is never a claim about all of $\No[i]$. The
definition also concerns the family \emph{before} cancellation between
different indices: one cannot first cancel a nonsummable family and then
declare it summable, and convergence of a numerical coefficient series does
not license infinitely many contributions at one Hahn exponent.

\begin{definition}[Order unit]\label{hol:def:orderunit}
A positive $\mu\in\Gamma$ is an \emph{order unit} if for every $\gamma\in\Gamma$
there is $N\in\N$ with $\abs\gamma\le N\mu$; equivalently, if the positive
integer multiples of $\mu$ are cofinal in $\Gamma$. Zero is not an order unit.
\end{definition}

\begin{definition}[The two operator classes]\label{hol:def:operators}
Let $D_z$ be the formal derivative in $z$, so $D_z(a)=0$ for every $a\in\K$,
and for $q\in\K^\times$ let $\sigma_qf(z)=f(qz)$. A series $f\in\K[[z]]$ is
\emph{D-finite over $\K(z)$} if
\begin{equation}\label{hol:eq:dlinear}
 \sum_{r=0}^R p_r(z)D_z^rf(z)=0,\qquad p_r\in\K[z],\qquad (p_0,\ldots,p_R)\ne0,
\end{equation}
and \emph{linearly $q$-difference finite} if
\begin{equation}\label{hol:eq:qlinear}
 \sum_{\ell=0}^r p_\ell(z)f(q^\ell z)=0,\qquad p_\ell\in\K[z],
 \qquad (p_0,\ldots,p_r)\ne0.
\end{equation}
These are formal identities in $\K[[z]]$, not assertions about limits of
partial sums. Rational coefficients give the same notions after denominators
are cleared.
\end{definition}

We say ``linearly $q$-difference finite'' rather than silently identifying
several meanings of ``$q$-holonomic''. In the sequence literature $q$ is often
an indeterminate with coefficients in $\Q(q)$; here $q$ is a \emph{fixed} Hahn
element, possibly with infinitely many terms. The differential terminology is
the usual characteristic-zero one~\cite{stanley}. \emph{Nontorsion} and
``of infinite multiplicative order'' are used interchangeably.

\subsection{The principal results}

\begin{maintheorem}[Differential rigidity, with a uniform exterior obstruction]
\label{hol:main:ode}
Let $\Gamma$ be as in Convention~\ref{hol:conv:group} and let
$0\ne L=\sum_{r,k}c_{r,k}z^kD_z^r$ be a nonzero linear differential operator
with coefficients in $\K[z]$. There are $B_L\in\Gamma$ with $B_L\ge0$ and
$N_L\in\N$, both determined by the finitely many coefficients of $L$, such
that every formal solution of $Lf=0$ that is strongly evaluable at even one
$x\ne0$ with $v(x)\le-B_L$ has degree less than $N_L$. Consequently
\[
 \{f\in\E:f\text{ is D-finite over }\K(z)\}=\K[z].
\]
\end{maintheorem}

\begin{maintheorem}[Sharp dilation criterion]\label{hol:main:q}
Let $\Gamma$ be as in Convention~\ref{hol:conv:group} and let $q\in\K^\times$
have infinite multiplicative order. The following are equivalent.
\begin{enumerate}[label=\textup{(\roman*)}]
\item There is a nonpolynomial $f\in\E$ that is linearly $q$-difference finite
over $\K(z)$.
\item $v(q)\ne0$ and $\abs{v(q)}$ is an order unit of $\Gamma$.
\end{enumerate}
When \textup{(ii)} fails, every fixed equation \eqref{hol:eq:qlinear} has a
uniform exterior obstruction and a uniform polynomial degree bound, exactly as
in Theorem~\ref{hol:main:ode}. When \textup{(ii)} holds, an explicit partial
theta series with parameter $q$ or $q^{-1}$ is a nonpolynomial entire
solution.
\end{maintheorem}

\begin{maintheorem}[Several-variable rigidity]\label{hol:main:multi}
Let $\Gamma$ be as in Convention~\ref{hol:conv:group} and let
$F\in\K[[z_1,\ldots,z_d]]$ be strongly Hahn summable at every point of $\K^d$.
If the mixed partial derivatives of $F$ span a finite-dimensional vector space
over $\K(z_1,\ldots,z_d)$, then $F$ is a polynomial.
\end{maintheorem}

The fourth principal theorem comes from the third source manuscript and is
proved in Sections~\ref{hol:nl:sec:setting}--\ref{hol:nl:sec:first}. It
allows an arbitrary coefficient field of characteristic zero; see
Convention~\ref{hol:nl:conv:field}, which also explains why that generality
is not transferred to Theorems~\ref{hol:main:ode}--\ref{hol:main:multi}.

\begin{maintheorem}[First-order nonlinear rigidity]\label{hol:nl:main:first}
Let $\Gamma$ be as in Convention~\ref{hol:conv:group}, let $k$ be any field of
characteristic zero, and let $\KK=k((t^\Gamma))$, with strong evaluation and
entireness defined as in Definition~\ref{hol:def:entire} with $k$ in place of
$\C$. For every nonzero $P\in\KK[z,Y_0,Y_1]$, every strongly entire formal
solution $f\in\KK[[z]]$ of
\[
 P(z,f,D_zf)=0
\]
is a polynomial. More strongly, if a nonpolynomial formal solution exists,
there is $\delta_0\in\Gamma_{>0}$, depending on the equation and on finitely
many coefficients of that solution, such that $t^{-\delta}\notin\Dom(f)$ for
every $\delta\ge\delta_0$. Consequently, for every nonpolynomial strongly
entire $f$, the series $f$ and $D_zf$ are algebraically independent over
$\KK(z)$.
\end{maintheorem}

Theorem~\ref{hol:nl:main:first} has no separant, irreducibility or
solved-form hypothesis on $P$. Its exclusion bound depends on the solution,
and Theorem~\ref{hol:nl:thm:nouniform} shows that this cannot be avoided,
in contrast with the operator-only bound of Theorem~\ref{hol:main:ode}. It
does not say that a nonpolynomial entire function is differentially
transcendental in all orders: independence of $f$ and $D_zf$ does not exclude
a relation involving $D_z^2f$. Question~\ref{hol:q:nonlinear} is re-scoped
accordingly.

The hypotheses separate two obstructions. Differential equations have bounded
valuation costs after passing to their coefficient recurrence. Dilation
equations can have costs increasing linearly with the index, but only along
the single scale $v(q)$; a scale that is not an order unit cannot reach all of
$\Gamma$, even if its multiples are unbounded within their own Archimedean
component. Strong entireness detects that failure.

Beyond the first three, the report proves: a mixed differential--dilation rigidity
theorem at every nonzero noncofinal dilation scale
(Theorem~\ref{hol:thm:mixed}); a recurrence exclusion bound that is
\emph{attained}, boundary included, by an explicit formal exponential
(Corollary~\ref{hol:cor:periodic} and
Proposition~\ref{hol:prop:expdomain}); a coarsened exclusion bound valid at
every nonzero noncofinal dilation valuation (Theorem~\ref{hol:thm:coarse}); and the exact
strong evaluation domain of the partial theta series for an arbitrary Hahn
parameter, including arbitrary tails (Theorem~\ref{hol:thm:thetadomain}).

\subsection{Where divisibility of \texorpdfstring{$\Gamma$}{Gamma} is used}
\label{hol:sub:divisibility}

The core arguments --- the escape mechanism, the orbit valuation lemmas, the
two conversions, the theta domain theorem, and all of
Theorems~\ref{hol:main:ode}--\ref{hol:main:multi} --- use only
Convention~\ref{hol:conv:group}. Divisibility of $\Gamma$ enters in exactly
three threshold constructions, with the hypothesis stated where it is used.
\begin{enumerate}[label=\textup{(\arabic*)}]
\item The refined recurrence threshold $B_{\mathrm{rec}}$ of
\eqref{hol:eq:radiusbound}, whose defining maximum divides a valuation
difference by the jump length $j$. Corollary~\ref{hol:cor:periodic} states it
under divisibility; Corollary~\ref{hol:cor:boundedcost} is the
divisibility-free predecessor, with a coarser threshold and the same
qualitative conclusion.
\item The refined coarsened threshold $\overline B_{\mathrm{rec}}$ of
\eqref{hol:eq:coarsebound}, for the same reason, in the quotient group.
Theorem~\ref{hol:thm:coarse} states the divisibility-free form and the
divisible refinement separately.
\item The worked sparse threshold in Example~\ref{hol:ex:sparse}, an instance
of (1).
\end{enumerate}
The torsion eigenspace equivalence in Proposition~\ref{hol:prop:torsion}
needs no divisibility: inward stability of strong evaluation replaces the
root-extraction argument of the source manuscripts.

The nonlinear part,
Sections~\ref{hol:nl:sec:setting}--\ref{hol:nl:sec:consequences}, uses no
divisibility anywhere: its only divisions are of ordinary integers in $\Q$
(the degree cutoff \eqref{hol:nl:eq:kappa}), and its exclusion bounds are
chosen from finitely many strict inequalities with positive integer
multipliers.

No theorem in this report states a conclusion under
Convention~\ref{hol:conv:group} while resting on a proof step that divides by
an integer in $\Gamma$. The division-free predecessors in (1) and (2) are
weaker in the constant only, not in the conclusion.

\subsection{Position relative to the collection and the literature}
\label{hol:sub:position}

\paragraph{Provenance.}
This report is built from three manuscripts, numbered as they arrived in the
collection. Sources~08 and~09 were written independently on 22 September 2026
and both prove the differential rigidity theorem and the order-unit dilation
criterion (Theorems~\ref{hol:main:ode} and~\ref{hol:main:q}); source~08 is
the base, because it assumes only $\Gamma\ne0$ while source~09 assumed
divisibility, and every proof taken from~09 was re-read for a hidden division
(Section~\ref{hol:sub:divisibility}). Together they supply
Sections~\ref{hol:sec:intro}--\ref{hol:sec:certificates}. Source~10, delivered
later the same day, supplies the nonlinear part,
Sections~\ref{hol:nl:sec:setting}--\ref{hol:nl:sec:consequences} and
Theorem~\ref{hol:nl:main:first}. The repository snapshot consulted by
sources~08 and~09 is
\begin{center}\small\texttt{a3124af79f66b8b9c196d76b4cbc5ac3938907c4},\end{center}
and the one consulted by source~10 is
\begin{center}\small\texttt{048b72cf7cbfc8ab246e4f73788c10460cb3f6e0}.\end{center}
At the second pin this report consisted of sources~08 and~09 alone and its text
was identical to the text into which source~10 was merged, so source~10's
references to ``Question~15.1'' and ``Theorems~A--C'' of this report were
accurate there; that question is now Question~\ref{hol:q:nonlinear}, re-scoped
in Section~\ref{hol:sec:status}. The choices the merge made (renamed symbols,
results printed once, the placement of the nonlinear part after the linear
theory) are recorded in Section~\ref{hol:nl:sub:conventions}.

\paragraph{The companion entire-function report.}
The collection already contains a merged report on cofinality, factorization
and scalar extension for entire Hahn functions at arbitrary
rank~\cite{repoentire}. It fixes the same class $\E$ over the same kind of
workspace, classifies the ring structure of $\E$, and computes the exact
strong evaluation domain after an enlargement of the value group. None of its
ring-theoretic conclusions is used to prove anything here: the present proofs
start from Hahn support arithmetic. Two points of contact matter.

First, we do not reprove the order-unit coarsening machinery. That report's
Lemma~9.1, labelled \texttt{ent:lem:coarsening}, states that if $h>0$ is an
order unit of $\Gamma$ there is an order-preserving homomorphism
$\pi:\Gamma\to\R$ with $\pi(h)=1$ whose kernel is the convex subgroup
$H=\{\gamma:n\abs\gamma<h\text{ for all }n\ge1\}$, so that
$\abs a_\pi=\exp(-\pi(v(a)))$ is a nontrivial real-valued non-Archimedean
absolute value inducing the intrinsic valuation topology. That is exactly the
sense in which sources~08 and~09 of this report use the phrase
``coarsening'', and we cite it rather than repeat it; see
Remark~\ref{hol:rem:coarsening}.

Second, and this is easy to get wrong:
\emph{the two order-unit dichotomies are not the same theorem} --- the
companion report's order unit decides Hermite interpolation and the B\'ezout
property of the ring $\E$ (its Theorem~9.3), whereas the order unit here
decides which annihilating operators a nonpolynomial member of $\E$ can
satisfy. Neither dichotomy implies the other, and the coincidence of the
invariant is not a coincidence of statement. In particular the existence of
nonpolynomial entire functions at countable cofinality, and the fact that
$\Gamma_\infty=\bigoplus_{j\ge0}\Q\omega^j$ has countable cofinality and no
order unit, are background from that report~\cite{repoentire}, not targets of
this one.

The collection's differential-equations report uses a derivation on surreal
\emph{scalars}~\cite{repodifferential}; here $D_z$ kills every scalar and
differentiates only the formal variable, and no theorem about that scalar
derivation is restated as ordinary D-finite rigidity. The rank-one analytic
report~\cite{reporankone} already contains the quadratic-exponent example
$\sum t^{n^2}z^n$; our partial theta witness is explicitly classical and is
not counted as a new function.

\paragraph{External literature.}
Stanley's D-finite/P-recursive correspondence is the standard algebraic
starting point~\cite{stanley}. We make the coefficient conversion explicit
because the valuations of the resulting recurrence, rather than ordinary
coefficient growth, drive the proof. Ramis treats growth of ordinary complex
entire solutions of linear $q$-difference equations and also studies
simultaneous differential and difference constraints~\cite{ramis}; the
arbitrary-rank strong summability condition here is a different global
condition, and none of its order-unit conclusions is inferred from an ordinary
complex growth estimate. Garoufalidis proves eventual quadratic
quasi-polynomial degree behaviour for $q$-holonomic sequences in a
rational-function and Puiseux setting~\cite{garoufalidis}; that is an
important antecedent to the idea that dilation equations constrain coefficient
scales, we do not claim to originate it, and we do not strengthen its
quadratic quasi-polynomial conclusion. Di Vizio's ultrametric
Maillet--Malgrange theorem treats Gevrey behaviour of formal solutions of
nonlinear ultrametric $q$-difference equations, including the valuation-zero
case with Diophantine conditions~\cite{divizio}; we make no nonlinear
generalization of it. The partial theta series belongs to the classical
partial-theta literature~\cite{andrewswarnaar}, and Higman's finite-word
theorem is a standard external ingredient in the positive-support monoid
lemma~\cite{higman}.

Terminology needs care in two further directions. Hahn meromorphic functions
in the analytic setting of M\"uller and Strohmaier involve a convergence
framework different from the formal strong evaluation used
here~\cite{mueller}. The Berarducci--Mantova derivation differentiates surreal
coefficients themselves and is not $D_z$~\cite{bm}. Work on operation-closed
surreal subfields gives context for set-sized workspaces~\cite{bg}, and the
Conway normal-form conventions used in Section~\ref{hol:sub:embedding} follow
current surreal field work~\cite{kks}; closure under a scalar derivation is
assumed in no theorem below.

The inspections of sources~08 and~09 were targeted: repository documentation,
the report catalogue, the entire-function report, and the primary sources
named above. Neither was a line-by-line audit of every repository manuscript
or an exhaustive bibliography search, and a repository content search for
\texttt{holonomic} returning no matches is not proof of absence. Source~10's
repository and literature comparison, equally targeted, is recorded in
Section~\ref{hol:nl:sub:predecessors}. Section~\ref{hol:sec:status} records
exactly what is and is not claimed.

\section{Hahn support and what evaluation means}\label{hol:sec:support}

\subsection{Support calculus}

\begin{lemma}[Hahn--Neumann support calculus and the positive monoid]
\label{hol:lem:support}
Let $A,B$ be well-ordered subsets of an ordered abelian group.
\begin{enumerate}[label=\textup{(\arabic*)}]
\item $A+B$ is well ordered, and each $\gamma$ has only finitely many
representations $\gamma=a+b$ with $(a,b)\in A\times B$.
\item If $S$ is a well-ordered subset of $\Gamma_{>0}$, the additive monoid
\[
 M(S)=\{s_1+\cdots+s_k:k\ge0,\ s_i\in S\}
\]
is well ordered, and each of its elements has only finitely many
representations as a finite word in letters from $S$.
\item Multiplication by a fixed Hahn element preserves strong summability, and
multiplication by a nonzero one preserves and reflects it.
\end{enumerate}
\end{lemma}

\begin{proof}
An infinite sequence in a well-ordered set has an infinite nondecreasing
subsequence. Given a strictly decreasing sequence $a_n+b_n$, pass first to a
subsequence with $a_n$ nondecreasing and then to one with $b_n$ nondecreasing,
a contradiction. If infinitely many distinct pairs sum to a fixed $\gamma$,
the distinct first coordinates contain an increasing sequence, forcing a
decreasing sequence of second coordinates. This is part~(1).

For part~(2), order finite words over $S$ by subsequence embedding with
coordinatewise increase. Higman's theorem says that words over a
well-quasi-ordered alphabet are well quasi ordered~\cite{higman}. An embedding
cannot decrease the sum, and increases it strictly unless the two words are
identical, since inserted letters are positive and a strict coordinate
increase is positive. A descending sequence of sums therefore gives a bad
sequence of words, and so do infinitely many distinct words with one fixed
sum. Both are impossible.

For part~(3), let $A$ be the union of the supports of a strongly summable
family and $B$ the support of the multiplier. The new supports lie in $A+B$.
Part~(1) leaves finitely many admissible pairs at a specified exponent, and
for each first coordinate only finitely many family members contribute. For a
nonzero multiplier apply the same to its inverse.
\end{proof}

In part~(2) words carry their order; a fixed finite multiset has only finitely
many permutations, so the commutative monoid version follows. The argument
does not assume that $S$ is bounded below by an order unit, and remains valid
at arbitrary valuation rank.

\begin{lemma}[The nonincreasing leading-exponent obstruction]
\label{hol:lem:nonincreasing}
A family of nonzero Hahn elements possessing an infinite subsequence whose
leading exponents are \emph{nonincreasing} is not strongly summable.
\end{lemma}

\begin{proof}
If that sequence of leading exponents takes infinitely many distinct values it
has an infinite strictly decreasing subsequence, so the union of the supports
is not well ordered. Otherwise it is eventually constant, and one exponent
lies in the supports of infinitely many members, so local finiteness fails.
\end{proof}

The second alternative is essential and is what makes the exclusion boundary
below a \emph{closed} region. Merely excluding a descending support does not
establish strong summability: the family $(1\cdot z^n)_{n\ge0}$ at $z=1$ has
support union $\{0\}$, yet every summand contributes at exponent zero.

\begin{lemma}[Two sufficient certificates]\label{hol:lem:certificate}
Let $(b_n)_{n\ge0}$ be a family in $\K$.
\begin{enumerate}[label=\textup{(\alph*)}]
\item If $v(b_n)$ tends cofinally to $+\infty$, meaning that for each
$\beta\in\Gamma$ all but finitely many nonzero $b_n$ satisfy $v(b_n)>\beta$,
then $(b_n)$ is strongly summable.
\item If $\supp(b_n)\subseteq h_n+M(S)$ for a well-ordered
$S\subseteq\Gamma_{>0}$, where the set of values of $h_n$ is well ordered and
each value is taken by only finitely many indices, then $(b_n)$ is strongly
summable.
\end{enumerate}
Moreover, grouping a strongly summable family into finite blocks preserves
strong summability and its sum.
\end{lemma}

\begin{proof}
For (a) fix $\beta$. Only finitely many supports meet $(-\infty,\beta]$, and a
finite union of well-ordered subsets of a linear order is well ordered. A
descending sequence in the full union would lie below its first member, hence
inside such a finite union. The same finiteness bounds the multiplicity of an
exponent. For (b) use parts~(1) and~(2) of Lemma~\ref{hol:lem:support}: a
fixed target exponent admits finitely many pairs $(h,m)\in\{h_n\}\times M(S)$
with $h+m$ equal to it, and finitely many indices carry each such $h$; actual
supports can only be smaller than the indicated sets. For the last assertion,
each block support is contained in the union of the old supports and the
blocks contributing at any exponent form a finite set, while coefficientwise
associativity is finite.
\end{proof}

Criterion~(a) and criterion~(b) are genuinely different tools and both are
used below. Criterion~(a) is the shorter route whenever the term valuations
really do escape every bound. Criterion~(b) is what survives at the boundary,
where the term valuations increase strictly but remain trapped below a fixed
element of $\Gamma$; Theorem~\ref{hol:thm:thetadomain} is exactly such a case,
and no leading-valuation argument alone can decide it.

\begin{lemma}[Inward stability of strong evaluation]\label{hol:lem:inward}
Let $f=\sum_na_nz^n\in\K[[z]]$, and let $x,y\in\K^\times$ with
$v(y)\ge v(x)$. If $x\in\Dom(f)$, then $y\in\Dom(f)$.
Thus membership of a nonzero argument in $\Dom(f)$ depends only on its
valuation, and entireness can be tested on monomial arguments.
\end{lemma}

\begin{proof}
Put $b_n=a_nx^n$ and $u=y/x$, so $(b_n)$ is strongly summable and
$v(u)=\eta\ge0$. Write $u=c\,t^\eta(1+w)$, where $c\in\C^\times$
and $w$ is zero or has positive valuation. Let $S$ be $\supp(w)$,
with $\eta$ adjoined when $\eta>0$, and put $M=M(S)$.
This is a well-ordered positive alphabet, so Lemma~\ref{hol:lem:support}
makes $M$ well ordered. Every finite power $u^n$ has support in $M$.

Let $A=\bigcup_n\supp(b_n)$. The supports of
$a_ny^n=b_nu^n$ lie in $A+M$, a well-ordered set.
For a fixed exponent $\gamma$ there are only finitely many pairs
$(\alpha,\mu)\in A\times M$ with $\alpha+\mu=\gamma$.
For each such $\alpha$, strong summability of $(b_n)$ leaves only
finitely many indices $n$. Hence $\gamma$ occurs in only finitely many
supports of $b_nu^n$, proving strong summability. Equal valuations allow
the argument in both directions; compare an arbitrary nonzero $y$ with
$t^{v(y)}$ to obtain the monomial test. Evaluation at zero is always finite.
\end{proof}

The statement fixes the full coefficient series $a_n$. It does not say
that their leading valuations alone determine the domain: higher
coefficient tails can still matter at a boundary scale.

\subsection{Positive valuation is not cofinal decay}

For $\beta>0$ write
\begin{equation}\label{hol:eq:delta}
 \Delta_\beta=\{\gamma\in\Gamma:\abs\gamma\le N\beta\text{ for some }N\in\N\},
\end{equation}
the smallest convex subgroup containing $\beta$. It is all of $\Gamma$ exactly
when $\beta$ is an order unit. If it is proper, there is $\gamma>0$ outside
it, and then
\begin{equation}\label{hol:eq:dominates}
 \gamma>N\beta\qquad\text{for every ordinary }N\in\N,
\end{equation}
since a positive element bounded by some $N\beta$ would lie in $\Delta_\beta$.

\begin{remark}[Sufficiency is not necessity]\label{hol:rem:positive}
Let $0<\epsilon<H$ in $\Gamma$ with $n\epsilon<H$ for every ordinary integer
$n$; for instance $\Gamma=\Q H\oplus\Q\epsilon$ ordered by the highest nonzero
coordinate. The family $(t^{n\epsilon})_{n\ge0}$ is strongly summable, since
its support is an increasing sequence and each exponent occurs once, and its
sum is the Hahn element $(1-t^\epsilon)^{-1}$. Yet its valuations do not tend
cofinally to $+\infty$. This is the most elementary warning against
identifying positive valuation with topological nilpotence, and it is why
criterion (b) of Lemma~\ref{hol:lem:certificate} is needed at all. No appeal
to convergence of partial sums at the $H$ scale is involved.
\end{remark}

\begin{remark}[The real-valued coarsening is imported, not reproved]
\label{hol:rem:coarsening}
When $\Gamma$ does have an order unit $h$, the companion
report~\cite[Lemma~9.1, \texttt{ent:lem:coarsening}]{repoentire} supplies an
order-preserving $\pi:\Gamma\to\R$ with $\pi(h)=1$ and kernel the convex
subgroup $H=\{\gamma:n\abs\gamma<h\text{ for all }n\ge1\}$, hence a nontrivial
real-valued non-Archimedean absolute value $\abs a_\pi=\exp(-\pi(v(a)))$
inducing the intrinsic topology. We cite that lemma and do not reprove it. No
proof in this report uses it: every argument below is carried out in $\Gamma$
itself, which is what makes the arguments valid when no order unit exists.
The coarsening is mentioned twice, as an optional rank-one picture of the
order-unit case (Remark~\ref{hol:rem:thetacoarse}) and for the warning of
Section~\ref{hol:sub:position} that the interpolation/B\'ezout dichotomy of
that report and the annihilating-operator dichotomy of this one share an
invariant but are different theorems.
\end{remark}

\subsection{The entire class and its closure properties}

Polynomials are entire, because all the evaluation families are finite. An
ordinary complex entire Taylor series need not be entire in the sense of
Definition~\ref{hol:def:entire}: at a monomial argument $x=t^\gamma$ the
family
\begin{equation}\label{hol:eq:monomialexp}
 \left(\frac{x^n}{n!}\right)_{n\ge0}
\end{equation}
is strongly summable exactly when $\gamma>0$, since $\gamma=0$ puts infinitely
many contributions at exponent zero and $\gamma<0$ gives a descending
sequence. Proposition~\ref{hol:prop:expdomain} computes the exact domain of
the corresponding family with an arbitrary Hahn argument.

\begin{proposition}[Closure]\label{hol:prop:closure}
$\E$ is a $\K$-algebra, $\sigma_q(\E)\subseteq\E$ for every $q\in\K^\times$,
and $D_z(\E)\subseteq\E$.
\end{proposition}

\begin{proof}
For addition, the evaluation supports lie in the union of two well-ordered
sets with finite exponent multiplicities. For multiplication, the doubly
indexed product family is strongly summable by
Lemma~\ref{hol:lem:support}\,(1) and (3), and grouping by ordinary total
degree preserves strong summability by the last clause of
Lemma~\ref{hol:lem:certificate}. Dilation is evaluation at $qx$. For $x\ne0$
the derivative family is
\[
 (n+1)a_{n+1}x^n=x^{-1}\cdot(n+1)a_{n+1}x^{n+1}.
\]
Discarding one term and multiplying each term by a nonzero ordinary integer
enlarges neither supports nor multiplicities, and multiplication by $x^{-1}$
preserves summability by Lemma~\ref{hol:lem:support}\,(3). At $x=0$ only the
constant term contributes. Iterate.
\end{proof}

\begin{proposition}[Faithfulness of evaluation]\label{hol:prop:faithful}
If $f\in\E$ vanishes at infinitely many distinct ordinary complex arguments,
then $f=0$ as a formal series. Consequently a linear equation with
rational-function coefficients that holds pointwise wherever its coefficients
are defined is also a formal equation in the sense of
Definition~\ref{hol:def:operators}.
\end{proposition}

\begin{proof}
Strong evaluation at $1$ says that the coefficient family $(a_n)$ has a
well-ordered union of supports and is locally finite at every exponent. Hence
for each $\gamma$
\[
 P_\gamma(T)=\sum_n(a_n)_\gamma T^n\in\C[T]
\]
is a polynomial, and its value at an ordinary complex $c$ is the coefficient
of $t^\gamma$ in $f(c)$. If $f$ vanishes at infinitely many such $c$, every
$P_\gamma$ has infinitely many roots and is zero, so every $a_n$ is zero. For
the last assertion, clear denominators; the left-hand side is in $\E$ by
Proposition~\ref{hol:prop:closure}, and nonzero polynomial denominators
exclude at most finitely many ordinary complex arguments.
\end{proof}

The nonexistence proofs may therefore be run on formal coefficients without
weakening their reading as functional obstructions, and no interchange of
infinite limits is needed anywhere.

The theorems below use only the valuation laws
\eqref{hol:eq:valuationlaws}, the existence of every monomial $t^\gamma$,
triviality of the valuation on $\C^\times$, and the support facts of this
section. They do not require algebraic closedness, spherical completeness,
Weierstrass preparation, or an embedding of $\Gamma$ into $\R$.

\subsection{Embedding the workspace into the surcomplex numbers}
\label{hol:sub:embedding}

If $\Gamma$ is an additive ordered subgroup of $\No$, the Conway normal-form
map
\begin{equation}\label{hol:eq:embedding}
 \sum_{\gamma\in S}c_\gamma t^\gamma
 \longmapsto\sum_{\gamma\in S}c_\gamma\omega^{-\gamma}
\end{equation}
identifies the real-coefficient Hahn field with a subfield of $\No$ and its
complexification with a subfield of $\No[i]$; a well-ordered $S$ becomes a
reverse-well-ordered set of Conway exponents. This viewpoint and its
set-sized support discipline are standard~\cite{bg,kks}. Here $t^\gamma$ is a
formal monomial and $\omega^{-\gamma}$ a Conway monomial; the notation does
not silently identify every formal power with a choice of surreal
exponentiation.

The results may therefore be read twice. Abstractly they concern Hahn fields
and need no surreal construction. Concretely they restrict the exact linear
descriptions available for surcomplex-valued power series on specified surreal
scales. Formula~\eqref{hol:eq:embedding} is the only bridge used, and no
global surcomplex exponential is invoked.

\section{The finite-step escape mechanism}\label{hol:sec:escape}

The following isolates the mechanism used throughout. It needs no formula for
the recurrence coefficients, no asymptotic slope for the valuations, and no
assumption that cancellations are rare.

\begin{lemma}[Escape chain]\label{hol:lem:escape}
Suppose a sequence $(a_n)$ in $\K$ satisfies, for every $n\ge N$,
\begin{equation}\label{hol:eq:generalrecurrence}
 c_0(n)a_n+\sum_{j=1}^sc_j(n)a_{n+j}=0,\qquad c_0(n)\ne0,
\end{equation}
and suppose $a_{n_0}\ne0$ for some $n_0\ge N$. Then $s\ge1$ and there is an
infinite sequence $n_0<n_1<n_2<\cdots$ with $j_r:=n_{r+1}-n_r\in\{1,\ldots,s\}$,
with $a_{n_r}\ne0$ and $c_{j_r}(n_r)\ne0$ for every $r$, and with
\begin{equation}\label{hol:eq:escapeineq}
 v(a_{n_{r+1}})-v(a_{n_r})\le v(c_0(n_r))-v(c_{j_r}(n_r)).
\end{equation}
\end{lemma}

\begin{proof}
Let $n\ge N$ with $a_n\ne0$. The leftmost term of
\eqref{hol:eq:generalrecurrence} is nonzero, so $s\ge1$ and at least one later
term is nonzero. If every nonzero later term had valuation strictly larger
than $v(c_0(n)a_n)$, the leading monomial of the leftmost term could not
cancel, contradicting \eqref{hol:eq:valuationlaws}. So some
$j\in\{1,\ldots,s\}$ has $c_j(n)a_{n+j}\ne0$ and
\[
 v(c_j(n)a_{n+j})\le v(c_0(n)a_n),
\]
which is \eqref{hol:eq:escapeineq}. Taking the least eligible $j$ makes the
recursion definite. The successor coefficient is nonzero and its index is
again at least $N$, so the construction never stops.
\end{proof}

\begin{remark}[The direction of the inequality matters]
\label{hol:rem:direction}
Solving the recurrence for its \emph{last} coefficient typically yields a
\emph{lower} bound on a valuation, which obstructs nothing. The comparison
above is made from the first nonzero coefficient forward and produces a later
coefficient whose valuation is controlled from above. Cancellation among the
other terms can raise the valuation of their sum, but cannot make every term
on the right have valuation strictly larger than the single nonzero term on
the left. Zeros and arbitrarily large cancellations elsewhere therefore do not
invalidate the chain; the chain records only that finite inequality,
repeatedly.
\end{remark}

\begin{theorem}[Evaluation exclusion criterion]\label{hol:thm:escapeexclusion}
In the situation of Lemma~\ref{hol:lem:escape}, let $x\in\K^\times$ and
suppose that for every $n\ge N$ and every $j\in\{1,\ldots,s\}$ with
$c_j(n)\ne0$,
\begin{equation}\label{hol:eq:boundincrements}
 v(c_0(n))-v(c_j(n))+jv(x)\le0.
\end{equation}
If $a_{n_0}\ne0$ for some $n_0\ge N$, then $(a_nx^n)_{n\ge0}$ is not strongly
summable. Equivalently: if $(a_nx^n)_{n\ge0}$ is strongly summable, then
$a_n=0$ for every $n\ge N$.
\end{theorem}

\begin{proof}
Build the chain of Lemma~\ref{hol:lem:escape} from $n_0$ and add
$n_{r+1}v(x)$ to each side of \eqref{hol:eq:escapeineq}, using
$n_{r+1}=n_r+j_r$:
\[
 v(a_{n_{r+1}}x^{n_{r+1}})
 \le v(a_{n_r}x^{n_r})+v(c_0(n_r))-v(c_{j_r}(n_r))+j_rv(x)
 \le v(a_{n_r}x^{n_r}).
\]
The evaluated leading exponents along the chain are therefore nonincreasing,
and Lemma~\ref{hol:lem:nonincreasing} applies.
\end{proof}

Note that \eqref{hol:eq:boundincrements} is a weak inequality. This is the
difference between an open and a closed excluded region, and it is available
only because Lemma~\ref{hol:lem:nonincreasing} also rules out a constant
leading exponent. Section~\ref{hol:sub:expdomain} exhibits an equation whose
excluded region is exactly the closed one, and
Section~\ref{hol:sub:coarseboundary} exhibits a coarsened comparison where the
boundary genuinely must be left in.

\begin{corollary}[Uniform bounded cost; no divisibility]
\label{hol:cor:boundedcost}
Suppose in \eqref{hol:eq:generalrecurrence} there is a single $B\in\Gamma$ with
$v(c_0(n))-v(c_j(n))\le B$ for all $n\ge N$ and all $j\ge1$ with $c_j(n)\ne0$.
Put $B^{+}=\max(B,0)\in\Gamma$. Then for every $x\ne0$ with $v(x)\le-B^{+}$,
strong summability of $(a_nx^n)_n$ forces $a_n=0$ for all $n\ge N$. Such $x$
exist: $x=t^{-B^{+}}$, and $x=t^{-\delta}$ for every $\delta\ge B^{+}$.
\end{corollary}

\begin{proof}
Since $B^{+}\ge0$ we have $v(x)\le0$, so $jv(x)\le v(x)\le-B^{+}\le-B$ for
$j\ge1$. Hence
$v(c_0(n))-v(c_j(n))+jv(x)\le B-B=0$, and
Theorem~\ref{hol:thm:escapeexclusion} applies. Every ordered abelian group
admits a maximum of finitely many elements and the monomial $t^{-B^{+}}$
exists, so no cofinal cyclic subgroup is needed.
\end{proof}

The absence of a cofinality requirement here is exactly why the mechanism
proves rigidity in groups of very large rank. It is also why the refined
bound below, which divides by $j$, is a refinement of the constant and never
of the conclusion.

\begin{corollary}[Refined periodic threshold; requires $\Gamma$ divisible]
\label{hol:cor:periodic}
Assume in addition that $\Gamma$ is divisible. Assume each slot of
\eqref{hol:eq:generalrecurrence} is either identically zero for $n\ge N$ or
nonzero for all $n\ge N$ with periodic valuation there of common period
$m\ge1$; enlarge $N$ past any initial nonperiodic range first.
Let $\alpha_{j,r}$ be the valuation of slot $j$ on the residue class
$r\bmod m$, the identically zero slots being omitted. Put
\begin{equation}\label{hol:eq:radiusbound}
 B_{\mathrm{rec}}
 =\max_{\substack{0\le r<m\\ 1\le j\le s,\ c_j\not\equiv0}}
   \frac{\alpha_{0,r}-\alpha_{j,r}}{j}\ \in\Gamma.
\end{equation}
If there is no active forward slot $j\ge1$, set $B_{\mathrm{rec}}=0$;
the recurrence then forces $a_n=0$ for all $n\ge N$ directly.
Then for every $x\ne0$ with $v(x)\le-B_{\mathrm{rec}}$, strong summability of
$(a_nx^n)_n$ forces $a_n=0$ for all $n\ge N$. In particular every solution
with $a_n\ne0$ for some $n\ge N$ --- so in particular every nonpolynomial
solution --- satisfies
\begin{equation}\label{hol:eq:excludedball}
 \Dom(f)\subseteq\{0\}\cup\{x\in\K^\times:v(x)>-B_{\mathrm{rec}}\},
 \qquad f=\sum_na_nz^n.
\end{equation}
\end{corollary}

\begin{proof}
The case with no forward slot was handled in the statement. Otherwise
divisibility makes each quotient $(\alpha_{0,r}-\alpha_{j,r})/j$ an element of
$\Gamma$, so $B_{\mathrm{rec}}$ is defined. If $v(x)\le-B_{\mathrm{rec}}$ then
$jv(x)\le-j\,(\alpha_{0,r}-\alpha_{j,r})/j=-(\alpha_{0,r}-\alpha_{j,r})$,
because multiplying a weak inequality by the positive integer $j$ preserves
it. Hence \eqref{hol:eq:boundincrements} holds and
Theorem~\ref{hol:thm:escapeexclusion} applies.
\end{proof}

The inequality in \eqref{hol:eq:excludedball} is strict on the right: at
$v(x)=-B_{\mathrm{rec}}$ the chain may keep a constant leading exponent rather
than descend, which is precisely why both halves of
Definition~\ref{hol:def:entire} are needed. The bound is an \emph{exclusion}
bound. It does not assert that every point on the other side lies in the
domain.

\section{Valuations on integer and geometric orbits}\label{hol:sec:orbits}

\subsection{Ordinary integers cannot create unbounded valuation loss}

\begin{lemma}[Integer evaluation]\label{hol:lem:integer}
Let $0\ne P(X)=\sum_kp_kX^k\in\K[X]$ and put $\beta(P)=\min_{p_k\ne0}v(p_k)$.
For all but finitely many nonnegative ordinary integers $n$,
\begin{equation}\label{hol:eq:integerstable}
 P(n)\ne0\qquad\text{and}\qquad v(P(n))=\beta(P).
\end{equation}
\end{lemma}

\begin{proof}
The polynomial $t^{-\beta(P)}P$ has coefficients of nonnegative valuation and
at least one of valuation zero, so its coefficientwise reduction
\[
 \overline P(X)=\sum_{v(p_k)=\beta(P)}\lc(p_k)X^k\in\C[X]
\]
is nonzero and has finitely many roots. The ordinary integers have distinct
images in $\C$. At every integer $n$ with $\overline P(n)\ne0$, the Hahn
coefficient of $t^{-\beta(P)}P(n)$ at exponent zero is nonzero, and no smaller
exponent occurs.
\end{proof}

Residue characteristic zero does real work here. The proof does not apply to
$p$-adic integer inputs, where positive integers can have arbitrarily large
valuation. Conversely no Diophantine approximation over $\C$ is relevant to
$v$: every nonzero complex number, however small in ordinary modulus, has
valuation zero.

\subsection{A geometric orbit of nonzero valuation}

\begin{lemma}[Eventually affine valuation]\label{hol:lem:affine}
Let $q\in\K^\times$ with $\lambda=v(q)\ne0$ and let
$0\ne P(X)=\sum_kp_kX^k\in\K[X]$. There are an index $k_P$ with $p_{k_P}\ne0$,
the value $\beta_P=v(p_{k_P})$, and $N_P\in\N$, such that
\begin{equation}\label{hol:eq:affine}
 v(P(q^n))=\beta_P+k_Pn\lambda\qquad(n\ge N_P).
\end{equation}
In particular $P(q^n)$ is eventually nonzero.
\end{lemma}

\begin{proof}
The nonzero terms have valuations $v(p_k)+kn\lambda$. For two distinct indices
the difference $v(p_k)-v(p_l)+(k-l)n\lambda$ is strictly monotone in $n$, so
its sign is eventually constant and it vanishes for at most one integer,
whether or not its values are cofinal in any convex subgroup. Finitely many
pairs give one common threshold beyond which every pairwise comparison is
strict and fixed, hence a unique minimizing index $k_P$. Its term cannot
cancel against terms of strictly larger valuation.
\end{proof}

\begin{remark}[A rank-one shortcut that fails]\label{hol:rem:shortcut}
The minimizing index need not be the smallest when $v(q)>0$. In
$\Gamma=\Q H\oplus\Q\epsilon$ with $H$ dominating every $n\epsilon$, take
$q=t^\epsilon$ and $P(X)=X+t^{-H}$: the constant term has smaller valuation
than $q^n$ for every $n$, however large. Conversely with $P(X)=1+t^{-H}X$ the
nonconstant term is the minimum for every $n$. The same phenomenon in the
other sign convention: in $\Gamma=\Q e_0\oplus\Q e_1$ with $e_1\gg e_0>0$,
$q=t^{e_0}$ and $P(X)=t^{e_1}+X$ give $v(P(q^n))=ne_0$ for every $n\ge0$, so
the nonzero constant coefficient never becomes the least-valued term. The
proof above uses the eventual ordering of finitely many affine functions, not
the Archimedean assertion that $n\epsilon$ eventually exceeds $H$, and not the
rank-one slogan that $q^n$ eventually beats every coefficient ratio.
\end{remark}

\subsection{Valuation-zero dilations, including residues that are roots of unity}

Distinct powers of a valuation-zero $q$ all have the same valuation, so the
previous lemma says nothing. Polynomial evaluation nevertheless has only a
finite periodic valuation pattern.

\begin{theorem}[Unit-orbit valuations]\label{hol:thm:unitorbit}
Let $q\in\K^\times$ have $v(q)=0$ and infinite multiplicative order, and let
$0\ne P\in\K[X]$. Then $v(P(q^n))$ is eventually finite and periodic. More
precisely, with $\zeta=\res(q)\in\C^\times$:
\begin{enumerate}[label=\textup{(\alph*)}]
\item if $\zeta$ is not a root of unity, $v(P(q^n))$ is eventually constant,
equal to the minimum valuation of the coefficients of $P$;
\item if $\zeta$ has exact order $h$, then $v(P(q^n))$ is eventually constant
on each residue class modulo $h$, so its eventual period divides $h$.
\end{enumerate}
\end{theorem}

\begin{proof}
For (a) write $P(X)=\sum_kp_kX^k$ and $\beta=\min_kv(p_k)$. The coefficient of
$P(q^n)$ at exponent $\beta$ is $\overline P(\zeta^n)$ with
$\overline P(X)=\sum_{v(p_k)=\beta}\lc(p_k)X^k$, a nonzero complex polynomial.
The points $\zeta^n$ are distinct, so this vanishes only finitely often, and
no exponent below $\beta$ occurs.

For (b) write
\[
 q=\zeta(1+u),\qquad \eta=v(u)>0,\qquad d=\lc(u).
\]
Here $u=q/\zeta-1$ has positive valuation because $q/\zeta$ has valuation zero
and residue $1$, and $u\ne0$ because otherwise $q=\zeta$ would have finite
order. For every positive ordinary integer $n$ the finite binomial expansion
gives
\begin{equation}\label{hol:eq:yn}
 y_n:=(1+u)^n-1,\qquad v(y_n)=\eta,\qquad \lc(y_n)=nd,
\end{equation}
the leading coefficient being nonzero in characteristic zero. There is no
convergence assertion: for each fixed $n$ only finitely many powers occur.

Fix $r\in\{0,\ldots,h-1\}$. For $n\equiv r\pmod h$ we have $\zeta^n=\zeta^r$,
hence $q^n=\zeta^r(1+y_n)$. Expand the nonzero polynomial
\begin{equation}\label{hol:eq:unitexpand}
 P(\zeta^r(1+Y))=\sum_{k=0}^{\deg P}b_{r,k}Y^k
\end{equation}
and set
\begin{equation}\label{hol:eq:betar}
 \beta_r=\min_{b_{r,k}\ne0}\bigl(v(b_{r,k})+k\eta\bigr),
 \qquad J_r=\{k:v(b_{r,k})+k\eta=\beta_r\}.
\end{equation}
Then $P(q^n)=\sum_kb_{r,k}y_n^k$, and its coefficient at exponent $\beta_r$ is
the value at $n$ of
\begin{equation}\label{hol:eq:residuepoly}
 R_r(T)=\sum_{k\in J_r}\lc(b_{r,k})d^kT^k\in\C[T].
\end{equation}
This polynomial is nonzero: its terms have distinct degrees and all the
indicated coefficients are nonzero. Outside its finitely many integer roots
the coefficient at $\beta_r$ survives, so $v(P(q^n))=\beta_r$ eventually on
that class. Finitely many classes give one finite exceptional range.
\end{proof}

\begin{remark}[A second normalization of the same proof]
\label{hol:rem:unitalt}
Sources~08 and~09 both prove Theorem~\ref{hol:thm:unitorbit}\,(b), by two
slightly different normalizations, and we record the other one because it is
occasionally more convenient. Instead of factoring out the complex root of
unity, put $q^h=1+u$ directly, with $\eta=v(u)>0$ and $u\ne0$ for the same
reason. Write $n=r+hk$ and expand the polynomial $Q_r(Y)=P(q^rY)$ as
$Q_r(1+V)=\sum_kb_{r,k}V^k$. The finite binomial expansion gives
$V_k=(1+u)^k-1$ with $v(V_k)=\eta$ and $\lc(V_k)=k\lc(u)$, and
$P(q^n)=Q_r(1+V_k)$. The minimum construction \eqref{hol:eq:betar} and the corresponding
nonzero residue polynomial \eqref{hol:eq:residuepoly}, now in the variable
$k$ rather than $n$, finish the argument. Neither version sums a binomial
series to a noninteger exponent; every expansion is finite, and arbitrary
supports in $q$ and in the coefficients of $P$ enter only through the leading
valuation and leading coefficient of finitely many Hahn elements.
\end{remark}

\begin{example}[A genuine two-phase pattern]\label{hol:ex:unitperiod}
Let $q=-(1+t^\eta)$ with $\eta>0$. Then $q$ is nontorsion, $\res(q)=-1$ has
order two, and for $P(X)=X-1$ and $n\ge1$
\[
 v(P(q^n))=v(q^n-1)=
 \begin{cases}
 0,&n\text{ odd},\\
 \eta,&n\text{ even}.
 \end{cases}
\]
For odd $n$ the residue is $-2$; for even positive $n$ the leading term is
$nt^\eta$. The valuation sequence is periodic, not the dilation itself, and a
claim of eventual \emph{constancy} without splitting residue classes would be
false. The escape argument uses the bounded valuation costs; it does not
replace $q$ by its residue.
\end{example}

\begin{example}[Finite exceptional indices really occur]
\label{hol:ex:unitexception}
For $q=1+t^\eta$ and $P(X)=X-1-7t^\eta$, the leading coefficient of $P(q^n)$
at exponent $\eta$ is $n-7$, so $v(P(q^n))=\eta$ for positive $n\ne7$ while
$v(P(q^7))=2\eta$. With $P(X)=X-1-3t^\eta$ the exceptional index is $n=3$,
where $P(q^3)=3t^{2\eta}+t^{3\eta}$. The exceptions are harmless for the
escape argument, since it needs only an eventual range, but they cannot be
ignored in a statement asserted for all indices. This is why
Lemma~\ref{hol:lem:integer} and Theorem~\ref{hol:thm:unitorbit} both allow
finitely many exceptions even when $\res(q)=1$.
\end{example}

For an ordinary complex $q$ with $\abs q_\C\ne1$, the complex modulus does not
alter these conclusions: as an element of the Hahn coefficient field, $q$
still has valuation zero. The ordinary modulus and the scale valuation answer
different questions.

\section{Differential rigidity}\label{hol:sec:ode}

\subsection{From an operator to a forward recurrence}

A formal series is D-finite exactly when its derivatives span a
finite-dimensional space over $\K(z)$. The D-finite/P-recursive equivalence is
classical~\cite[Theorem~1.5]{stanley}; the following records the exact finite
conversion, because the valuations of the resulting recurrence, not ordinary
coefficient growth, are what the proof consumes.

\begin{lemma}[Differential coefficient recurrence]\label{hol:lem:drec}
Write
\begin{equation}\label{hol:eq:operator}
 L=\sum_{r=0}^R\sum_{k=0}^mc_{r,k}z^kD_z^r,\qquad c_{r,k}\in\K,
\end{equation}
with at least one $c_{r,k}\ne0$, and let $f=\sum_{n\ge0}a_nz^n$. For $n\ge m$
the coefficient of $z^n$ in $Lf$ is
\begin{equation}\label{hol:eq:odecoeff}
 \sum_{s=-m}^{R}Q_s(n)a_{n+s},\qquad
 Q_s(X)=\sum_{r-k=s}c_{r,k}\fall{X+s}{r},
\end{equation}
where $\fall Yr=Y(Y-1)\cdots(Y-r+1)$ and $\fall Y0=1$. At least one $Q_s$ is
nonzero. Writing $s_0$ for the least active shift and $u=n+s_0$, every formal
solution of $Lf=0$ satisfies, for all sufficiently large $u$,
\begin{equation}\label{hol:eq:precursive}
 \sum_{j=0}^{s}P_j(u)a_{u+j}=0,\qquad P_j(X)=Q_{s_0+j}(X-s_0),
 \qquad P_0\ne0,
\end{equation}
with $s=R-s_0$. If $s=0$, every formal solution is a polynomial.
\end{lemma}

\begin{proof}
The operator $z^kD_z^r$ carries the term of index $n+r-k$ to degree $n$ and
multiplies it by $\fall{n+s}{r}$ with $s=r-k$, which gives
\eqref{hol:eq:odecoeff}. For a fixed shift $s$ the distinct derivative orders
$r$ contribute polynomials of distinct degrees in $X$, so a nonempty group
containing a nonzero coefficient cannot cancel identically and some $Q_s$ is
nonzero. Reindexing gives \eqref{hol:eq:precursive}. If $s=0$ the recurrence
reads $P_0(u)a_u=0$ with $P_0\ne0$, so by Lemma~\ref{hol:lem:integer} all but
finitely many $a_u$ vanish.
\end{proof}

\begin{proposition}[Uniform obstruction for P-recursive coefficients]
\label{hol:prop:precursive}
Suppose $(a_n)$ satisfies \eqref{hol:eq:precursive} for all $n\ge N_1$, with
$P_j\in\K[X]$ and $P_0\ne0$. Let $\beta_j=\beta(P_j)$ for the nonzero $P_j$
and put
\[
 B=\max\bigl(\{0\}\cup\{\beta_0-\beta_j:j\ge1,\ P_j\ne0\}\bigr)\ \ge0.
\]
There is $N\ge N_1$, beyond the finitely many exceptional integers of
Lemma~\ref{hol:lem:integer} for the finitely many nonzero $P_j$, such that for
every $x\ne0$ with $v(x)\le-B$, strong summability of $(a_nx^n)_n$ forces
$a_n=0$ for all $n\ge N$. If $\Gamma$ is divisible one may replace $B$ by the
smaller threshold
\begin{equation}\label{hol:eq:refinedthreshold}
 B_{\mathrm{rec}}
 =\max\Bigl(\Bigl\{\tfrac{\beta_0-\beta_j}{j}:j\ge1,\ P_j\ne0\Bigr\}\Bigr),
\end{equation}
which is \eqref{hol:eq:radiusbound} with period $m=1$.
If the displayed set is empty, take $B_{\mathrm{rec}}=0$; then every
solution terminates beyond $N$ without an evaluation hypothesis.
\end{proposition}

\begin{proof}
Discard the identically zero $P_j$. By Lemma~\ref{hol:lem:integer} there is
$N$ beyond which $v(P_j(n))=\beta_j$ for all the finitely many nonzero $P_j$
simultaneously. Apply Corollary~\ref{hol:cor:boundedcost} with this $B$, or
Corollary~\ref{hol:cor:periodic} with $m=1$ when $\Gamma$ is divisible. If no
$j\ge1$ remains, the recurrence forces the same vanishing directly.
\end{proof}

These are certified exterior obstructions, not a computation of the exact
evaluation domain of each solution. Individual solutions can have stricter
restrictions or can terminate much earlier.

\begin{proof}[Proof of Theorem~\ref{hol:main:ode}]
Lemma~\ref{hol:lem:drec} converts $L$ into \eqref{hol:eq:precursive}. Take
$B_L=B$ from Proposition~\ref{hol:prop:precursive} and take $N_L$ to be an
positive integer beyond the recurrence's starting index and beyond every nonnegative
integer root of the finitely many residue polynomials $\overline{P_j}$. If a
formal solution is strongly evaluable at one $x\ne0$ with $v(x)\le-B_L$, then
$a_n=0$ for all $n\ge N_L$, so its degree is less than $N_L$. An entire
solution is in particular strongly evaluable at the available monomial
$t^{-B_L}$, hence is a polynomial. Conversely a polynomial is entire and is
annihilated by a sufficiently high derivative, which gives the displayed
equality.
\end{proof}

The closed region $v(x)\le-B_L$, rather than an open one, comes from the weak
inequality in Theorem~\ref{hol:thm:escapeexclusion}, hence from
Lemma~\ref{hol:lem:nonincreasing}. A polynomial solution of degree at least
$N_L$ would in any case contradict the recurrence at its last nonzero
coefficient, since $P_0(n)$ no longer vanishes there.

Theorem~\ref{hol:main:ode} is stronger than merely excluding a globally
extended exponential. The operator coefficients may themselves be arbitrary
Hahn series carrying infinite surreal scales. There are nevertheless only
finitely many of them, so their leading valuations give a finite obstruction
that one additional exterior scale defeats.

\subsection{An exact boundary, not only a nonexistence result}
\label{hol:sub:expdomain}

\begin{proposition}[Formal exponential domain]\label{hol:prop:expdomain}
For $\eta\in\Gamma$ let
\[
 E_\eta(z)=\sum_{n\ge0}\frac{t^{n\eta}}{n!}z^n.
\]
Then $D_zE_\eta=t^\eta E_\eta$ and
\begin{equation}\label{hol:eq:expdomain}
 \Dom(E_\eta)=\{0\}\cup\{x\in\K^\times:v(x)>-\eta\}.
\end{equation}
\end{proposition}

\begin{proof}
The differential identity is a formal coefficient comparison. Put $y=t^\eta x$
for $x\ne0$, so the $n$th evaluation term is $y^n/n!$.

If $v(y)>0$, write $y=c\,t^{v(y)}(1+u)$ with $c\in\C^\times$ and $u$ zero or
of positive valuation, and put $S=\supp(u)\subseteq\Gamma_{>0}$. Since
$(1+u)^n$ is a finite integer power, $\supp(y^n)\subseteq h_n+M(S)$ with
$h_n=nv(y)$, and $(h_n)$ is strictly increasing, so its values are well
ordered and each is taken once. Lemma~\ref{hol:lem:certificate}\,(b) makes the
family strongly summable; the factor $1/n!$ is a nonzero complex scalar and
changes no support.

If $v(y)<0$ the leading exponents $nv(y)$ strictly decrease. If $v(y)=0$, then
$y^n/n!$ has the nonzero coefficient $\lc(y)^n/n!$ at exponent zero for every
$n$, so local finiteness fails. These cases exhaust $x\ne0$, and $v(y)>0$ is
$v(x)>-\eta$.
\end{proof}

The recurrence of $E_\eta$ is $(n+1)a_{n+1}-t^\eta a_n=0$, with one active
shift $j=1$, so $\alpha_0=\eta$, $\alpha_1=0$, and
\eqref{hol:eq:radiusbound} gives $B_{\mathrm{rec}}=\eta$ with no division by
an integer at all. Comparing with \eqref{hol:eq:expdomain}: the excluded
region $v(x)\le-\eta$ of Corollary~\ref{hol:cor:periodic} is here
\emph{exactly} the complement of the domain, boundary included. The general
exclusion bound can therefore be attained, and the closed form of the boundary
is not an artefact of a lossy estimate. This example is the reason for keeping
Lemma~\ref{hol:lem:nonincreasing} rather than the strictly descending version
alone.

There is no conflict with any scalar exponential on a surreal field: $E_\eta$
is an ordinary power series in $z$, not a total exponential evaluated by an
independent construction. Taking $\eta=0$ recovers the elementary observation
\eqref{hol:eq:monomialexp}.

\subsection{Algebraic and system consequences}

\begin{corollary}[Algebraic entire series]\label{hol:cor:algebraic}
An element of $\E$ algebraic over $\K(z)$ belongs to $\K[z]$.
\end{corollary}

\begin{proof}
Let $F=\K(z,f)$ be the finite extension generated by the algebraic series. In
characteristic zero its minimal polynomial is separable, and differentiating
that polynomial expresses $D_zf$ inside $F$, because the derivative with
respect to the algebraic variable is nonzero at $f$ and hence invertible in
$F$. So $D_z$ preserves $F$, all derivatives of $f$ lie in the
finite-dimensional $\K(z)$-space $F$, and $f$ is D-finite. Apply
Theorem~\ref{hol:main:ode}.
\end{proof}

This is not offered as a new algebraic rigidity theorem; it records that the
linear result subsumes the algebraic case without any factorization theorem
for $\E$.

\begin{corollary}[Rational differential systems]\label{hol:cor:systems}
If a finite vector $Y$ of strongly entire formal series satisfies
$D_zY=A(z)Y$ with $A\in M_r(\K(z))$, every component of $Y$ is a polynomial.
The same holds for an inhomogeneous system with a rational forcing vector,
provided all its solution components are strongly entire.
\end{corollary}

\begin{proof}
The $\K(z)$-span of the components is stable under $D_z$, so the successive
derivatives of each component span a finite-dimensional space. For $Y'=AY+b$
append the constant component $1$ and use
$\left(\begin{smallmatrix}A&b\\0&0\end{smallmatrix}\right)$.
\end{proof}

\begin{proposition}[Rational inhomogeneous equations]
\label{hol:prop:inhomogeneous}
The polynomiality conclusions of this report remain valid when the right-hand
side of a linear equation is a rational function of $z$, provided the
corresponding homogeneous hypothesis on $L$ or on $q$ holds.
\end{proposition}

\begin{proof}
Clear denominators in the coefficients and in the right-hand side. The result
has a polynomial right-hand side, so the coefficient recurrence is homogeneous
for all sufficiently large indices, and the escape statements were formulated
for eventual recurrences. The left-hand operator is still nonzero, having been
multiplied by a nonzero polynomial.
\end{proof}

This applies, for instance, to the inhomogeneous first-order theta identity
\eqref{hol:eq:thetafirst} when the parameter valuation is noncofinal: it is
not necessary to homogenize it in order to see the obstruction.

\begin{example}[A sparse recurrence does not escape the obstruction]
\label{hol:ex:airy}
The formal Airy-type equation $f''-zf=0$ has the recurrence
\[
 a_{n+3}=\frac{a_n}{(n+3)(n+2)}\quad(n\ge0),\qquad a_2=0.
\]
With $a_0=1$ and $a_1=0$ there are infinitely many nonzero coefficients, at
indices $3n$, all of valuation zero. At $x=t^{-\delta}$ their leading
exponents are $-3n\delta$, strictly decreasing for every $\delta>0$. Gaps
between nonzero coefficients do not help: the escape chain follows only the
surviving congruence class.
\end{example}

\begin{remark}[What transcendence has actually been proved]
\label{hol:rem:transcendence}
Every nonpolynomial member of $\E$ is transcendental over $\K(z)$ and
satisfies no nonzero \emph{linear} differential equation over $\K(z)$, and no
finite family $f,D_zf,\ldots,D_z^Rf$ is linearly dependent over $\K(z)$. This
does not imply differential transcendence in the nonlinear sense. A relation
\[
 P\bigl(z,f,D_zf,\ldots,D_z^Rf\bigr)=0
\]
involving products of $f$ and its derivatives lies outside the present proofs
unless it can separately be reduced to a finite linear system. The
distinction is not cosmetic: the escape mechanism exploits a finite
\emph{linear} cancellation in which one nonzero coefficient forces a later
nonzero coefficient, whereas nonlinear coefficient equations involve
convolutions of many earlier coefficients and supply no such forward step
without further work.

That further work is the nonlinear part of this report
(Sections~\ref{hol:nl:sec:setting}--\ref{hol:nl:sec:consequences}, from
source~10). It does not repair the escape chain; it replaces the recurrence by
a finite initial polynomial at one large scale. Every relation of order
$R\le1$ is thereby excluded (Theorem~\ref{hol:nl:main:first}), and so is every
relation of higher order whose residue corner polynomial is nonzero
(Theorem~\ref{hol:nl:thm:corner}). Relations of order at least two with a
vanishing residue corner remain outside all proofs in this report.
\end{remark}

\section{Generic lines and several-variable rigidity}\label{hol:sec:multi}

A single differential equation in several variables is not enough to force
polynomiality: it could ignore most of the variables. The correct hypothesis
is finite-dimensionality of the span of \emph{all} mixed partial derivatives,
which is what Theorem~\ref{hol:main:multi} assumes.

\begin{lemma}[Avoiding countably many polynomial conditions]
\label{hol:lem:generic}
For any countable family of nonzero polynomials in $\K[z_1,\ldots,z_d]$ there
is $c\in\C^d$ at which none of them vanishes.
\end{lemma}

\begin{proof}
For one polynomial $H$, take the least Hahn exponent among its finitely many
nonzero coefficients; the coefficient at that exponent is a nonzero
$\overline H\in\C[z_1,\ldots,z_d]$, and $\overline H(c)\ne0$ implies
$H(c)\ne0$. The coefficients of the countably many polynomials $\overline H$
form a countable subset of $\C$, and the field they generate over $\Q$ is
countable. Choose $c_1,\ldots,c_d$ algebraically independent over that field,
successively outside the algebraic closure of a countable field, which remains
countable while $\C$ is uncountable.
\end{proof}

The same countable-field device is used in the companion report held in
\path{docs/surcomplex/entire-functions-at-arbitrary-rank/}
(\texttt{ent:mv:thm:exceptional}) to show that the set of polynomial
directions of a nonpolynomial series is not all of projective space.

\begin{proof}[Proof of Theorem~\ref{hol:main:multi}]
Suppose $F$ is not a polynomial and write it by total degree,
\begin{equation}\label{hol:eq:homogeneous}
 F(z)=\sum_{n\ge0}F_n(z),\qquad F_n\text{ homogeneous of degree }n,
\end{equation}
so infinitely many $F_n$ are nonzero. Let $g_1,\ldots,g_\sigma$ be a basis of
the finite-dimensional $\K(z_1,\ldots,z_d)$-span of all mixed derivatives of
$F$, taken to be mixed derivatives themselves with $g_1=F$. The span is stable
under differentiation, so there are rational matrices $A_i$ with
\begin{equation}\label{hol:eq:pfaffian}
 \partial_ig=A_i(z)g\qquad(1\le i\le d).
\end{equation}
These matrices have finitely many nonzero denominator polynomials; choose one
nonzero homogeneous component of each. Apply Lemma~\ref{hol:lem:generic} to
every nonzero $F_n$, to those finitely many chosen components, and to the
coordinate polynomials $z_i$, obtaining $c\in\C^d$ with $F_n(c)\ne0$ whenever
$F_n\ne0$ and with no denominator becoming the zero polynomial under
$z=wc$. Then
\[
 h(w)=F(wc)=\sum_{n\ge0}F_n(c)w^n
\]
is nonpolynomial.

For each $x\in\K$ the family
$(a_\alpha c^\alpha x^{\abs\alpha})_{\alpha\in\N^d}$ is strongly summable, by
entireness of $F$ at $(c_1x,\ldots,c_dx)$. Group it into the finite blocks of
fixed total degree; by the last clause of Lemma~\ref{hol:lem:certificate},
$\sum_nF_n(c)x^n$ is strongly summable, so $h\in\E$.

The formal chain rule, applied after clearing denominators in
\eqref{hol:eq:pfaffian}, gives
\begin{equation}\label{hol:eq:restricted}
 \frac{d}{dw}g(wc)=\Bigl(\sum_{i=1}^dc_iA_i(wc)\Bigr)g(wc).
\end{equation}
Every denominator here is a nonzero element of $\K[w]$. It may vanish at
$w=0$, which is harmless: read the identity in the Laurent series field
$\K((w))$, with matrices in $M_\sigma(\K(w))$. The finite $\K(w)$-span of the
components of $g(wc)$ is preserved by $d/dw$, so its first component $h$ is
D-finite over $\K(w)$, and Theorem~\ref{hol:main:ode} forces $h$ to be a
polynomial. This contradicts its construction.
\end{proof}

\begin{corollary}[Multivariate algebraic rigidity]\label{hol:cor:multialg}
If a strongly entire $F\in\K[[z_1,\ldots,z_d]]$ is algebraic over
$\K(z_1,\ldots,z_d)$, it is a polynomial.
\end{corollary}

\begin{proof}
Each partial derivation extends to the finite separable extension generated by
$F$, by differentiation of its minimal polynomial, so all mixed partial
derivatives lie in one finite-dimensional space. Apply
Theorem~\ref{hol:main:multi}.
\end{proof}

The generic line is an existence device, not a promised effective choice of
ordinary complex constants. The proof does not say that one prescribed line
detects every nonpolynomial multivariate series; the line is chosen only after
the countably many nonzero homogeneous parts and the finitely many rational
connection matrices are known.

\section{Dilation equations and the necessity of a cofinal scale}
\label{hol:sec:qnecessity}

\subsection{Exact coefficient conversion}

\begin{lemma}[Dilation coefficient recurrence]\label{hol:lem:qrec}
Let $q\in\K^\times$ have infinite multiplicative order and let $f=\sum a_nz^n$
satisfy \eqref{hol:eq:qlinear}. Write $p_\ell(z)=\sum_kp_{\ell,k}z^k$ and
$Q_k(X)=\sum_\ell p_{\ell,k}X^\ell$, let $d$ and $e$ be the largest and
smallest $k$ with $Q_k\ne0$, and put $s=d-e$. Then for every $n\ge0$
\begin{equation}\label{hol:eq:qrec}
 \sum_{h=0}^sR_h(q^n)a_{n+h}=0,\qquad
 R_h(X)=Q_{d-h}(q^hX),
\end{equation}
with $R_0=Q_d\ne0$ and $R_s=Q_e(q^s\,\cdot\,)\ne0$. Each nonzero $R_h$
satisfies $R_h(q^n)\ne0$ for all but finitely many $n$. If $s=0$, every formal
solution is a polynomial.
\end{lemma}

\begin{proof}
The coefficient of $z^{n+d}$ in $\sum_\ell p_\ell(z)f(q^\ell z)$ is
$\sum_{k=e}^{d}Q_k(q^{n+d-k})a_{n+d-k}$; substituting $h=d-k$ and using
$Q_{d-h}(q^{n+h})=Q_{d-h}(q^h\cdot q^n)$ gives \eqref{hol:eq:qrec}. Since $q$
has infinite order the points $q^n$ are distinct, and a nonzero univariate
polynomial vanishes at only finitely many of them; this already gives eventual
nonvanishing, independently of the finer valuation results. If $s=0$ the
relation reads $R_0(q^n)a_n=0$ for every $n$, so all but finitely many $a_n$
vanish.
\end{proof}

The powers $q^h$ in \eqref{hol:eq:qrec} matter: dropping them changes the
recurrence and can change its valuation constants. Equivalently, taking
$m=\max_\ell\deg p_\ell$ and reading off the coefficient of $z^{n+m}$ gives
$\sum_jP_j(q^n)a_{n+j}=0$ with
$P_j(X)=\sum_\ell p_{\ell,m-j}q^{\ell j}X^\ell$ and $P_0=Q_m\ne0$; both
conversions are checked on finite exact examples by the programs of
Section~\ref{hol:sec:verification}. If $m=0$ the equation already reads
$P_0(q^n)a_n=0$ and all solutions are polynomials with no entireness
assumption at all, so we may assume $s\ge1$ below.

\subsection{The valuation-zero case}

\begin{corollary}[Every unit dilation is globally rigid]
\label{hol:cor:unitrigid}
Let $q\in\K^\times$ have $v(q)=0$ and infinite multiplicative order. Every
linearly $q$-difference finite member of $\E$ is a polynomial, and each fixed
equation \eqref{hol:eq:qlinear} has a uniform exterior obstruction and a
uniform degree bound. If $\Gamma$ is in addition divisible, every
nonpolynomial formal solution has its domain bounded as in
\eqref{hol:eq:excludedball} by the constant \eqref{hol:eq:radiusbound} formed
from the eventual residue-class valuations of \eqref{hol:eq:qrec}.
\end{corollary}

\begin{proof}
Apply Theorem~\ref{hol:thm:unitorbit} simultaneously to the finitely many
nonzero $R_h$ of Lemma~\ref{hol:lem:qrec}. Their valuations take only finitely
many values after a common starting index, with common period one or the exact
order $h_0$ of $\res(q)$. All differences $v(R_0(q^n))-v(R_h(q^n))$ are
therefore bounded above by a single element of $\Gamma$, and
Corollary~\ref{hol:cor:boundedcost} gives the exterior obstruction and, via
the residue polynomials, the degree bound. When $\Gamma$ is divisible,
Corollary~\ref{hol:cor:periodic} applies with $m=1$ if the residue has
infinite order, and $m=h_0$ otherwise. In both cases take the starting
index beyond the exceptional range.
\end{proof}

This includes parameters arbitrarily close to a root of unity without being
one. The finite exceptional cancellations of Theorem~\ref{hol:thm:unitorbit}
can move the starting index but do not remove the bound.

\subsection{Nonzero but noncofinal valuation}

\begin{proposition}[Noncofinal dilation scales force polynomiality]
\label{hol:prop:noncofinal}
Let $q\in\K^\times$ with $\beta=v(q)\ne0$ and suppose $\abs\beta$ is not an
order unit of $\Gamma$. Every linearly $q$-difference finite member of $\E$ is
a polynomial, with a uniform exterior obstruction and a uniform degree bound
for each fixed equation. Divisibility of $\Gamma$ is not used.
\end{proposition}

\begin{proof}
By Lemma~\ref{hol:lem:affine} applied to the finitely many nonzero $R_h$ of
\eqref{hol:eq:qrec}, there is a common index beyond which
\[
 v(R_h(q^n))=\alpha_h+k_hn\beta,\qquad k_h\in\N.
\]
Put $V=\abs\beta>0$, choose $B\in\Gamma$ with $B\ge0$ and
$B\ge\abs{\alpha_0-\alpha_h}$ for all relevant $h$, and put
$M=\max_h\abs{k_0-k_h}\in\N$. Then for $h\ge1$
\begin{equation}\label{hol:eq:linearbound}
 v(R_0(q^n))-v(R_h(q^n))=\alpha_0-\alpha_h+(k_0-k_h)n\beta\le B+MnV.
\end{equation}
Because $\abs\beta$ is not an order unit, \eqref{hol:eq:dominates} supplies
$\gamma>0$ with $\gamma>NV$ for every $N\in\N$; in particular $MnV<\gamma$ for
all $n$, trivially so if $M=0$. Hence the right side of
\eqref{hol:eq:linearbound} is bounded above by the single element $B+\gamma$.
Corollary~\ref{hol:cor:boundedcost} with this bound excludes every $x\ne0$
with $v(x)\le-(B+\gamma)$, and the explicit witness $x=t^{-(B+\gamma)}$ is
available. Since an entire solution is evaluable there, it satisfies $a_n=0$
beyond the common starting index, hence is a polynomial of degree less than
that index. Nothing in this argument divides an element of $\Gamma$ by an
integer.
\end{proof}

The proof covers two different higher-rank situations. The group may have no
order unit at all, so every nonzero dilation valuation is noncofinal;
alternatively the group may have an order unit while the chosen $q$ lives on a
smaller scale. Global $q$-difference finiteness depends on the valuation of
that particular dilation, not only on the group. The argument also requires no
asymptotic formula for $v(a_n)$ and no eventually quadratic coefficient
valuation: a single infinite escape chain suffices, which is why cancellations
and sparse zero coefficients do not force a separate Skolem--Mahler--Lech
analysis.

\subsection{An intrinsic coarsened exclusion bound}

The previous proof produces more than one bad point. Let $\beta=v(q)\ne0$, let
$\Delta=\Delta_{\abs\beta}$ be the convex subgroup \eqref{hol:eq:delta}, and
assume $\Delta\ne\Gamma$. Let $\pi:\Gamma\to\Gamma/\Delta$ be the ordered
quotient and put $w=\pi\circ v$. With the eventual constants $\alpha_h$ of the
previous proof, put
\begin{equation}\label{hol:eq:coarsebound0}
 \overline B^{\flat}_{\mathrm{rec}}
 =\max\bigl(\{0\}\cup\{\pi(\alpha_0)-\pi(\alpha_h):1\le h\le s,\ R_h\ne0\}\bigr),
\end{equation}
and, when $\Gamma$ is divisible so that $\Gamma/\Delta$ is a divisible ordered
group,
\begin{equation}\label{hol:eq:coarsebound}
 \overline B_{\mathrm{rec}}
 =\max_{1\le h\le s,\ R_h\ne0}\frac{\pi(\alpha_0)-\pi(\alpha_h)}{h}.
\end{equation}

\begin{theorem}[Coarsened exclusion]\label{hol:thm:coarse}
Let $f=\sum a_nz^n$ be a nonpolynomial formal solution of
\eqref{hol:eq:qlinear} with $v(q)\ne0$ and $\Delta=\Delta_{\abs{v(q)}}$
proper.
\begin{enumerate}[label=\textup{(\alph*)}]
\item Under Convention~\ref{hol:conv:group} alone, every $x\ne0$ with
$-w(x)>\overline B^{\flat}_{\mathrm{rec}}$ lies outside $\Dom(f)$.
\item If $\Gamma$ is in addition divisible, every $x\ne0$ with
$-w(x)>\overline B_{\mathrm{rec}}$ lies outside $\Dom(f)$.
\end{enumerate}
\end{theorem}

\begin{proof}
Coarsening by $\Delta$ kills $n\beta$, since $n\beta\in\Delta$ for every $n$,
so the projected valuations of the recurrence coefficients are the constants
$\pi(\alpha_h)$. In case~(b), $-w(x)>\overline B_{\mathrm{rec}}$ gives
$hw(x)<-(\pi(\alpha_0)-\pi(\alpha_h))$ for every relevant $h\ge1$. In case~(a)
we have $w(x)<-\overline B^{\flat}_{\mathrm{rec}}\le0$, hence
$hw(x)\le w(x)<-\overline B^{\flat}_{\mathrm{rec}}
 \le-(\pi(\alpha_0)-\pi(\alpha_h))$. Either way the projected leading
valuation along an escape chain strictly decreases at every step. A strict
decrease in the ordered quotient is a strict decrease in $\Gamma$, so the
chain produces an infinite strictly descending sequence in the union of the
supports.
\end{proof}

The boundary is deliberately \emph{not} excluded here. Equal leading values in
the quotient can correspond to a strictly increasing sequence in $\Gamma$, and
hence to a strongly summable family; Section~\ref{hol:sub:coarseboundary}
exhibits exactly that. Unlike Corollary~\ref{hol:cor:periodic}, a coarsened
comparison alone cannot decide local finiteness at an exponent of $\Gamma$.

\section{Partial theta series: sharp witness and exact domain}
\label{hol:sec:theta}

\subsection{The classical identities}

For $p\in\K^\times$ with $\beta=v(p)>0$ set
\begin{equation}\label{hol:eq:theta}
 \Theta_p(z)=\sum_{n\ge0}p^{\binom n2}z^n
 =\sum_{n\ge0}p^{n(n-1)/2}z^n.
\end{equation}
This is a standard partial theta series up to a harmless sign
convention~\cite{andrewswarnaar}. None of its coefficients vanishes, so it is
not a polynomial. The coefficient identity
$p^{n(n-1)/2}=p^{\,n-1}p^{(n-1)(n-2)/2}$ for $n\ge1$ gives
\begin{equation}\label{hol:eq:thetafirst}
 \Theta_p(z)=1+z\Theta_p(pz),
\end{equation}
and applying $\sigma_p-1$ to $(1-z\sigma_p)\Theta_p=1$, with
$\sigma_pz=pz\sigma_p$, yields the homogeneous equation
\begin{equation}\label{hol:eq:thetahom}
 \Theta_p(z)-(1+z)\Theta_p(pz)+pz\Theta_p(p^2z)=0,
\end{equation}
which is also obtained by substituting $pz$ for $z$ in
\eqref{hol:eq:thetafirst} and subtracting the two inhomogeneous identities.
The existence of these identities is not the new part. What matters is
precisely where the series is strongly evaluable, for an arbitrary-rank value
group and an arbitrary Hahn parameter.

\subsection{Exact domain with arbitrary coefficient tails}

\begin{theorem}[Exact partial theta domain]\label{hol:thm:thetadomain}
Let $\Gamma$ be as in Convention~\ref{hol:conv:group} and let $p\in\K^\times$
with $\beta=v(p)>0$, and let $\Delta_\beta$ be as in \eqref{hol:eq:delta}.
Then
\begin{align}
 \Dom(\Theta_p)
 &=\{0\}\cup\{x\in\K^\times:v(x)\ge-N\beta\text{ for some }N\in\N\}
 \label{hol:eq:thetadomain}\\
 &=\{0\}\cup\{x\in\K^\times:(v\bmod\Delta_\beta)(x)\ge0\}.\nonumber
\end{align}
In particular $\Theta_p\in\E$ if and only if $\beta$ is an order unit of
$\Gamma$. No monomial hypothesis on $p$ or on $x$ is needed, and $\Gamma$ need
not be divisible.
\end{theorem}

\begin{proof}
Evaluation at $0$ is finite. For $x\ne0$ put $\delta=v(x)$ and write
\[
 p=c\,t^\beta(1+u),\qquad x=d\,t^\delta(1+w),\qquad c,d\in\C^\times,
\]
where each of $u,w$ is zero or has strictly positive valuation. Let
$S=\supp(u)\cup\supp(w)\subseteq\Gamma_{>0}$ and $M=M(S)$. The $n$th
evaluation term is
\[
 p^{\binom n2}x^n
 =c^{\binom n2}d^{\,n}\,t^{h_n}\,(1+u)^{\binom n2}(1+w)^n,
 \qquad h_n=\binom n2\beta+n\delta,
\]
and the last two factors are finite integer powers, so
\begin{equation}\label{hol:eq:thetasupport}
 \supp\bigl(p^{\binom n2}x^n\bigr)\subseteq h_n+M,
 \qquad v\bigl(p^{\binom n2}x^n\bigr)=h_n.
\end{equation}

Suppose first that $\delta\ge-N\beta$ for some $N\in\N$. Then
\[
 h_{n+1}-h_n=n\beta+\delta\ge(n-N)\beta>0\qquad(n>N),
\]
so $(h_n)$ is eventually strictly increasing. Its set of values is therefore
well ordered, including the finite initial segment, and each value is taken by
finitely many indices. Lemma~\ref{hol:lem:certificate}\,(b) applied to
\eqref{hol:eq:thetasupport} proves strong summability.

If no such $N$ exists, then $\delta<-N\beta$ for every $N$, so
$h_{n+1}-h_n=n\beta+\delta<0$ for every $n$. By
\eqref{hol:eq:thetasupport} the leading exponents are exactly $h_n$, and they
form a strictly decreasing sequence, so strong summability fails by
Lemma~\ref{hol:lem:nonincreasing}.

Finally, $\delta\ge-N\beta$ for some $N$ says exactly that the image of
$\delta$ in $\Gamma/\Delta_\beta$ is nonnegative: a nonnegative coset either
is positive, whence $\delta>0$, or is zero, whence $\delta\in\Delta_\beta$ and
$\delta\ge-N\beta$ for some $N$; a negative coset lies below every element of
$\Delta_\beta$. The last assertion is the case $\Delta_\beta=\Gamma$.
\end{proof}

This proof controls the whole support, not only the leading valuation. Leading
valuation formulas by themselves do not in general determine a Hahn evaluation
domain, because hidden tails can matter; here the common positive monoid $M$
is the extra certificate that makes the formula exact. Note also that the
entire calculation is carried out in $\Gamma$: nothing is divided by an
integer, and no real-valued surrogate for $v$ is introduced.

\begin{remark}[A second, shorter proof of the order-unit half]
\label{hol:rem:thetasecond}
When $\mu=\beta$ is an order unit, entireness also follows from the
cofinal-valuation criterion Lemma~\ref{hol:lem:certificate}\,(a) alone, and
the routes of sources~08 and~09 are both kept here because they are genuinely
different. For $x\ne0$ with $\alpha=v(x)$, the $n$th evaluated term has
valuation
\begin{equation}\label{hol:eq:thetaval}
 \frac{n(n-1)}2\mu+n\alpha.
\end{equation}
Given a target $\beta'\in\Gamma$, choose ordinary nonnegative integers $M',B'$
with $\alpha\ge-M'\mu$ and $\beta'\le B'\mu$, which is possible exactly
because $\mu$ is an order unit. Then \eqref{hol:eq:thetaval} is at least
$\bigl(\tfrac{n(n-1)}2-M'n\bigr)\mu$, which exceeds $B'\mu$ for all
sufficiently large ordinary $n$. So the term valuations tend cofinally to
$+\infty$ and Lemma~\ref{hol:lem:certificate}\,(a) applies. Only leading term
valuations are used, so unrestricted higher Hahn tails in $p$ do not
invalidate the conclusion. This shorter route is available only in the
order-unit case: at the boundary $v(x)=-N\beta$ with $\beta$ noncofinal, the
term valuations increase strictly but stay trapped below a fixed element of
$\Gamma$, so criterion (a) says nothing and the monoid certificate is
essential.
\end{remark}

\begin{remark}[The order-unit case has an optional rank-one picture]
\label{hol:rem:thetacoarse}
When $\beta$ is an order unit, the real-valued coarsening
$\pi:\Gamma\to\R$ of~\cite[Lemma~9.1]{repoentire}, normalized by
$\pi(\beta)=1$, gives the same field $\K$ a rank-one absolute value.
For $x\ne0$, its theta terms have coarsened valuations
$\binom n2+n\pi(v(x))$, tending to $+\infty$.
This is an optional normed-field picture of the order-unit case,
not a map that replaces the Hahn exponents by their images in $\R$.

Such a coefficientwise replacement need not define a Hahn series.
In $\Q H\oplus\Q\epsilon$ with $H$ dominating all $n\epsilon$,
the coarsening normalized by $H$ kills $\epsilon$.
The valid series $\sum_{n\ge0}t^{n\epsilon}$ would then place infinitely
many coefficients at exponent zero. Coarsening a valuation retains
the field and changes the value map; it is not an embedding into a
complex-coefficient Hahn field with smaller exponent group.
No proof here requires this optional coarsening.
\end{remark}

\subsection{Completion of the sharp criterion}

\begin{proof}[Proof of Theorem~\ref{hol:main:q}]
Assume (i) and suppose (ii) fails. If $v(q)=0$, Corollary
\ref{hol:cor:unitrigid} makes every entire solution a polynomial. If
$v(q)\ne0$ and $\abs{v(q)}$ is not an order unit,
Proposition~\ref{hol:prop:noncofinal} does the same. Either way (i) fails,
which proves (i) $\Rightarrow$ (ii) and supplies the stated uniform
obstruction and degree bound.

Assume (ii). If $v(q)>0$ is an order unit take $p=q$: by
Theorem~\ref{hol:thm:thetadomain}, $\Theta_q\in\E\setminus\K[z]$, and
\eqref{hol:eq:thetahom} is a nonzero linear $q$-difference equation with
polynomial coefficients. If $v(q)<0$ and $-v(q)$ is an order unit take
$p=q^{-1}$, so $v(p)=-v(q)>0$ is an order unit and $\Theta_{q^{-1}}\in\E$ by
the same theorem. Substituting $q^2z$ for $z$ in \eqref{hol:eq:thetahom},
with $p=q^{-1}$ so that $pz\mapsto qz$ and $p^2z\mapsto z$, gives
\begin{equation}\label{hol:eq:inverseq}
 \Theta_{q^{-1}}(q^2z)-(1+q^2z)\Theta_{q^{-1}}(qz)
 +qz\,\Theta_{q^{-1}}(z)=0,
\end{equation}
which uses only nonnegative powers of $\sigma_q$ and has polynomial
coefficients that are not all zero. Both signs of $v(q)$ are covered.
\end{proof}

The partial theta function is used as an established type of construction,
not advertised as a newly discovered special function. Its role is precise: it
realizes the positive side of a necessary and sufficient criterion, including
nonmonomial $q$ and arbitrary valuation rank.

\subsection{Two scales in the same rank-two field}

\begin{example}[One field, opposite answers]\label{hol:ex:ranktwo}
Let $\Gamma=\Q e_0\oplus\Q e_1$ with $e_1\gg e_0>0$, ordered by the highest
nonzero coordinate. Then $e_1$ is an order unit and $e_0$ is not.

For $q=t^{e_0}$, no nonpolynomial strongly entire series satisfies a nonzero
linear $q$-difference equation. The series $\Theta_q$ itself is not entire:
by Theorem~\ref{hol:thm:thetadomain} its exact domain is the valuation ring
obtained by forgetting the $e_0$ coordinate, which contains every $x$ with
$v(x)=-Ne_0$ and excludes $x=t^{-e_1}$.

For $q=t^{e_1}$, $\Theta_q$ is strongly entire, though the value group still
has rank two. The two conclusions occur in the \emph{same} Hahn field: the
distinction is the dilation scale, not the rank of the field.
\end{example}

Under $e_0=1$, $e_1=\omega$ and $t^\gamma=\omega^{-\gamma}$, the first dilation
is multiplication by $\omega^{-1}$ and the second by $\omega^{-\omega}$.
Quadratic multiples of the first valuation never reach the second scale;
quadratic multiples of the second dominate every scale in this workspace. A
fuller table for the same group, including unit dilations and the torsion
exception, appears in Section~\ref{hol:sub:tables}.

\subsection{The coarse boundary really can be included}
\label{hol:sub:coarseboundary}

For the theta series the forward recurrence is $a_{n+1}=p^na_n$, read off from
\eqref{hol:eq:thetafirst}, whose inhomogeneity affects only the index $n=0$
and is removed by Proposition~\ref{hol:prop:inhomogeneous}. Its two slots have
valuations $n\beta$ and $0$, both of which project to $0$ modulo
$\Delta_\beta$, so the projected gap is zero and
$\overline B_{\mathrm{rec}}=\overline B^{\flat}_{\mathrm{rec}}=0$.
Theorem~\ref{hol:thm:coarse} therefore excludes the points of negative
coarsened valuation, while Theorem~\ref{hol:thm:thetadomain} includes
\emph{all} points of zero coarsened valuation. The strict inequality in the
coarsened exclusion theorem cannot be replaced by a weak one in general. This
is the exact mirror image of Section~\ref{hol:sub:expdomain}, where the
uncoarsened bound was attained with its boundary excluded from the domain. The
two boundaries behave differently because one comparison is made in $\Gamma$
and the other in a quotient that has forgotten the scale doing the work.

\section{Torsion parameters and why they are excluded}\label{hol:sec:torsion}

If $q^{h}=1$ then $\sigma_q^{h}-1$ annihilates \emph{every} formal power
series, so the unrestricted definition of linear $q$-difference finiteness is
vacuous for such a parameter. This is not a technical inconvenience removable
by a better proof; it is a genuine counterexample to dropping the nontorsion
hypothesis in Theorem~\ref{hol:main:q}.

The difficulty persists even after excluding operators with zero action. For
$h>1$ and
\[
 f(z)=z^rg(z^{h}),
\]
one has $f(qz)=q^rf(z)$, and the operator $\sigma_q-q^r$ is not zero on all
formal series when $h>1$. A nonpolynomial entire $g$ produces a nonpolynomial
entire $f$, since evaluation of $f$ at $x$ is evaluation of $g$ at $x^{h}$
followed by multiplication by $x^r$. Finite-order dilations therefore have a
real periodic-support exception and are not a limiting case of
Theorem~\ref{hol:main:q}. In particular ``the residue of $q$ is a root of
unity'' and ``$q$ is a root of unity'' are fundamentally different
hypotheses: Example~\ref{hol:ex:unitperiod} is covered by the theorem and this
section is not.

There is nevertheless a meaningful eigenspace statement.

\begin{proposition}[Torsion covariance]\label{hol:prop:torsion}
Let $q\in\C^\times$ have exact finite order $m$ and let $0\le r<m$.
\begin{enumerate}[label=\textup{(\alph*)}]
\item A formal series satisfies $f(qz)=q^rf(z)$ if and only if
$f(z)=z^rG(z^m)$ for a unique $G\in\K[[z]]$.
\item If $G\in\E$ then $f=z^rG(z^m)\in\E$.
\item The converse of \textup{(b)} also holds, so
$f\in\E$ if and only if $G\in\E$, without divisibility of $\Gamma$.
\end{enumerate}
Moreover, every root of unity of $\K$ lies in $\C$, so the hypothesis
$q\in\C^\times$ loses no torsion element of the Hahn field. All parts use only Convention~\ref{hol:conv:group}.
\end{proposition}

\begin{proof}
For (a), coefficient comparison gives $(q^n-q^r)a_n=0$, so only degrees
congruent to $r$ modulo $m$ occur, and $G$ is read off uniquely.

For (b), evaluation of $z^rG(z^m)$ at $x$ is $x^r$ times the family for $G$ at
$x^m$, reindexed; multiplication by the fixed element $x^r$ preserves strong
summability by Lemma~\ref{hol:lem:support}\,(3), and at $x=0$ the evaluation
is finite.

For (c), suppose $f\in\E$ and let $y\in\K^\times$.
Set $\delta=v(y)$, $\gamma=\min(\delta,0)$, and $x=t^\gamma$.
Since $m\ge1$, we have $v(x^m)=m\gamma\le\delta$.
Strong evaluation of $f$ at $x$, restricted to the degrees $r+mn$
and multiplied by $x^{-r}$, gives strong evaluation of $G$ at $x^m$.
Lemma~\ref{hol:lem:inward} then gives strong evaluation at $y$.
At $y=0$ evaluation is finite. No root of $y$ is extracted.

Finally, let $\zeta\in\K$ be a root of unity. Then $v(\zeta)=0$; dividing by
$\res(\zeta)$ gives $1+u$ with $u$ zero or of positive valuation, and
$(1+u)^{m'}=1$ for some $m'\ge1$. If $u\ne0$ then $(1+u)^{m'}-1$ has valuation
$v(u)$ and leading coefficient $m'\lc(u)\ne0$ in characteristic zero, a
contradiction. So $u=0$ and $\zeta=\res(\zeta)\in\C$.
\end{proof}

The converse uses only that the values $m\gamma$ are coinitial in $\Gamma$,
as the explicit choice $\gamma=\min(\delta,0)$ shows. Actual $m$th roots
need not exist in the original Hahn field; this stronger algebraic
requirement is unnecessary for the entireness equivalence.

\section{An extension to mixed differential--dilation equations}
\label{hol:sec:mixed}

For a nonzero noncofinal dilation valuation, derivatives do not repair the
missing-scale problem. The extension is useful when a symbolic system allows
both differentiation and rescaling of the argument.

\begin{lemma}[Polynomial dependence on the index as well as the dilation]
\label{hol:lem:bivariate}
Let $q\in\K^\times$ with $\lambda=v(q)\ne0$ and let
$0\ne P(U,X)\in\K[U,X]$. There are $\beta\in\Gamma$, $k\in\N$ and $N\in\N$
with
\[
 v(P(n,q^n))=\beta+kn\lambda\qquad(n\ge N).
\]
\end{lemma}

\begin{proof}
Write $P(U,X)=\sum_kC_k(U)X^k$. By Lemma~\ref{hol:lem:integer} each nonzero
$C_k(n)$ has the fixed valuation $\beta_k=\beta(C_k)$ for all sufficiently
large $n$. The finitely many summands then have valuations
$\beta_k+kn\lambda$, and the pairwise-order argument of
Lemma~\ref{hol:lem:affine} produces a unique eventual minimizing index, whose
term cannot cancel.
\end{proof}

\begin{theorem}[Mixed rigidity at a noncofinal nonzero scale]
\label{hol:thm:mixed}
Let $\Gamma$ be as in Convention~\ref{hol:conv:group}, let $q\in\K^\times$
with $v(q)\ne0$, and suppose $\abs{v(q)}$ is not an order unit of $\Gamma$. If
a strongly entire $f$ satisfies a nonzero equation
\begin{equation}\label{hol:eq:mixed}
 \sum_{r,\ell}p_{r,\ell}(z)D_z^r\sigma_q^\ell f(z)=0,\qquad p_{r,\ell}\in\K[z],
\end{equation}
where the sum is finite and the coefficients are not all zero, then $f$ is a
polynomial. Each fixed equation has a uniform exterior obstruction and a
uniform degree bound.
\end{theorem}

\begin{proof}
Expand $p_{r,\ell}(z)=\sum_kc_{r,\ell,k}z^k$. With $s=r-k$, the term of
coefficient index $n+s$ contributes to degree $n$ the quantity
$c_{r,\ell,k}\fall{n+s}{r}q^{\ell(n+s)}a_{n+s}$, so the coefficient recurrence
has shift polynomials
\begin{equation}\label{hol:eq:mixedshift}
 P_s(U,X)=\sum_{r-k=s}\sum_\ell c_{r,\ell,k}\,q^{\ell s}\fall{U+s}{r}X^\ell,
\end{equation}
evaluated at $(U,X)=(n,q^n)$. For fixed $s$ and $\ell$, distinct $r$ give
distinct degrees in $U$, so a nonempty group of nonzero coefficients cannot
cancel identically and some $P_s$ is nonzero. Reindex from the smallest active
shift; replacing $U$ by $U-s_0$ and $X$ by $q^{-s_0}X$ preserves polynomiality
and nonzeroness, so the leftmost coefficient is $A_0(n)=P'_0(n,q^n)$ for a
nonzero bivariate $P'_0$, and Lemma~\ref{hol:lem:bivariate} makes it
eventually nonzero.

Lemma~\ref{hol:lem:bivariate} gives eventual affine valuations for all the
nonzero coefficients, so their differences obey a bound $B+Mn\abs{v(q)}$
exactly as in \eqref{hol:eq:linearbound}. Since $\abs{v(q)}$ is not an order
unit, \eqref{hol:eq:dominates} bounds this by one element of $\Gamma$, and
Corollary~\ref{hol:cor:boundedcost} finishes the proof.
\end{proof}

The hypothesis $v(q)\ne0$ is essential to \emph{this} proof. For $v(q)=0$ the
separate pure differential and pure dilation conclusions have been proved
above, but they are not a proof of the mixed conclusion, and no mixed
unit-valuation classification is asserted anywhere in this report; see
Question~\ref{hol:q:mixedunit}.

\section{Surreal workspaces that separate the theories}
\label{hol:sec:surreal}

The results so far are stated in set-sized Hahn fields, which keeps the
foundations transparent and allows a literal surreal reading without
quantifying over a proper class as though it were a set.

\subsection{The workspace without an order unit}

Take
\begin{equation}\label{hol:eq:gammainfty}
 \Gamma_\infty=\bigoplus_{j\ge0}\Q e_j,
 \qquad 0<e_0\ll e_1\ll e_2\ll\cdots,
\end{equation}
ordered so that the highest index with a nonzero coordinate determines the
sign, equivalently $e_{j+1}>Ne_j>0$ for all $j\ge0$ and $N\ge1$. Every element
has only finitely many nonzero coordinates. One may identify $e_j$ with the
surreal $\omega^j$.

The group has no order unit: if $\gamma>0$ has largest active index $j$, then
$e_{j+1}>n\gamma$ for every ordinary $n$. The sequence $(e_n)$ is nevertheless
cofinal in the whole group, so this is a countable cofinality that cannot be
generated by repeatedly adding one positive element. This group and this
distinction appear in the companion entire-function
report~\cite{repoentire} and are background here, not a new observation.

\begin{corollary}[Universal nontorsion dilation rigidity]
\label{hol:cor:universal}
Suppose $\Gamma$ has no order unit. Then every nonpolynomial $f\in\E$
satisfies, simultaneously for all nontorsion $q\in\K^\times$ and all $r\ge0$,
\[
 f(z),f(qz),\ldots,f(q^rz)\quad\text{are linearly independent over }\K(z).
\]
Such fields can contain nonpolynomial strongly entire series, as
\eqref{hol:eq:gammainfty} and Theorem~\ref{hol:thm:explicit} show.
\end{corollary}

\begin{proof}
A linear dependence over $\K(z)$ becomes, after clearing denominators, a
nonzero equation \eqref{hol:eq:qlinear}. Since $\Gamma$ has no order unit,
no $\abs{v(q)}$ is an order unit, so Theorem~\ref{hol:main:q} forces $f$ to be
a polynomial.
\end{proof}

\begin{theorem}[An explicit universally nonholonomic entire series]
\label{hol:thm:explicit}
In $K_{\Gamma_\infty}$ the series
\begin{equation}\label{hol:eq:explicitF}
 F(z)=\sum_{n\ge0}t^{e_n}z^n
     =\sum_{n\ge0}\omega^{-\omega^n}z^n
\end{equation}
is strongly entire and nonpolynomial. It satisfies no nonzero linear
differential equation over $K_{\Gamma_\infty}(z)$ and, for every nontorsion
$q\in K_{\Gamma_\infty}^\times$, no nonzero polynomial-coefficient linear
$q$-difference equation. For every $q$ of nonzero valuation it satisfies no
nonzero mixed equation \eqref{hol:eq:mixed} either.
\end{theorem}

\begin{proof}
Fix $x\ne0$ and write $x=c\,t^\delta(1+w)$ with $c\in\C^\times$ and $w$ zero
or of positive valuation. Choose $M$ with $\delta$ involving only the
coordinates $e_0,\ldots,e_M$. The $n$th evaluation term has leading exponent
$h_n=e_n+n\delta$ and support contained in $h_n+M(\supp w)$. For $n>M$ the
highest coordinate of $h_n$ is $e_n$ with coefficient $1$, so $h_n>0$, and the
highest coordinate of
\[
 h_{n+1}-h_n=e_{n+1}-e_n+\delta
\]
is $e_{n+1}$ with coefficient $1$, so $(h_n)$ is eventually strictly
increasing. Lemma~\ref{hol:lem:certificate}\,(b) proves strong summability; at
$x=0$ only the constant term remains. All coefficients are nonzero, so $F$ is
nonpolynomial.

Alternatively, for any $\beta\in\Gamma_\infty$ choose $j$ larger than every
active coordinate index of $\delta$ and of $\beta$; for $n>j$ the largest
active coordinate of $h_n-\beta$ is $n$ with coefficient $1$, so $h_n>\beta$,
and Lemma~\ref{hol:lem:certificate}\,(a) applies for arbitrary $w$.
Cofinal leading valuations control the whole support, since every exponent
of a Hahn term is at least its leading exponent.

The differential conclusion is Theorem~\ref{hol:main:ode}. Since
$\Gamma_\infty$ has no order unit, Theorem~\ref{hol:main:q} --- equivalently
Corollary~\ref{hol:cor:universal} --- excludes every nontorsion dilation
equation, and Theorem~\ref{hol:thm:mixed} excludes the mixed equations at
every $q$ with $v(q)\ne0$. Rational coefficient equations reduce to the
polynomial-coefficient case by clearing denominators.
\end{proof}

The universal quantifier over $q$ matters. The parameter may be chosen after
$F$ is fixed and may itself have infinite support; each individual $v(q)$
still has finitely many coordinates in $\Gamma_\infty$, hence is noncofinal.
So the theorem excludes a linear relation for \emph{each} such parameter, not
merely for an advance list of simple dilations. The series has a simple
explicit coefficient rule, and this is a statement about finite linear
annihilating equations, not a claim that the rule is uncomputable or lacks a
finite description.

\subsection{Why the same formula is not entire on all surcomplex numbers}
\label{hol:sub:notallsurcomplex}

Let $e_\ast=\omega^\omega$, which dominates every integer multiple of every
positive element of $\Gamma_\infty$, and enlarge the group to
$\Gamma_\infty+\Q e_\ast$. At the surcomplex point $x=\omega^{e_\ast}
=t^{-e_\ast}$ the evaluated leading exponents are $e_n-ne_\ast$, whose
successive differences $e_{n+1}-e_n-e_\ast$ are negative. The defining family
is therefore not strongly summable, by
Lemma~\ref{hol:lem:nonincreasing}.

This explicit obstruction prevents the fixed-workspace theorem from being
misread as an assertion of a nonpolynomial power series entire on all of
$\No[i]$. The example illustrates two independent restrictions: a workspace
can be large enough that no single dilation scale is cofinal, yet small enough
to admit a countable sequence of coefficient scales escaping all its
arguments, and enlarging it further can destroy the very entireness on which
the theorem operates. The companion report studies this kind of boundary in
general~\cite{repoentire}; the calculation here is included only to state the
domain of the present example unambiguously, and no new universal
scalar-extension theorem is claimed.

\subsection{Two classification tables}\label{hol:sub:tables}

Let $\Gamma=\Q H\oplus\Q\epsilon$ with $H$ dominating every $n\epsilon$, so
that $H$ is an order unit and $\epsilon$ is not.

\begin{center}
\begin{tabularx}{\textwidth}{@{}l l X@{}}
\toprule
Dilation $q$ & Its valuation & Nonpolynomial entire solution of a nonzero
linear $q$-difference equation?\\
\midrule
$t^H$ & $H$ & Yes: $\Theta_{t^H}$, by \eqref{hol:eq:thetahom}.\\
$t^{-H}$ & $-H$ & Yes: $\Theta_{t^{H}}$, by \eqref{hol:eq:inverseq}.\\
$t^\epsilon$ & $\epsilon$ & No: positive, but noncofinal.\\
$1+t^H$ & $0$ & No: an infinite-order unit.\\
$2$ & $0$ & No: complex size does not change the Hahn valuation.\\
$-1$ & $0$ & Outside the theorem: a torsion parameter, with
periodic-support exceptions.\\
\bottomrule
\end{tabularx}
\end{center}

\begin{center}
\small
\begin{tabularx}{\textwidth}{@{}l l l X@{}}
\toprule
Workspace group & Parameter & Solution exists? & Reason\\
\midrule
$\Q$ & $v(q)\ne0$ & Yes & Every positive nonzero scale is an order unit.\\
$\Q$ & $v(q)=0$, nontorsion & No & Unit-orbit valuations are eventually
periodic.\\
$\Q e_0\oplus\Q e_1$, $e_1\gg e_0$ & $v(q)=e_0$ & No & A lower scale is not
cofinal.\\
the same group & $v(q)=e_1$ & Yes & The top scale is an order unit.\\
$\bigoplus_{n\ge0}\Q e_n$ & any nontorsion $q$ & No & No scale is an order
unit, although $F$ is entire.\\
\bottomrule
\end{tabularx}
\end{center}

In every row, nonpolynomial strongly entire D-finite series are impossible.
Torsion parameters are deliberately absent from the second table: their
annihilating operator may be an identity on every series. Finite valuation
rank alone is therefore not the criterion, and the distinction persists for
nonmonomial dilations, because the theorem uses the valuation rather than a
requirement that $q$ be a monomial.

\begin{corollary}[Detecting order units by functional equations]
\label{hol:cor:detect}
For a nonzero set-sized ordered abelian group $\Gamma$ the following are
equivalent:
\begin{enumerate}[label=\textup{(\roman*)}]
\item $\Gamma$ has an order unit;
\item some nontorsion $q\in\K^\times$ admits a nonpolynomial strongly entire
solution of a nonzero polynomial-coefficient linear $q$-difference equation.
\end{enumerate}
\end{corollary}

\begin{proof}
Theorem~\ref{hol:main:q} gives (ii) $\Rightarrow$ (i). For the converse take
$q=t^\mu$ for an order unit $\mu$; then $v(q)=\mu$ is an order unit and $q$ is
nontorsion.
\end{proof}

This recovers a property of the ordered group from a global functional
equation existence question. The workspace \eqref{hol:eq:gammainfty} shows why
the existence of nonpolynomial entire functions is by itself strictly weaker
than the existence of an order unit.

\subsection{Several meanings of transcendence, kept apart}

For the explicit $F$ the established conclusions are: transcendence over
$K_{\Gamma_\infty}(z)$ by Corollary~\ref{hol:cor:algebraic}; linear
differential independence of every finite family of derivatives; and linear
independence along every single nontorsion geometric dilation orbit. These are
substantial restrictions, and they do \emph{not} imply that no polynomial
relation $P(z,F,D_zF,\ldots,D_z^RF)=0$ exists. Nonlinear differential
equations lie outside these proofs, and a relation involving $F(q_1z)$ and
$F(q_2z)$ for multiplicatively independent parameters is not covered merely by
applying the single-$q$ theorem twice.

The nonlinear part adds one more conclusion by a separate proof: $F$ and
$D_zF$ are algebraically independent over $K_{\Gamma_\infty}(z)$
(Corollary~\ref{hol:nl:cor:infiniterank}), so every nonlinear relation with
$R\le1$ is excluded. Relations with $R\ge2$ are still not excluded unless
their residue corner polynomial is nonzero, and no statement about several
dilations follows.

\section{Coefficient fields, degenerate value groups, characteristic}
\label{hol:sec:fields}

The one-variable proofs work over any characteristic-zero coefficient field
with trivial valuation: the root-of-unity argument uses only finitely many
ordinary polynomial roots and needs no algebraic closedness, and the
integer-evaluation lemma needs only residue characteristic zero. All theorems
also hold over $\R((t^\Gamma))$, the generic-line proof of
Theorem~\ref{hol:main:multi} using uncountability of $\R$ in place of
uncountability of $\C$. The results therefore concern surreal as well as
surcomplex workspaces. More generally the core versions work over
$k((t^\Gamma))$ for any characteristic-zero field $k$, and, as
Section~\ref{hol:sub:divisibility} records, without divisibility of $\Gamma$.

No statement here is a theorem for arbitrary valued fields of positive residue
characteristic. In that setting Lemma~\ref{hol:lem:integer} fails and the
valuations of near-torsion geometric powers can behave differently, so any
comparison with the ultrametric literature must keep the coefficient field and
the residue characteristic visible. Example~\ref{hol:nl:ex:charp}, from
source~10, shows that characteristic zero is needed even for the single
linear equation $D_zf=0$: over $\mathbb F_p((t^\Q))$ a nonpolynomial strongly
entire series has zero derivative.

The nonlinear part (Convention~\ref{hol:nl:conv:field}) is proved by its
source over an arbitrary field of characteristic zero. That is a statement
about Sections~\ref{hol:nl:sec:setting}--\ref{hol:nl:sec:consequences} only;
the generality of the linear theorems is governed by the first paragraph of
this section and its own argument.

We assumed $\Gamma\ne0$ in order to speak of positive exterior scales. When
$\Gamma=0$ the field is $\C$ with the trivial valuation, strong evaluation at
$x=1$ already permits only finitely many nonzero coefficients because every
nonzero term has support $\{0\}$, and entire series are polynomials for that
immediate support reason.

\section{Quantitative certificates and symbolic representation}
\label{hol:sec:certificates}

\subsection{What one can extract from an annihilating equation}

The classification has a direct representational consequence. A system whose
nonpolynomial special functions are represented only by finite linear
differential equations cannot cover the nonpolynomial strong entire functions
of a Hahn workspace. Allowing linear dilation equations enlarges the class only
when a cofinal dilation scale is available. In the no-order-unit example, an
exact coefficient rule such as \eqref{hol:eq:explicitF} lies outside both pure
equation formats.

This does not rule out exact symbolic computation with such functions. It
specifies what a faithful representation must keep separate: the
coefficient-generation rule, the exponent group, the strong evaluation domain,
and the type of finite equation being asserted. A field-independent label such
as ``entire theta function'' discards relevant information; for $\Theta_p$ the
proof certificate is not merely $v(p)>0$, it includes that $v(p)$ is an order
unit of the selected group.

For a differential operator the reduction yields a concrete certificate.
Compute the finite shift polynomials $Q_s$ of \eqref{hol:eq:odecoeff} and
shift the least active index to zero. For each resulting $P_j$ record its
least coefficient valuation $\beta_j$ and the corresponding nonzero complex
residue polynomial $\overline{P_j}$. Record an integer $N$ beyond their
nonnegative integer roots and beyond the recurrence's starting index. Record
the threshold $B$ of Proposition~\ref{hol:prop:precursive}, or
\eqref{hol:eq:refinedthreshold} when $\Gamma$ is divisible. The escape
statements then certify that every formal solution strongly evaluable at some
$x\ne0$ with $v(x)\le-B$ has degree less than $N$.

For a valuation-zero dilation the same plan uses \eqref{hol:eq:qrec}: if
$\res(q)$ has infinite order use the minimum coefficient valuation, and if it
has order $h$ use \eqref{hol:eq:unitexpand}--\eqref{hol:eq:residuepoly} on
each residue class. For a nonzero-valuation dilation, finitely many affine
comparisons select the eventual winner of each $R_h(q^n)$; project their
constants modulo $\Delta_{\abs{v(q)}}$ to obtain the coarsened bound
\eqref{hol:eq:coarsebound0}, or \eqref{hol:eq:coarsebound} when $\Gamma$ is
divisible. When the parameter scale is noncofinal one can instead display
$\gamma\gg\abs{v(q)}$ and the explicit bad argument $t^{-(B+\gamma)}$ of
Proposition~\ref{hol:prop:noncofinal}.

These are mathematical certificate schemas relative to exact coefficient data
and order comparisons. They are \emph{not} a terminating algorithm for
arbitrary encoded surreals. Deciding whether an arbitrary complex residue is a
root of unity, comparing arbitrary surreal values, extracting arbitrary
supports, or deciding whether an abstract scale is cofinal can require
information not effectively available in a chosen representation, and none of
this has been made computable by the proofs.

\subsection{A valuation recurrence with zero coefficients allowed}

\begin{example}[Jump lengths, not just jumps]\label{hol:ex:sparse}
Suppose a sequence that is not eventually zero satisfies
\[
 t^Aa_n+(n-2)t^Ba_{n+1}+t^Ca_{n+3}=0\qquad(n\ge0),
\]
with $A,B,C\in\Gamma$. The middle coefficient vanishes at $n=2$, a finite
exceptional index; beyond it the three valuations are $A,B,C$. The escape
chain may move forward by one or by three, and no assumption that all $a_n$
are nonzero is needed. Corollary~\ref{hol:cor:boundedcost} excludes the region
$v(x)\le-\max\{A-B,\,A-C,\,0\}$ under
Convention~\ref{hol:conv:group} alone. \emph{If $\Gamma$ is divisible},
Corollary~\ref{hol:cor:periodic} gives the sharper excluded region
\[
 v(x)\le-\max\Bigl\{A-B,\ \frac{A-C}{3}\Bigr\}.
\]
The maximum must be divided by the actual jump length: replacing every jump by
one need not preserve the intended threshold, especially when the valuation
gaps have different signs. This is a worked instance of item (1) of
Section~\ref{hol:sub:divisibility}.
\end{example}

\subsection{Current Lean coverage and remaining dependency order}

The ledger \texttt{docs/FORMALIZATION.md} records the formalized escape
mechanism in \path{Surreal/HahnSeries/EscapeChain.lean}:
the nonincreasing-order obstruction, escape step and infinite chain,
evaluation exclusion, and the uniform bounded-cost corollary. It also covers
the univariate strong-domain definitions and the formal exponential
identity with its exact domain and attained exclusion boundary.
Part~(2) of Lemma~\ref{hol:lem:support} is mapped separately to
\path{Surreal/HahnSeries/NeumannWords.lean}.

The next dependencies are integer valuation stabilization and the unit-orbit
residue-polynomial argument, followed by the operator-to-recurrence
conversions and the differential and dilation classifications.
The refined periodic threshold, the complete support certificates, the
new inward-stability lemma, the strengthened torsion equivalence and the
generic-line extension remain unformalized. The several-variable proof also
needs countable polynomial avoidance and restriction of a rational
differential system.

Nothing in the nonlinear part,
Sections~\ref{hol:nl:sec:setting}--\ref{hol:nl:sec:consequences}, is
formalized. The build and axiom audit cover only imported declarations
recorded in the ledger. The Python checks are finite checks of formulas and examples; they
do not establish any remaining infinite statement.

%======================================================================
% Nonlinear part: source 10 (nonlinear rigidity). Labels use hol:nl:.
%======================================================================
\section{Nonlinear equations: finite initial polynomials}
\label{hol:nl:sec:setting}

This section and the next three contain the material of source~10. They
answer the first-order case of Question~\ref{hol:q:nonlinear}, which
sources~08 and~09 had left open, together with a nondegenerate higher-order
class and positive-weight Euler equations in several variables. The method is
not the escape chain of Section~\ref{hol:sec:escape}. The warning in
Remark~\ref{hol:rem:transcendence} stands: a nonlinear equation gives
coefficient convolutions, not a bounded-step linear recurrence. Instead, a
whole scaled series is reduced to a finite polynomial at one large scale, and
an exact leading-term computation gives the contradiction.

\subsection{Coefficient field, conventions and the merge's choices}
\label{hol:nl:sub:conventions}

\begin{convention}[Coefficient field of the nonlinear part]
\label{hol:nl:conv:field}
In Sections~\ref{hol:nl:sec:setting}--\ref{hol:nl:sec:consequences}, $k$ is an
arbitrary field of characteristic zero with the trivial valuation, $\Gamma$ is
as in Convention~\ref{hol:conv:group}, and
\[
 \KK=k((t^\Gamma)).
\]
The valuation $v$, leading coefficient $\lc$, residue $\res$, strong
summability, $\Dom(f)$ and entireness are defined exactly as in
Section~\ref{hol:sub:conventions} and Definition~\ref{hol:def:entire}, with
$k$ in place of $\C$, except that, following source~10, $\res(a)$ is defined
for every $a$ of nonnegative valuation as its coefficient at exponent zero, so
that $\res(a)=0$ when $v(a)>0$. Write $\EK$ for the strongly entire series in
$\KK[[z]]$. In several variables, strong entireness of
$f=\sum_\alpha a_\alpha x^\alpha$ means strong summability of the \emph{joint}
family $(a_\alpha b^\alpha)_{\alpha\in\N^m}$ at every $b\in\KK^m$. For $k=\C$
one has $\KK=\K$ and $\EK=\E$. The derivative $D_z$ kills $\KK$, and
$f'=D_zf$; no scalar surreal derivation, analytic continuation or
fine-topology limit is involved.
\end{convention}

\paragraph{Why the coefficient field differs.}
Source~10 proves every statement of the nonlinear part over an arbitrary
characteristic-zero field~$k$, and it uses neither algebraic nor real
closedness. Its hypothesis on the coefficient field is therefore weaker than
the hypothesis $k=\C$ under which
Theorems~\ref{hol:main:ode}--\ref{hol:main:multi} and the rest of
Sections~\ref{hol:sec:intro}--\ref{hol:sec:certificates} are stated. This
generality is recorded for the nonlinear part only. It is \emph{not}
transferred to the linear theorems: they remain stated over~$\C$, and their
possible extension to other coefficient fields is the separate remark of
Section~\ref{hol:sec:fields}, resting on its own argument. In the other
direction, the nonlinear part borrows from Section~\ref{hol:sec:support} only
facts about well-ordered subsets of~$\Gamma$ and finite coefficient sums
(Lemmas~\ref{hol:lem:support} and~\ref{hol:lem:certificate}), whose proofs do
not use any property of~$\C$; source~10 proves the facts it needs over~$k$
itself, except the positive-support lemma used in
Theorem~\ref{hol:nl:thm:nouniform}, which it imports as standard. The group $\Gamma$ is not assumed divisible, and no order unit, no
countable cofinal subset and no real-valued norm is used anywhere in the
nonlinear part.

\paragraph{Renamed symbols.}
Several symbols of source~10 would collide with symbols of
Sections~\ref{hol:sec:intro}--\ref{hol:sec:certificates}; they are renamed
here, in the new material only.
\begin{center}
\small
\begin{tabularx}{\textwidth}{@{}l l X@{}}
\toprule
Source~10 & Here & Reason\\
\midrule
$D=\max_\nu d_\nu$, $W$ & $d_P$, $w_P$ & $D_z$ is the derivative.\\
$H_P$ (residue corner polynomial) & $\chi_P$ & $H$ is an element of $\Gamma$
in Sections~\ref{hol:sec:support}--\ref{hol:sec:surreal}.\\
$B_P$ (degree cutoff, in $\Q$) & $\kappa_P$ & $B$, $B_L$, $B_{\mathrm{rec}}$
are elements of $\Gamma$.\\
$B$ (set of right sides) & $\mathcal R$ & The same.\\
$r$, $r_0$ (scale, ``radius'') & $\delta$, $\delta_0$ & $r$ indexes
derivative orders in \eqref{hol:eq:dlinear}; $t^{-\delta}$ is the notation
of Corollary~\ref{hol:cor:boundedcost}.\\
$h_r$, $p_r$ (Gauss value, initial polynomial) & $\gamma_\delta$,
$\phi_\delta$ & $h$ is a torsion order, $p_r$ an operator coefficient.\\
$M$ (an integer bound) & $n_\ast$ & $M(S)$ is a support monoid.\\
$E_q$ (Euler operator) & $\vartheta_q$ & $E_\eta$ is the formal exponential
of Proposition~\ref{hol:prop:expdomain}.\\
Theorems~A, B, C & Theorems~\ref{hol:nl:main:first}, \ref{hol:nl:thm:euler},
\ref{hol:nl:thm:corner} & Theorems~A--C of this report are different
theorems.\\
\bottomrule
\end{tabularx}
\end{center}
The shipped audit files \path{10-nonlinear-rigidity-PROOF_AUDIT.md} and
\path{10-nonlinear-rigidity-SOURCES_AND_SCOPE.md} keep the source's
notation and numbering.

\paragraph{Words with two meanings.}
\emph{Escape} keeps the meaning of Section~\ref{hol:sec:escape} (a chain
through a linear recurrence); source~10's lemma ``escape of active degree''
is a different statement and is called Lemma~\ref{hol:nl:lem:active} here.
An \emph{exclusion bound} is, as in Section~\ref{hol:sec:escape}, an element
$\delta_0\in\Gamma$ such that evaluation fails at the arguments $t^{-\delta}$
with $\delta\ge\delta_0$; the source's word ``radius'' is not used. A
\emph{corner} is always the finite set of terms of
Section~\ref{hol:nl:sub:corner}, never a full differential Newton polygon.
The \emph{residue polynomials} $\overline P$ and $R_r$ of
Section~\ref{hol:sec:orbits} are unrelated to the residue corner polynomial
$\chi_P$.

\paragraph{Results printed once.}
Source~10 reproves several facts that Sections~\ref{hol:sec:support}
and~\ref{hol:sec:surreal} already contain. They are printed once and cited:
the support calculus inside Lemma~\ref{hol:nl:lem:algebra}
(Lemma~\ref{hol:lem:support}\,(1),(3) and Proposition~\ref{hol:prop:closure});
its ``cofinal escape of coefficient minima'' lemma, which is
Lemma~\ref{hol:lem:certificate}\,(a); the positive-support geometric-series
lemma, which source~10 imports and Lemma~\ref{hol:lem:support}\,(2) proves;
the surreal embedding, which is Section~\ref{hol:sub:embedding} with the same
monomial convention $t^\gamma\mapsto\omega^{-\gamma}$; strong entireness of
the series \eqref{hol:eq:explicitF}, which is Theorem~\ref{hol:thm:explicit};
and the remark on $\Gamma=0$, which is in Section~\ref{hol:sec:fields}.
Everything else in source~10 is printed in
Sections~\ref{hol:nl:sec:setting}--\ref{hol:nl:sec:consequences}, with its
proofs.

\subsection{The coefficient-family algebra}\label{hol:nl:sub:algebra}

For a finite tuple $X=(X_1,\ldots,X_m)$ of indeterminates put
\[
 \Sfam_m=\Bigl\{F=\sum_\alpha b_\alpha X^\alpha\in\KK[[X]]:
 (b_\alpha)_\alpha\text{ is strongly summable}\Bigr\},
\]
and let $\Sfam_m^+$ be the subset whose coefficients all have nonnegative
valuation. The condition concerns the family of coefficients, not convergence
in the $X$-adic topology.

\begin{lemma}[Hahn-family algebra and its residue]\label{hol:nl:lem:algebra}
$\Sfam_m$ is a $\KK$-subalgebra of $\KK[[X]]$, stable under the formal partial
derivatives and under the Euler operators $\sum_\nu q_\nu X_\nu\partial_\nu$
with $q_\nu\in k$. Coefficientwise residue is a ring homomorphism
\[
 \rho:\Sfam_m^+\longrightarrow k[X],\qquad
 \rho\Bigl(\sum_\alpha b_\alpha X^\alpha\Bigr)=\sum_\alpha\res(b_\alpha)X^\alpha,
\]
whose values are polynomials, not arbitrary formal power series; $\rho$
commutes with these differential operators.
\end{lemma}

\begin{proof}
Closure under sums, scalar multiples and products is the support
calculation of Lemma~\ref{hol:lem:support}\,(1) and~(3), as in the proof of
Proposition~\ref{hol:prop:closure}. For strongly summable $(b_\alpha)$ and
$(c_\epsilon)$ a fixed exponent has finitely many representations as a sum of
two support exponents, each carried by finitely many indices, so the double
family $(b_\alpha c_\epsilon)$ is strongly summable, and regrouping by
$\alpha+\epsilon$ gives the Cauchy product. Source~10 proves this over $k$ by
the same two-coordinate sequence argument; like the Hahn-field construction
itself, that argument uses ordinary set-theoretic choice. A formal derivative
or an Euler operator with coefficients in $k$ reindexes a subfamily and
multiplies its members by elements of the trivially valued field $k$; it
creates no new support exponent. For $F\in\Sfam_m^+$ the exponent zero lies
in only finitely many coefficient supports, so $\rho(F)$ is a polynomial. Each
$X$-coefficient of a product is a finite sum of products of coefficients, the
residue is additive and multiplicative on elements of nonnegative valuation,
and the local finiteness just proved justifies every interchange; so $\rho$ is
a ring homomorphism. Differentiation commutes with residue coefficientwise.
\end{proof}

For $0\ne F=\sum_\alpha b_\alpha X^\alpha\in\Sfam_m$ put
\[
 \gamma(F)=\min_{b_\alpha\ne0}v(b_\alpha),\qquad
 \inpoly(F)=\rho\bigl(t^{-\gamma(F)}F\bigr).
\]
The minimum exists because it is the least element of the nonempty
well-ordered union of supports. Then $t^{-\gamma(F)}F\in\Sfam_m^+$ and the
\emph{initial polynomial} $\inpoly(F)$ is nonzero.

\begin{corollary}[Gauss multiplicativity]\label{hol:nl:cor:gauss}
For nonzero $F,G\in\Sfam_m$,
\[
 \gamma(FG)=\gamma(F)+\gamma(G),\qquad
 \inpoly(FG)=\inpoly(F)\,\inpoly(G).
\]
\end{corollary}

\begin{proof}
Normalize both factors and apply Lemma~\ref{hol:nl:lem:algebra}. The product
of the two nonzero residues is nonzero because $k[X]$ is an integral domain,
so some coefficient of $t^{-\gamma(F)-\gamma(G)}FG$ has valuation exactly zero.
\end{proof}

\subsection{Scaled initial polynomials}\label{hol:nl:sub:scaled}

Let $f=\sum_{n\ge0}a_nz^n\in\KK[[z]]$ and $\delta\in\Gamma$ with
$t^{-\delta}\in\Dom(f)$. Then
\[
 F_\delta(X)=f(t^{-\delta}X)=\sum_{n\ge0}a_nt^{-n\delta}X^n\in\Sfam_1.
\]
For $f\ne0$ put
\begin{equation}\label{hol:nl:eq:gaussdata}
 \gamma_\delta=\min_{a_n\ne0}\bigl(v(a_n)-n\delta\bigr),\qquad
 g_\delta=t^{-\gamma_\delta}F_\delta,\qquad
 \phi_\delta=\rho(g_\delta),\qquad
 N_\delta=\deg\phi_\delta.
\end{equation}
The \emph{active} indices are the $n$ attaining $\gamma_\delta$; there are
finitely many, $N_\delta$ is the largest, and $A_\delta=\lc(\phi_\delta)\in
k^\times$.

The element $\gamma_\delta$ is a \emph{coefficient Gauss value}. It need not
equal $v(f(t^{-\delta}))$: $\phi_\delta(1)$ can vanish although
$\phi_\delta\ne0$. The proofs below use identities in $\Sfam_1$ and in $k[X]$
only, never an identification of these two quantities, and never replace the
initial polynomial by its value at~$1$.

\begin{lemma}[Exact derivative scaling]\label{hol:nl:lem:derivatives}
For every $j\ge0$,
\begin{equation}\label{hol:nl:eq:scalederivative}
 (D_z^jf)(t^{-\delta}X)=t^{\gamma_\delta+j\delta}\,g_\delta^{(j)}(X),
 \qquad \rho\bigl(g_\delta^{(j)}\bigr)=\phi_\delta^{(j)}.
\end{equation}
If $N_\delta\ge j$, then $\phi_\delta^{(j)}$ has degree $N_\delta-j$ and
leading coefficient $A_\delta\fall{N_\delta}{j}$.
\end{lemma}

\begin{proof}
The scaling identity is a formal coefficient computation:
$D_X^jF_\delta(X)=t^{-j\delta}(D_z^jf)(t^{-\delta}X)$. The residue identity is
Lemma~\ref{hol:nl:lem:algebra}. The falling factorial is a nonzero element of
$k$ for $N_\delta\ge j$ because $k$ has characteristic zero.
\end{proof}

\begin{lemma}[Large active degrees]\label{hol:nl:lem:active}
Let $f$ be nonpolynomial, let $n_\ast\in\N$, and choose $m>n_\ast$ with
$a_m\ne0$. There is $\delta_1\in\Gamma_{>0}$ such that, whenever
$\delta\ge\delta_1$ and $t^{-\delta}\in\Dom(f)$, every active index exceeds
$n_\ast$. Moreover $\gamma_\delta\le v(a_m)-m\delta$.
\end{lemma}

\begin{proof}
Choose $\delta_1>0$ with
\begin{equation}\label{hol:nl:eq:lowindices}
 (m-n)\delta_1>v(a_m)-v(a_n)\qquad(0\le n\le n_\ast,\ a_n\ne0).
\end{equation}
There are finitely many inequalities, and they need no division in $\Gamma$:
take $\delta_1$ positive and larger than the absolute values of all right-hand
sides, since each multiplier $m-n$ is a positive integer. They persist for
$\delta\ge\delta_1$ and say that each such low-index term has scaled valuation
$v(a_n)-n\delta$ strictly larger than that of the $m$th term. The last
assertion is the definition of $\gamma_\delta$.
\end{proof}

This is not a limit $N_\delta\to\infty$ along a countable cofinal sequence,
which need not exist; it concerns a final segment of $\Gamma_{>0}$ and only
those scales at which strong evaluation is available. It is also unrelated to
the escape chain of Lemma~\ref{hol:lem:escape}: no recurrence is used.

\subsection{The finite corner of a differential polynomial}
\label{hol:nl:sub:corner}

Let $P$ be a nonzero differential polynomial of order at most $s$, written in
combined-monomial form
\begin{equation}\label{hol:nl:eq:P}
 P(z,Y_0,\ldots,Y_s)=\sum_{\nu\in J}c_\nu z^{k_\nu}\prod_{j=0}^sY_j^{m_{\nu j}},
\end{equation}
with $J$ finite, every $c_\nu\in\KK^\times$, and identical monomials already
combined. Write $P[f]=P(z,f,D_zf,\ldots,D_z^sf)$. Define the integers
\begin{equation}\label{hol:nl:eq:dw}
 d_\nu=\sum_{j=0}^sm_{\nu j},\qquad
 w_\nu=k_\nu-\sum_{j=0}^sj\,m_{\nu j},\qquad
 d_P=\max_\nu d_\nu,\qquad
 w_P=\max_{d_\nu=d_P}w_\nu,
\end{equation}
and
\[
 \mathcal C_P=\{\nu:d_\nu=d_P,\ w_\nu=w_P\},\qquad
 \beta_P=\min_{\nu\in\mathcal C_P}v(c_\nu),\qquad
 \bar c_\nu=\res\bigl(t^{-\beta_P}c_\nu\bigr)\in k\quad(\nu\in\mathcal C_P).
\]
Thus $(d_P,w_P)$ is chosen lexicographically: first the largest total
differential degree, then the largest shift $w$. This is a finite corner
selection, not the construction of a differential Newton polygon. The
\emph{residue corner polynomial} and the \emph{full corner polynomial} are
\begin{equation}\label{hol:nl:eq:chiI}
 \chi_P(T)=\sum_{\nu\in\mathcal C_P}\bar c_\nu
 \prod_{j=0}^s\fall Tj^{m_{\nu j}}\in k[T],\qquad
 I_P(T)=\sum_{\nu\in\mathcal C_P}c_\nu
 \prod_{j=0}^s\fall Tj^{m_{\nu j}}\in\KK[T].
\end{equation}
In general $\chi_P$ can vanish although some $\bar c_\nu\ne0$, because
different products of falling factorials can be linearly dependent. If
$\chi_P\ne0$, then $I_P\ne0$, and
\begin{equation}\label{hol:nl:eq:rootcontainment}
 I_P(n)=0\ \Longrightarrow\ \chi_P(n)=0\qquad(n\in\N),
\end{equation}
because $t^{-\beta_P}I_P$ has coefficients of nonnegative valuation and residue
commutes with substitution of an ordinary integer.

\begin{lemma}[The exact corner identity]\label{hol:nl:lem:corner}
Let $\phi\in k[X]$ have degree $N\ge s$ and leading coefficient $A\ne0$, and
put
\[
 \Omega_\phi(X)=\sum_{\nu\in\mathcal C_P}\bar c_\nu X^{k_\nu}
 \prod_{j=0}^s\bigl(\phi^{(j)}(X)\bigr)^{m_{\nu j}}.
\]
Then $\deg\Omega_\phi\le d_PN+w_P$ and
\begin{equation}\label{hol:nl:eq:corneridentity}
 [X^{d_PN+w_P}]\,\Omega_\phi=A^{d_P}\chi_P(N).
\end{equation}
In particular $\chi_P(N)\ne0$ implies $\Omega_\phi\ne0$.
\end{lemma}

\begin{proof}
Each corner summand has degree at most
$k_\nu+\sum_jm_{\nu j}(N-j)=d_PN+w_P$, and its coefficient in that degree is
$\bar c_\nu A^{d_P}\prod_j\fall Nj^{m_{\nu j}}$. Summing gives the formula.
\end{proof}

Only finitely many nonnegative integers are roots of a nonzero $\chi_P$, even
when $k$ is not algebraically closed. The active degree is an ordinary
integer, not an element of $\Gamma$, and no root extraction or division in
$\Gamma$ occurs in this step.

\section{First-order algebraic differential equations}\label{hol:nl:sec:first}

\subsection{A finite certificate excluding all large scales}

The following theorem carries the proof of Theorem~\ref{hol:nl:main:first}
and is also a practical audit object: its scale inequalities use only finitely
many coefficients of the solution. It applies to equations of any order.

\begin{theorem}[Nonlinear finite exclusion certificate]
\label{hol:nl:thm:certificate}
Let $P$ be as in \eqref{hol:nl:eq:P} with $d_P\ge1$ and $\chi_P\ne0$, and let
$f=\sum_na_nz^n$ be a nonpolynomial formal solution of $P[f]=0$. Choose an
integer $n_\ast\ge s$ at least as large as every nonnegative integer root of
$\chi_P$, and choose $m>n_\ast$ with $a_m\ne0$ and
\begin{equation}\label{hol:nl:eq:slope}
 (d_P-d_\nu)m+w_P-w_\nu>0\qquad(d_\nu<d_P).
\end{equation}
Every $\delta\in\Gamma_{>0}$ satisfying the finitely many strict inequalities
\begin{align}
 (m-n)\delta&>v(a_m)-v(a_n)
 &&(0\le n\le n_\ast,\ a_n\ne0),\label{hol:nl:eq:cert1}\\
 (w_P-w_\nu)\delta&>\beta_P-v(c_\nu)
 &&(d_\nu=d_P,\ w_\nu<w_P),\label{hol:nl:eq:cert2}\\
 \bigl((d_P-d_\nu)m+w_P-w_\nu\bigr)\delta
 &>\beta_P-v(c_\nu)+(d_P-d_\nu)v(a_m)
 &&(d_\nu<d_P)\label{hol:nl:eq:cert3}
\end{align}
has $t^{-\delta}\notin\Dom(f)$. Such $m$ and $\delta$ exist, and the
inequalities persist when $\delta$ increases. Consequently $\Dom(f)$ omits
$t^{-\delta}$ for every $\delta$ in a final segment of $\Gamma_{>0}$, and every
strongly entire solution of $P[f]=0$ is a polynomial.
\end{theorem}

\begin{proof}
An integer $n_\ast$ exists because $\chi_P\ne0$. For $d_\nu<d_P$ the
coefficient $d_P-d_\nu$ of $m$ in \eqref{hol:nl:eq:slope} is a positive
integer, so all sufficiently large $m$ satisfy it, and a nonpolynomial $f$ has
nonzero coefficients of arbitrarily large index. Every multiplier on the left
of \eqref{hol:nl:eq:cert1}--\eqref{hol:nl:eq:cert3} is then a positive
integer, so any $\delta$ that is positive and larger than the absolute values
of all right-hand sides satisfies them, as does every larger $\delta$.

Suppose $t^{-\delta}\in\Dom(f)$ for such a $\delta$, and form
$\gamma=\gamma_\delta$, $g=g_\delta$, $\phi=\phi_\delta$, $N=N_\delta$ as in
\eqref{hol:nl:eq:gaussdata}. By \eqref{hol:nl:eq:cert1} and
Lemma~\ref{hol:nl:lem:active},
\begin{equation}\label{hol:nl:eq:hbound}
 N>n_\ast,\qquad \gamma\le v(a_m)-m\delta.
\end{equation}
All derivatives of $g$ lie in $\Sfam_1^+$. Substituting $z=t^{-\delta}X$, a
ring homomorphism $\KK[[z]]\to\KK[[X]]$, Lemma~\ref{hol:nl:lem:derivatives}
turns the $\nu$th summand of $P[f]$ into
\[
 c_\nu\,t^{d_\nu\gamma-w_\nu\delta}\,X^{k_\nu}\prod_j\bigl(g^{(j)}\bigr)^{m_{\nu j}}.
\]
Multiply the whole identity by $t^{-\lambda}$ with
$\lambda=\beta_P+d_P\gamma-w_P\delta$. The scalar in front of the $\nu$th
product then has valuation
\begin{equation}\label{hol:nl:eq:scalarvalue}
 v(c_\nu)-\beta_P+(d_\nu-d_P)\gamma+(w_P-w_\nu)\delta.
\end{equation}
For $\nu\in\mathcal C_P$ this is $v(c_\nu)-\beta_P\ge0$. For $d_\nu=d_P$ and
$w_\nu<w_P$ it is positive by \eqref{hol:nl:eq:cert2}. For $d_\nu<d_P$,
multiplying the bound on $\gamma$ in \eqref{hol:nl:eq:hbound} by the negative
integer $d_\nu-d_P$ reverses it and gives the lower bound
\[
 v(c_\nu)-\beta_P-(d_P-d_\nu)v(a_m)
 +\bigl((d_P-d_\nu)m+w_P-w_\nu\bigr)\delta>0
\]
by \eqref{hol:nl:eq:cert3}. So the normalized identity lies in $\Sfam_1^+$,
and only the corner summands survive the residue map. Applying $\rho$ to the
formal identity gives
\begin{equation}\label{hol:nl:eq:residualequation}
 0=\sum_{\nu\in\mathcal C_P}\bar c_\nu X^{k_\nu}
 \prod_j\bigl(\phi^{(j)}\bigr)^{m_{\nu j}}=\Omega_\phi(X).
\end{equation}
But $N>n_\ast\ge s$ and $\chi_P(N)\ne0$, so by
Lemma~\ref{hol:nl:lem:corner} the coefficient of $X^{d_PN+w_P}$ in
$\Omega_\phi$ is $\lc(\phi)^{d_P}\chi_P(N)\ne0$, a contradiction. The last
assertion follows because a strongly entire series is evaluable at every
monomial argument.
\end{proof}

\paragraph{An explicit, division-free exclusion bound.}
Let $\mathcal R$ be the finite set of right-hand sides of
\eqref{hol:nl:eq:cert1}--\eqref{hol:nl:eq:cert3}; an empty list contributes
nothing. For any $\epsilon\in\Gamma_{>0}$,
\begin{equation}\label{hol:nl:eq:deltazero}
 \delta_0=\epsilon+\max\bigl(\{0\}\cup\{\abs b:b\in\mathcal R\}\bigr)
\end{equation}
satisfies all of them, since their multipliers are positive integers. This is
an explicit certificate; it is not claimed to be optimal. The proof uses the
full supports to form the initial polynomial but only finitely many leading
valuations to choose $\delta_0$. Leading valuations alone do \emph{not} in
general decide strong evaluation at a boundary argument, and the certificate
has no converse.

For $k=\C$, Lemma~\ref{hol:lem:inward} (inward stability) turns the monomial
conclusion into the exterior form used in Theorem~\ref{hol:main:ode}: a
nonpolynomial formal solution is then not strongly evaluable at any
$x\ne0$ with $v(x)\le-\delta_0$. This combination of the two statements is
made by the merge; source~10 states only the monomial form.

\begin{corollary}[Cofinal scales suffice]\label{hol:nl:cor:cofinal}
If $\chi_P\ne0$, a formal solution of $P[f]=0$ strongly evaluable at
$t^{-\delta}$ for a cofinal set of $\delta\in\Gamma_{>0}$ is a polynomial.
\end{corollary}

\begin{proof}
A cofinal subset meets the final segment excluded by
Theorem~\ref{hol:nl:thm:certificate} for a nonpolynomial solution.
\end{proof}

\paragraph{Dependency audit.}
The only infinite-support step is Lemma~\ref{hol:nl:lem:algebra}, with the
existence and finiteness of a leading coefficient layer; the rest is finite
polynomial algebra. The proof makes none of the following replacements: a
higher-rank group by a real-valued norm; strong summation by convergence of
finite subsums; a nonsummable family by a cancelled expression; the residue
polynomial by its value at~$1$; a nonlinear coefficient recurrence by a
linear one. The three critical checkpoints are the minimum in
\eqref{hol:nl:eq:gaussdata}, the order reversal when \eqref{hol:nl:eq:hbound}
is multiplied by $d_\nu-d_P<0$, and the nonvanishing of $\chi_P(N)$.

\subsection{Why the corner never degenerates in order one; proof of
\texorpdfstring{Theorem~\ref{hol:nl:main:first}}{Theorem D}}
\label{hol:nl:sub:firstorder}

Write a first-order equation in combined-monomial form as
\begin{equation}\label{hol:nl:eq:firstP}
 P(z,Y,Y')=\sum_{k,i,j}c_{kij}z^kY^i(Y')^j.
\end{equation}
For a term, $d=i+j$ and $w=k-j$. At the corner, fixing $j$ forces
$i=d_P-j$ and $k=w_P+j$, so there is at most one corner term for each $j$, and
\eqref{hol:nl:eq:chiI} becomes
\begin{equation}\label{hol:nl:eq:firstchi}
 \chi_P(T)=\sum_{\substack{i+j=d_P\\ k-j=w_P}}\bar c_{kij}T^j,\qquad
 I_P(T)=\sum_{\substack{i+j=d_P\\ k-j=w_P}}c_{kij}T^j.
\end{equation}
At least one $\bar c_{kij}$ is nonzero by the choice of $\beta_P$, and its
power $T^j$ cannot cancel against another summand. Hence $\chi_P\ne0$.

\begin{proof}[Proof of Theorem~\ref{hol:nl:main:first}]
If $d_P=0$, then $P$ is a nonzero polynomial in $z$ alone and has no formal
solutions. Otherwise the preceding observation and
Theorem~\ref{hol:nl:thm:certificate} give polynomiality and the exclusion
statement, with $\delta_0$ as in \eqref{hol:nl:eq:deltazero}. A nonzero
algebraic relation between $f$ and $D_zf$ over $\KK(z)$ becomes a nonzero
polynomial of the form \eqref{hol:nl:eq:firstP} after denominators are
cleared, which gives the independence statement.
\end{proof}

The uniqueness of a corner term for fixed $j$ is a structural fact of order
one. It is exactly what fails when several derivative orders can contribute
the same falling-factorial polynomial (Section~\ref{hol:nl:sub:vanishing}).

\subsection{Finite candidate degrees and the entire-solution locus}
\label{hol:nl:sub:degrees}

Once polynomiality is known, ordinary polynomial degree gives a sharper
finite candidate set that uses the \emph{full} coefficients. Assume $d_P\ge1$
and define the nonnegative rational number
\begin{equation}\label{hol:nl:eq:kappa}
 \kappa_P=\max\Bigl(\{0\}\cup\Bigl\{\frac{w_\nu-w_P}{d_P-d_\nu}:d_\nu<d_P\Bigr\}\Bigr)\in\Q.
\end{equation}
This divides ordinary integers in $\Q$, not elements of $\Gamma$.

\begin{theorem}[Finite candidate degrees]\label{hol:nl:thm:degrees}
For the first-order polynomial \eqref{hol:nl:eq:firstP} with $d_P\ge1$, every
nonconstant strongly entire solution has degree in the finite set
\begin{equation}\label{hol:nl:eq:candidates}
 \mathcal D_P=\{n\in\Z_{>0}:n\le\kappa_P\}\ \cup\ \{n\in\Z_{>0}:I_P(n)=0\}.
\end{equation}
The constant solutions are exactly the $c\in\KK$ with $P(z,c,0)=0$ as a
polynomial in $z$. The set \eqref{hol:nl:eq:candidates} is necessary; it does
not promise that each candidate occurs.
\end{theorem}

\begin{proof}
Polynomiality is Theorem~\ref{hol:nl:main:first}. Let $f=az^n+\cdots$ with
$n\ge1$ and $a\ne0$. A term of \eqref{hol:nl:eq:firstP} has degree
$k+in+j(n-1)=dn+w$ and leading coefficient $c_{kij}a^dn^j$. If $n>\kappa_P$,
every term with $d<d_P$ has degree below $d_Pn+w_P$, and so do the terms with
$d=d_P$, $w<w_P$. So the coefficient of $z^{d_Pn+w_P}$ in $P[f]$ is
$a^{d_P}I_P(n)$, which must vanish. $I_P$ is nonzero by
\eqref{hol:nl:eq:firstchi}, so its set of ordinary integer roots is finite.
The assertion on constants is immediate by substitution.
\end{proof}

\begin{remark}[What ``finite'' does and does not compute]
The set of ordinary integer roots of $I_P$ is well defined and finite for
arbitrary Hahn coefficients, but its existence is not a uniform algorithm for
arbitrary coefficient names: extracting a leading exponent, deciding
coefficient equality, or testing an arbitrary complex coefficient for exact
integrality can need information that no numerical approximation supplies.
For rational coefficients, or other exact representations with the needed
operations, \eqref{hol:nl:eq:candidates} is an actual finite algebraic
search. The shipped program uses only exact rational data.
\end{remark}

Replacing $I_P(n)=0$ by $\chi_P(n)=0$ gives a coarser necessary set, by
\eqref{hol:nl:eq:rootcontainment}. Keeping $I_P$ can be genuinely sharper.

\begin{example}[A resonance visible only in the first layer]
\label{hol:nl:ex:resonance}
Let $c=t^\eta$ with $\eta>0$, $N\in\Z_{>0}$, and $P=zY'-(N+c)Y$. Then
$\chi_P(T)=T-N$ but $I_P(T)=T-N-c$ has no ordinary nonnegative integer root.
The only formal solution in $\KK[[z]]$ is zero, since coefficient comparison
gives $(n-N-c)a_n=0$ for every $n$. A residue resonance alone is not a
polynomial solution.
\end{example}

\begin{corollary}[An algebraic entire-solution locus]\label{hol:nl:cor:locus}
For a first-order equation $P$ with $d_P\ge1$ put
$n_P=\max(\{0\}\cup\mathcal D_P)$. There is an explicitly defined affine
algebraic set $\mathcal V_P\subseteq\mathbb A_{\KK}^{n_P+1}$ whose
$\KK$-points are in bijection with all strongly entire formal solutions of
$P[f]=0$. The same equations parametrize all strongly entire solutions over
every Hahn-field extension $\KK\hookrightarrow\mathbb L=k'((t^{\Gamma'}))$
induced by an embedding of coefficient fields $k\hookrightarrow k'$ and an
embedding of ordered groups $\Gamma\hookrightarrow\Gamma'$.
\end{corollary}

\begin{proof}
With indeterminates $u_0,\ldots,u_{n_P}$ put $F_u(z)=\sum_{n\le n_P}u_nz^n$ and
expand $P(z,F_u,F_u')=\sum_\ell Q_\ell(u)z^\ell$ with
$Q_\ell\in\KK[u_0,\ldots,u_{n_P}]$. Let $\mathcal V_P$ be the common zero
locus of these finitely many $Q_\ell$. Its points are exactly the polynomial
solutions of degree at most $n_P$, each of which is strongly entire, and
Theorem~\ref{hol:nl:thm:degrees} shows there are no others. Under the stated
extension the monomial degrees and $\kappa_P$ are unchanged, and for an
ordinary integer $n$ the element $I_P(n)$ vanishes in $\mathbb L$ if and only
if it vanishes in $\KK$. Theorem~\ref{hol:nl:main:first} applies over
$\mathbb L$, whose coefficient field again has characteristic zero and whose
group is again nonzero, so the same degree bound and equations remain
exhaustive there.
\end{proof}

This is a finite-dimensional algebraic parametrization, not a claim that the
number of solutions is finite, and not a moduli statement over coefficient
rings with nilpotents. The extension class is exactly the one stated.

\subsection{Worked first-order equations}\label{hol:nl:sub:examples}

All identities use the formal derivative $D_z$. For $k=\C$ they are statements
about surcomplex Hahn workspaces; for $k=\R$ the real versions follow with the
non-real constant roots omitted. Consider
\begin{equation}\label{hol:nl:eq:riccati}
 D_zf=f^2+p(z),\qquad p\in\KK[z].
\end{equation}
The corner is the term $-Y^2$, so $d_P=2$, $w_P=0$ and $I_P=-1$. If $p$ has
degree $e\ge0$, then $\kappa_P=\max(0,e/2)$, so a nonconstant entire solution
has degree at most $\lfloor e/2\rfloor$. More precisely, a solution of degree
$n\ge1$ requires $\deg p=2n$ and the leading coefficient of $p$ to be minus
the square of that of $f$, since the top term of $f^2$ cannot be cancelled by
$D_zf$. These conditions are necessary only; the remaining coefficient
equations must still hold.

\begin{proposition}[Two exact Riccati classifications]\label{hol:nl:prop:riccati}
Over $\KK=\C((t^\Gamma))$, the strongly entire solutions of $D_zf=f^2+1$ are
exactly $f=i$ and $f=-i$, and the only strongly entire solution of
$D_zf=f^2+1-z^2$ is $f=z$.
\end{proposition}

\begin{proof}
In the first equation the degree bound excludes nonconstant solutions, and a
constant must satisfy $c^2+1=0$; in the field $\KK$, which contains $i$, the
roots are $\pm i$. In the second the degree bound is one and constants are
impossible. For $f=az+b$ coefficient comparison gives $a^2=1$, $2ab=0$ and
$a=b^2+1$, so $b=0$ and $a=1$, and $f=z$ is verified directly.
\end{proof}

The first equation is also \texttt{diff:cor:riccati} of the collection's
differential-equations report~\cite{repodifferential}, with the same two
solutions, but there $\partial u=1+u^2$ is solved in $\No[i]$ for a derivation
on surreal \emph{scalars}. The two statements concern different derivations
and different objects, and neither implies the other.

\begin{example}[The degree cutoff is attained]\label{hol:nl:ex:cutoff}
For every integer $m\ge1$ the equation $D_zf=f^2+mz^{m-1}-z^{2m}$ has the
strongly entire solution $f=z^m$. Here $I_P=-1$ and $\kappa_P=m$ exactly, so
the low-degree part of \eqref{hol:nl:eq:candidates} cannot in general be
replaced by $n<\kappa_P$.
\end{example}

\begin{example}[Arbitrarily large resonant degrees]\label{hol:nl:ex:resonant}
For $N\ge1$ the equation
\begin{equation}\label{hol:nl:eq:resonant}
 zf\,D_zf-Nf^2=0
\end{equation}
has $d_P=2$, $w_P=0$ and $I_P(T)=T-N$. Since $\KK[[z]]$ is an integral
domain, a nonzero solution satisfies $zD_zf=Nf$, and coefficient comparison
gives exactly $f=cz^N$ with $c\in\KK$ (including $c=0$). The exceptional
integer root is an actual degree, and no degree bound can depend only on the
positions of the monomials of $P$: the coefficient $N$ matters.
\end{example}

\subsection{The exclusion bound cannot be uniform}\label{hol:nl:sub:nouniform}

Theorem~\ref{hol:main:ode} has an exclusion bound depending only on the
operator. For nonlinear equations this is false, already for one fixed
equation of first order.

\begin{theorem}[A fixed Riccati equation defeats every uniform bound]
\label{hol:nl:thm:nouniform}
For $\lambda\in\Gamma$ let $c=t^\lambda$ and
\[
 f_\lambda(z)=\frac{c}{1-cz}=\sum_{n\ge0}c^{n+1}z^n\in\KK[[z]].
\]
Then $D_zf_\lambda=f_\lambda^2$, $f_\lambda$ is nonpolynomial, and
\begin{equation}\label{hol:nl:eq:riccatidomain}
 \Dom(f_\lambda)=\{0\}\cup\{x\in\KK^\times:v(x)+\lambda>0\}.
\end{equation}
Consequently, for every $\delta\in\Gamma$ some nonpolynomial solution of the
fixed equation $D_zf=f^2$ is strongly evaluable at $t^{-\delta}$, and no
exclusion bound depending on the equation alone excludes all its
nonpolynomial solutions.
\end{theorem}

\begin{proof}
The derivative identity is a coefficient comparison. For $x\ne0$ put $u=cx$,
so the $n$th evaluation term is $cu^n$. If $v(u)>0$, write
$u=b\,t^{v(u)}(1+w)$ with $b\in k^\times$ and $w$ zero or of positive
valuation; then $\supp(u^n)\subseteq nv(u)+M(\supp w)$, the leading exponents
$nv(u)$ increase strictly, and Lemma~\ref{hol:lem:certificate}\,(b), resting
on the positive-support monoid of Lemma~\ref{hol:lem:support}\,(2), gives
strong summability; multiplication by $c$ preserves it. (Source~10 cites the
Hahn--Neumann positive-support lemma here instead of proving it; no
Archimedean convergence of $n\,v(u)$ is used.) If $v(u)=0$, every $cu^n$ has
the nonzero coefficient $\lc(u)^n$ at exponent $\lambda$, so local finiteness
fails. If $v(u)<0$, the leading exponents $\lambda+nv(u)$ strictly decrease.
At $x=0$ only one member is nonzero. This proves
\eqref{hol:nl:eq:riccatidomain}. Given $\delta$, choose $\lambda>\delta$;
then $v(t^{-\delta})+\lambda=\lambda-\delta>0$.
\end{proof}

There is no conflict with Theorem~\ref{hol:nl:thm:certificate}, whose
finite data, and hence whose bound, may depend on the solution; for
$f_\lambda$ the exact boundary moves with $\lambda$. At the boundary
$v(x)=-\lambda$ geometric-series algebra cannot legalize infinitely many
contributions at one exponent. Nor is there a conflict with
Theorem~\ref{hol:main:ode}: each $f_\lambda$ also satisfies the linear
equation $(1-t^\lambda z)D_zf-t^\lambda f=0$, whose operator depends on
$\lambda$, exactly as the operator $D_z-t^\eta$ of the formal exponential of
Proposition~\ref{hol:prop:expdomain} depends on~$\eta$. (This comparison is
the merge's; it is a direct computation.)

\section{Higher order and positive-weight Euler equations}
\label{hol:nl:sec:higher}

\subsection{The nondegenerate higher-order class}

\begin{theorem}[Higher-order corner criterion and its boundary]
\label{hol:nl:thm:corner}
Let $P$ be a nonzero differential polynomial \eqref{hol:nl:eq:P} of any finite
order with $d_P\ge1$ and $\chi_P\ne0$. Then every strongly entire formal
solution of $P[f]=0$ is a polynomial, and every nonpolynomial formal solution
fails strong evaluation at $t^{-\delta}$ for all $\delta$ in a final segment
of $\Gamma_{>0}$, with the finite certificate of
Theorem~\ref{hol:nl:thm:certificate}. The condition $\chi_P\ne0$ holds
automatically in order one, but not in order two: a fixed second-order
equation can have polynomial solutions of every degree
(Proposition~\ref{hol:nl:prop:degenerate}).
\end{theorem}

\begin{proof}
The first two assertions are Theorem~\ref{hol:nl:thm:certificate}; the third
is Section~\ref{hol:nl:sub:firstorder} and
Proposition~\ref{hol:nl:prop:degenerate}.
\end{proof}

\begin{proposition}[Higher-order degree candidates]
\label{hol:nl:prop:higherdegrees}
Suppose $d_P\ge1$ and $\chi_P\ne0$ for an equation of order at most $s$. Every
nonzero strongly entire solution has degree in the finite set
\[
 \{0,\ldots,s-1\}\ \cup\ \{n\in\N:n\ge s,\ n\le\kappa_P\}\ \cup\
 \{n\in\N:n\ge s,\ I_P(n)=0\},
\]
with $\kappa_P$ as in \eqref{hol:nl:eq:kappa}. When $s=0$ the first set is
empty. Constant solutions can always be checked by substitution.
\end{proposition}

\begin{proof}
Polynomiality is Theorem~\ref{hol:nl:thm:certificate}. For degree $n\ge s$
every derivative leading-term formula of Lemma~\ref{hol:nl:lem:corner} applies,
and if also $n>\kappa_P$ the top coefficient of $P[f]$ is
$\lc(f)^{d_P}I_P(n)$, which must vanish; $I_P\ne0$ because $\chi_P\ne0$.
Degrees below $s$, where derivative leading terms can vanish, are listed
separately.
\end{proof}

\begin{proposition}[The first two Painlev\'e polynomial equations]
\label{hol:nl:prop:painleve}
In characteristic zero, the equation $D_z^2f=6f^2+z$ has no strongly entire
formal solution, and the equation $D_z^2f=2f^3+zf+\alpha$ with $\alpha\in\KK$
has a strongly entire solution exactly when $\alpha=0$, the only one then
being $f=0$.
\end{proposition}

\begin{proof}
The first equation has $d_P=2$, $w_P=0$, $\chi_P=-6$; the second has $d_P=3$,
$w_P=0$, $\chi_P=-2$. Theorem~\ref{hol:nl:thm:certificate} makes every entire
solution a polynomial. In the first equation a polynomial of degree $n\ge1$
gives $f^2$ of degree $2n$, larger than $\deg D_z^2f$ (when nonzero) and than
$\deg z=1$, so no cancellation is possible, and a constant cannot cancel $z$.
In the second, for $n\ge1$ one has $3n>n+1$ and $3n>\deg D_z^2f$, so $2f^3$
cannot cancel; a constant $c$ must satisfy $2c^3+zc+\alpha=0$ identically,
which forces $c=0$ and $\alpha=0$, and substitution proves sufficiency.
\end{proof}

The names only identify the displayed equations; no classification of their
classical meromorphic solutions is imported. The examples illustrate a class
beyond first order: a unique term of greatest total differential degree has a
nonzero corner whenever its falling-factorial product is nonzero. In
particular an equation $D_z^sf=Q(z,f)$ with $Q$ of degree at least two in $f$
has a nonzero corner and hence only polynomial strongly entire solutions; the
coefficient of the top power of $f$ may be a nonconstant polynomial in $z$,
whose highest $z$-term then selects the corner.

\subsection{A vanishing corner}\label{hol:nl:sub:vanishing}

Consider
\begin{equation}\label{hol:nl:eq:degenerate}
 z^2f\,D_z^2f+zf\,D_zf-z^2(D_zf)^2=0.
\end{equation}
All three terms have $d=2$ and $w=0$, and the residue corner polynomial is
\begin{equation}\label{hol:nl:eq:degeneratechi}
 T(T-1)+T-T^2=0,
\end{equation}
an identity in $k[T]$, not an exceptional integer resonance.

\begin{proposition}[Complete formal classification of the example]
\label{hol:nl:prop:degenerate}
The formal solutions of \eqref{hol:nl:eq:degenerate} in $\KK[[z]]$ are exactly
$0$ and $cz^n$ with $c\in\KK^\times$ and $n\in\N$. All are strongly entire,
and their degrees are unbounded for this single equation.
\end{proposition}

\begin{proof}
With $\vartheta=zD_z$, equation \eqref{hol:nl:eq:degenerate} reads
$f\,\vartheta^2f-(\vartheta f)^2=0$. For $f\ne0$ work in the Laurent series
field $\KK((z))$ and divide by $f^2$: $\vartheta(\vartheta f/f)=0$. The
constants of $\vartheta$ in $\KK((z))$ are exactly $\KK$, since
$\vartheta\sum b_nz^n=\sum nb_nz^n$ and characteristic zero kills no nonzero
integer. Hence $\vartheta f/f=c_0\in\KK$; if the first nonzero term of $f$
has exponent $n$, comparing lowest terms gives $c_0=n$. Then $\vartheta f=nf$
kills every coefficient except the $n$th. Substitution verifies the monomials
and zero.
\end{proof}

The example has two messages. A universal higher-order degree bound cannot
be obtained by declaring $\chi_P$ nonzero: it may vanish identically, and a
fixed equation can then have unbounded degrees. And the example is \emph{not}
a nonpolynomial entire solution, so it does not refute polynomial rigidity in
higher order; it locates a failure of this proof and of a stronger uniform
conclusion.

The corner is also not an invariant of the solution set. Multiplying an
equation by another differential polynomial, or passing to a differential
consequence, can change the corner data. The criterion $\chi_P\ne0$ is a
sufficient condition on the displayed combined polynomial: when $\chi_P=0$
another equation satisfied by the same function may have a nonzero corner, or
a logarithmic-derivative reduction as in
Proposition~\ref{hol:nl:prop:degenerate} may be available. A zero corner
cannot be repaired by retaining only monomial degrees, since the residue
coefficients and their exact cancellation are essential. The methods here do
not decide whether a higher-order equation with vanishing residue corner has
nonpolynomial strongly entire solutions; that is
Question~\ref{hol:q:nonlinear}.

\subsection{Positive-weight Euler equations in several variables}
\label{hol:nl:sub:euler}

Fix positive integers $q=(q_1,\ldots,q_m)$, write
$\wt(\alpha)=q\cdot\alpha$ for $\alpha\in\N^m$, and let
\[
 \vartheta_q=\sum_{\nu=1}^mq_\nu x_\nu\frac{\partial}{\partial x_\nu},
 \qquad \vartheta_qx^\alpha=\wt(\alpha)x^\alpha.
\]
Only finitely many multi-indices have weight at most a given integer; this,
with the eigenvalue identity, is the essential feature of positive weights.
Write a nonzero equation in combined-monomial form
\begin{equation}\label{hol:nl:eq:eulerP}
 P(x,Y,Z)=\sum_{\nu\in J}c_\nu x^{\alpha_\nu}Y^{i_\nu}Z^{j_\nu},
\end{equation}
to be read as $P(x,f,\vartheta_qf)=0$, and set
$d_\nu=i_\nu+j_\nu$, $w_\nu=\wt(\alpha_\nu)$, with $d_P$, $w_P$,
$\mathcal C_P$, $\beta_P$ and $\bar c_\nu$ as in
Section~\ref{hol:nl:sub:corner}. The corner polynomials now keep the spatial
variables:
\begin{equation}\label{hol:nl:eq:eulerchi}
 \chi(T,X)=\sum_{\nu\in\mathcal C_P}\bar c_\nu X^{\alpha_\nu}T^{j_\nu},\qquad
 I(T,X)=\sum_{\nu\in\mathcal C_P}c_\nu X^{\alpha_\nu}T^{j_\nu}.
\end{equation}
Both are nonzero: for fixed $(\alpha_\nu,j_\nu)$ the condition $d_\nu=d_P$
forces $i_\nu=d_P-j_\nu$, so no two combined corner terms give the same
monomial in $(X,T)$. Only finitely many integers $n$ make $\chi(n,X)$ the zero
polynomial, since they must all be roots of one nonzero $X$-coefficient of
$\chi$, and likewise for $I$.

\begin{theorem}[Positive-weight Euler rigidity]\label{hol:nl:thm:euler}
Let $f\in\KK[[x_1,\ldots,x_m]]$ satisfy $P(x,f,\vartheta_qf)=0$ for a nonzero
$P\in\KK[x_1,\ldots,x_m,Y,Z]$. If $f$ is nonpolynomial, it fails joint strong
evaluation at $(t^{-q_1\delta},\ldots,t^{-q_m\delta})$ for every sufficiently
large $\delta\in\Gamma_{>0}$. Consequently, if $f$ is strongly entire, or
merely jointly strongly evaluable at these points for a cofinal set of
$\delta\in\Gamma_{>0}$, then $f$ is a polynomial. Every fixed equation has a
finite set of possible positive weighted degrees
(Corollary~\ref{hol:nl:cor:eulerdegrees}).
\end{theorem}

\begin{proof}
If $d_P=0$, then $P$ is a nonzero polynomial in $x$ alone and there is no
solution; assume $d_P\ge1$. Choose $n_\ast\ge0$ with $\chi(n,X)\ne0$ for every integer $n>n_\ast$. Since
$f=\sum_\alpha a_\alpha x^\alpha$ is nonpolynomial and bounded-weight sets are
finite, choose $a_\mu\ne0$ of weight $n_\mu=\wt(\mu)>n_\ast$ so large that
$(d_P-d_\nu)n_\mu+w_P-w_\nu>0$ for $d_\nu<d_P$. Take $\delta>0$ satisfying
the finitely many inequalities
\begin{align*}
 (n_\mu-\wt(\alpha))\delta&>v(a_\mu)-v(a_\alpha)
 &&(\wt(\alpha)\le n_\ast,\ a_\alpha\ne0),\\
 (w_P-w_\nu)\delta&>\beta_P-v(c_\nu)
 &&(d_\nu=d_P,\ w_\nu<w_P),\\
 \bigl((d_P-d_\nu)n_\mu+w_P-w_\nu\bigr)\delta
 &>\beta_P-v(c_\nu)+(d_P-d_\nu)v(a_\mu)
 &&(d_\nu<d_P);
\end{align*}
as before they can be met in every nonzero ordered group and persist on a
final segment. Suppose the joint family is strongly summable at such a scale,
and put
\[
 F(X)=f(t^{-q_1\delta}X_1,\ldots,t^{-q_m\delta}X_m)\in\Sfam_m,\quad
 \gamma=\gamma(F),\quad g=t^{-\gamma}F,\quad \phi=\rho(g).
\]
Then $\gamma\le v(a_\mu)-n_\mu\delta$, and the first inequalities show that
$\phi$ has no monomial of weight at most $n_\ast$; its greatest weight $N$
exceeds $n_\ast$. Let $\phi_N\ne0$ be its weight-$N$ homogeneous part. The
Euler operator commutes with the scaling, so
$(\vartheta_qf)(t^{-q_1\delta}X_1,\ldots,t^{-q_m\delta}X_m)=t^\gamma\vartheta_qg$.
Normalizing the substituted equation by $t^{-\beta_P-d_P\gamma+w_P\delta}$,
the scalar computation \eqref{hol:nl:eq:scalarvalue} shows that only corner
terms survive the residue, whence
\begin{equation}\label{hol:nl:eq:eulerresidue}
 0=\sum_{\nu\in\mathcal C_P}\bar c_\nu X^{\alpha_\nu}\phi^{i_\nu}
 (\vartheta_q\phi)^{j_\nu}.
\end{equation}
The largest possible weight on the right is $d_PN+w_P$. Since the weight-$N$
part of $\vartheta_q\phi$ is $N\phi_N$, the homogeneous part of that weight is
\begin{equation}\label{hol:nl:eq:euleridentity}
 \phi_N(X)^{d_P}\,\chi(N,X),
\end{equation}
a product of two nonzero polynomials, hence nonzero. This contradicts
\eqref{hol:nl:eq:eulerresidue}.
\end{proof}

The proof uses \emph{joint} strong evaluation, not a scalar expression
obtained by first setting $x_1=\cdots=x_m$ and cancelling: a factor
$x_1-x_2$ can make a diagonal scalar expression zero while the ungrouped
monomial family is not summable. Keeping $X_1,\ldots,X_m$ in
\eqref{hol:nl:eq:euleridentity} avoids that error.

\begin{corollary}[Weighted-degree candidates]\label{hol:nl:cor:eulerdegrees}
With $\kappa_P$ defined by \eqref{hol:nl:eq:kappa} from the Euler data
$d_\nu,w_\nu$, every nonconstant strongly entire solution of
$P(x,f,\vartheta_qf)=0$ has greatest weight in the finite set
\[
 \{n\in\Z_{>0}:n\le\kappa_P\}\ \cup\ \{n\in\Z_{>0}:I(n,X)\equiv0\}.
\]
For a fixed equation all strongly entire solutions are therefore described by
finitely many polynomial equations in the coefficients of a bounded-weight
polynomial ansatz.
\end{corollary}

\begin{proof}
Polynomiality is Theorem~\ref{hol:nl:thm:euler}. For a polynomial of greatest
weight $N>\kappa_P$, the terms of lower differential degree have smaller total
weight, and the top-weight part of the equation is $f_N^{d_P}I(N,X)$, which
must vanish; so $I(N,X)=0$. The finite-root argument after
\eqref{hol:nl:eq:eulerchi} applies to $I$, and a bounded-weight ansatz has
finitely many coefficients.
\end{proof}

\begin{example}[An exactly described solution space]\label{hol:nl:ex:eulerspace}
For $N\ge0$ the solutions of $\vartheta_qf=Nf$ in $\KK[[x]]$ are exactly the
weighted homogeneous polynomials of weight $N$, since coefficient comparison
gives $(\wt(\alpha)-N)a_\alpha=0$. For $q=(1,2)$ this space has dimension
$\lfloor N/2\rfloor+1$ over $\KK$, with basis $x_1^{N-2j}x_2^j$,
$0\le j\le\lfloor N/2\rfloor$. The nonlinear equation
$f(\vartheta_qf-Nf)=0$ has the same solutions, because $\KK[[x]]$ is an
integral domain.
\end{example}

\paragraph{Why positive weights are needed.}
Let $\KK=\C((t^\Q))$ and $G(u)=\sum_{n\ge0}t^{n^2}u^n$, which is strongly
entire (Corollary~\ref{hol:nl:cor:theta}). The function $f(x,y)=G(y)$ is
nonpolynomial and strongly entire but satisfies $x\partial_xf=0$, so a zero
weight invalidates the theorem. For the mixed weights $(1,-1)$, the
nonpolynomial strongly entire function $f(x,y)=G(xy)$ satisfies
$(x\partial_x-y\partial_y)f=0$; at a fixed point its nonzero joint evaluation
terms are exactly $t^{n^2}(xy)^n$, which are strongly summable. Positivity is
therefore not a cosmetic substitute for nonzero weights, and a general
first-order partial differential equation need not have the rigidity of a
positive Euler operator.

\paragraph{Relation to Theorem~\ref{hol:main:multi}.}
Theorem~\ref{hol:main:multi} assumes that \emph{all} mixed partial
derivatives span a finite-dimensional space, and Section~\ref{hol:sec:multi}
notes that a single differential equation in several variables cannot force
polynomiality in general, since it may ignore variables; the zero-weight
example above is an instance. Theorem~\ref{hol:nl:thm:euler} instead assumes
one algebraic relation between $f$ and $\vartheta_qf$ with all weights
positive, and no finite-dimensional span. Neither theorem implies the other.

\section{Nonlinear consequences, hypotheses and predecessors}
\label{hol:nl:sec:consequences}

\subsection{Two-jet independence of explicit series}\label{hol:nl:sub:jets}

With the surreal embedding of Section~\ref{hol:sub:embedding}, for instance
$\C((\omega^{-\Q}))$ is a set-sized surcomplex workspace to which
Theorem~\ref{hol:nl:main:first} applies, and its coefficients and admissible
arguments are actual elements of $\No[i]$. The independence conclusions,
however, concern \emph{formal functions}: each value $f(x)$ lies in $\KK$ and
is of course not transcendental over $\KK$.

\begin{corollary}[Two independent jets of the quadratic-exponent series]
\label{hol:nl:cor:theta}
Over $\KK=\C((t^\Q))$ let $G(z)=\sum_{n\ge0}t^{n^2}z^n$. Then $G$ is
nonpolynomial and strongly entire, and $G$ and $D_zG$ are algebraically
independent over $\KK(z)$.
\end{corollary}

\begin{proof}
For $x\ne0$ with $v(x)=\sigma\in\Q$ the $n$th evaluated member has valuation
$n^2+n\sigma$, which eventually exceeds every rational bound, so
Lemma~\ref{hol:lem:certificate}\,(a) gives strong summability even when $x$
has an arbitrary infinite Hahn tail; at $x=0$ the family is finite. All
coefficients are nonzero. Theorem~\ref{hol:nl:main:first} gives the
independence.
\end{proof}

The series is not new: it is the rank-one example of the
rank-one-Berkovich report~\cite{reporankone}, Example
\texttt{found:ex:internal} of the foundations report~\cite{repofoundations},
and a rescaled partial theta series: $G(z)=\Theta_{t^2}(tz)$ in the notation
of \eqref{hol:eq:theta}. The contribution of source~10 is the nonlinear two-jet
consequence only.

\begin{corollary}[An example without an order unit]
\label{hol:nl:cor:infiniterank}
In $K_{\Gamma_\infty}$, with $\Gamma_\infty$ as in \eqref{hol:eq:gammainfty},
the strongly entire series $F=\sum_{n\ge0}t^{e_n}z^n=\sum_{n\ge0}\omega^{-\omega^n}z^n$
of \eqref{hol:eq:explicitF} satisfies: $F$ and $D_zF$ are algebraically
independent over $K_{\Gamma_\infty}(z)$.
\end{corollary}

\begin{proof}
Strong entireness is Theorem~\ref{hol:thm:explicit}; source~10 proves it
again by the cofinal-valuation route recorded in the second paragraph of that
proof. Apply Theorem~\ref{hol:nl:main:first}.
\end{proof}

$F$ and its linear rigidity are in Theorem~\ref{hol:thm:explicit}; the new
statement is the two-jet strengthening. It answers the first-order case of
the last sentence of Question~\ref{hol:q:nonlinear} with no coefficient
restriction, and says nothing about relations of order two or more.

\paragraph{Extension of scalars versus preservation of entireness.}
The degree classification is stable under the Hahn-field extensions of
Corollary~\ref{hol:nl:cor:locus}, but entireness of a nonpolynomial series
need not be. Enlarge $\Q$ to $\Q+\Q\omega$ with its surreal order and view $G$
over the larger field. At $z=t^{-\omega}$ the member exponents are
$n^2-n\omega$, with successive differences $(2n+1)-\omega<0$, so the support
union contains an infinite descending sequence and the evaluation is not
strong. The same computation appears in the rank-one-Berkovich
report~\cite{reporankone}, and for $F$ in
Section~\ref{hol:sub:notallsurcomplex}. No theorem here claims that $G$ or $F$ is
entire on the proper class $\No[i]$; hypotheses and derivative are relative to
a specified set-sized workspace.

\subsection{Which hypotheses are used}\label{hol:nl:sub:hypotheses}

\begin{example}[Positive characteristic destroys the conclusion]
\label{hol:nl:ex:charp}
Let $p$ be prime and $\KK=\mathbb F_p((t^\Q))$. The nonpolynomial series
$U(z)=\sum_{n\ge0}t^{n^2}z^{pn}$ is strongly entire: at a point of valuation
$\sigma$ its member valuations are $n^2+pn\sigma$, and
Lemma~\ref{hol:lem:certificate}\,(a) applies, its proof being independent of
the coefficient field. Yet $D_zU=0$.
\end{example}

So the conclusion fails in positive characteristic even for one linear
equation. The characteristic-zero requirements of the proof are substantive:
ordinary integer eigenvalues stay distinct, nonzero falling factorials stay
nonzero, and a nonzero polynomial cannot vanish at all ordinary integers.

The hypothesis $\Gamma\ne0$ is needed for the scale argument; the degenerate
case $\Gamma=0$ is treated in Section~\ref{hol:sec:fields}. Under ordinary
complex analytic summation, by contrast, $\exp z$ is a nonpolynomial entire
function with $f'=f$; these are different notions of infinite sum, and the
classical exponential does not contradict the Hahn theorem. The coefficient
field may be $\R$, $\C$ or any field of characteristic zero; neither algebraic
nor real closedness is used, no completeness of a real-valued norm is used,
and $\Gamma$ needs neither an order unit nor a countable cofinal subset.
Strong summability supplies exactly the finite initial layer the argument
needs.

\subsection{An exact finite workflow}\label{hol:nl:sub:workflow}

For first-order input $P$ the reduction is as follows. Combine like monomials
and remove zero coefficients. If no $Y$ or $Y'$ remains, decide the polynomial
identity directly. Otherwise compute $d_P$, $w_P$, $I_P$ and $\kappa_P$ from
the finite term list, find the positive integer roots of $I_P$ in the chosen
exact representation of the coefficients, take the largest candidate degree,
substitute a generic polynomial of that degree, and solve the resulting
finite algebraic system. This procedure is complete for \emph{entire}
solutions whenever its finite coefficient operations and algebraic solving are
available. It is not a general algorithm for arbitrary surreal names, it does
not enumerate the local formal solutions, and it does not decide arbitrary
strong-support membership. A negative answer after an exhaustive exact solve
is a genuine nonexistence certificate for entire solutions; failure of a
numerical solver is not.

For a nonpolynomial candidate formal solution a second workflow builds the
certificate \eqref{hol:nl:eq:cert1}--\eqref{hol:nl:eq:cert3} from one
sufficiently high nonzero coefficient and finitely many low-index ones; a
claimed certificate must supply their exact leading valuations. The shipped
program checks finite algebraic identities and sample certificates; it does
not decide infinitary summability.

\subsection{Predecessors, repository comparison and priority}
\label{hol:nl:sub:predecessors}

\paragraph{Inherited ideas.}
Hahn fields, the positive-support summation lemma, Conway normal forms, Gauss
initial forms, Euler eigenvalues and indicial-polynomial reasoning are
established ideas and are not relabelled as new. The organizing method is
related in spirit to maximum-term and central-index arguments in the
Wiman--Valiron tradition, surveyed by Hayman~\cite{hayman}; here the finite
residue polynomial gives an exact identity rather than an analytic
approximation, and no real-valued growth estimate is imported. Source~10
consulted Hayman's survey through its publisher metadata and abstract only,
and uses no theorem from its full text.

Newton-polygon methods for nonlinear equations have a substantial history.
Cano and Fortuny Ayuso construct initial and indicial polynomials for
nonlinear $q$-difference equations and obtain support and growth results for
generalized-series solutions~\cite{cano}. Their operator and setting are not
those of a formal $z$-series with arbitrary-rank Hahn \emph{coefficients}
evaluated strongly at every workspace scale. The conceptual precedent is
important, and no invention of differential Newton polygons is claimed.

Non-Archimedean value distribution already restricts nonlinear functional
and differential equations. Hu and Luan work over a complete algebraically
closed field with a nontrivial non-Archimedean absolute value and prove
Clunie-, Malmquist- and Mokhon'ko-type results~\cite{huluan}; source~10
consulted the parsed text of that paper for its field hypotheses and
Theorems~1.1, 1.3 and~1.5, reports that a remote screenshot service failed so
that no visual inspection of the PDF is claimed, and reproduces no figure or
table from it. Hu and Yang's monograph~\cite{huyang} is a broader
predecessor; only its publisher metadata and description were checked, so no
claim is made that all rank-one special cases were previously unknown.
Giansiracusa and Mereta develop differential enhancements of non-Archimedean
norms and a tropicalization of differential equations~\cite{giansiracusa};
the finite corner here records only the data this obstruction needs, and no
equivalence to their framework and no new tropical fundamental theorem is
asserted. Mantova and Matusinski's survey~\cite{mantovamatusinski} was
consulted for normal forms.

\paragraph{Repository comparison at source~10's pin.}
Source~10 compared its results with the collection at
\texttt{048b72c}; the comparison, with the corrections needed now, is as
follows.
\begin{center}
\small
\begin{tabularx}{\textwidth}{@{}>{\raggedright\arraybackslash}p{0.24\textwidth}
>{\raggedright\arraybackslash}X>{\raggedright\arraybackslash}X@{}}
\toprule
Material & Existing result or scope & Difference in the nonlinear part\\
\midrule
This report, Theorem~\ref{hol:main:ode} & Strongly entire D-finite series are
polynomials; the nonlinear case was the open Question~15.1, now
Question~\ref{hol:q:nonlinear}. & All first-order algebraic equations, a
higher-order corner criterion, and solution-dependent finite certificates.\\
This report, Theorem~\ref{hol:main:multi} & Finite-dimensional span of all
mixed derivatives forces polynomiality. & One algebraic relation in $f$ and
$\vartheta_qf$ suffices for positive weights; no finite-dimensional span.\\
Tail-spans report~\cite{repotail} & Particular Galois-supported constructions
have all-order differential independence, for scalar or fixed-domain analytic
derivations over specified smaller base fields. & A two-jet theorem for every
nonpolynomial strongly entire function over its own full Hahn field.\\
Documentation catalogue \texttt{docs/README.md} & Distinguishes workspaces, strong
summability, coefficient rings and scalar derivations. & Those distinctions
are kept; no global proper-class analytic claim is added.\\
\bottomrule
\end{tabularx}
\end{center}
The tail-span examples are neither weakened nor rediscovered: that report
proves stronger all-order independence for particular constructions under
different coefficient and function-class hypotheses, while the present
universal statement assumes strong entireness and concludes only two-jet
independence. Source~10 cited the quadratic-exponent series as occurring in
this report and in the catalogue; within the collection it occurs as the
rank-one-Berkovich example~\cite{reporankone}, as
\texttt{found:ex:internal}~\cite{repofoundations}, and as a case of
\eqref{hol:eq:theta}. Since that pin the collection has gained this nonlinear
part itself; no other report now contains a nonlinear rigidity theorem for
$D_z$, and the Riccati statement \texttt{diff:cor:riccati} of the
differential-equations report concerns a scalar derivation (compare
Proposition~\ref{hol:nl:prop:riccati}).

Source~10's repository audit covered the root README, the catalogue, contents
listings, this report's definitions, theorems and nonlinear-question passages,
and the README files of the tail-span, analytic-geometry and product-birthday
reports (the last two for context only). It was a targeted comparison, not a
line-by-line search of every report, Lean file, archive, branch or the Git
history, and no absence claim rests on a truncated preview.

\paragraph{Priority and proof status.}
Source~10's literature search, made on 22 September 2026, did not locate the
exact all-rank strong-family first-order theorem, the finite exclusion
certificate, or the positive-weight Euler formulation. Some searches were
poorly targeted, and no exhaustive MathSciNet, zbMATH, full-book or
citation-network review was made, so priority is provisional. The safest
description is a proposed new theorem package and a proved first-order answer
to a repository question, which is not a named historical conjecture; it is
not a certified first discovery of nonlinear non-Archimedean rigidity. The
written proofs are the evidence. The finite checks of
Section~\ref{hol:sec:verification} establish no arbitrary-support theorem;
no Lean formalization of the nonlinear part exists, and no independent human
referee has checked it.

\section{The three finite check suites}\label{hol:sec:verification}

Each source manuscript shipped an independent exact-arithmetic suite. The
suites of sources~08 and~09 were run on copies and passed, and neither
supersedes the other: they cover overlapping but distinct finite facts of the
linear theory, and they are kept side by side in \texttt{code/} for that
reason. The suite of source~10 checks finite facts of the nonlinear part
only; it was also run on a copy for this merge and reproduced its recorded
counts exactly. All three use only the Python standard library, integers, and
\texttt{fractions.Fraction}, with no floating point, no network access, and no
third-party packages.

\paragraph{\texttt{code/08-finite-recurrences-verify.py} --- 3{,}705 checks,
0 failures.}
It emphasises the operator-to-recurrence conversions in all three formats and
the higher-rank ordered-group arithmetic.
\begin{center}
\begin{tabularx}{\textwidth}{@{}X r@{}}
\toprule
Finite calculation & Checks\\
\midrule
Differential operator to coefficient recurrence & 648\\
Dilation operator to coefficient recurrence & 336\\
Mixed operator to coefficient recurrence & 276\\
Unit residues, parity, and finite exceptional cancellation & 400\\
Partial theta identities and finite quadratic estimates & 552\\
Higher-rank lexicographic valuation examples & 1{,}292\\
Sparse Airy recurrence and finite escape edges & 201\\
\midrule
Total & 3{,}705\\
\bottomrule
\end{tabularx}
\end{center}
The operator tests compare the grouped recurrence formulas against direct
finite polynomial differentiation, dilation and multiplication; the theta
tests compare Laurent monomial exponents exactly, including the $q^{-1}$-to-$q$
conversion of \eqref{hol:eq:inverseq}; the unit tests retain the exceptional
index $n=7$ rather than silently checking only generic values; and the
ordered-group checks use lexicographic integer tuples, not a real-valued
surrogate norm. This script takes a \emph{required} \verb|--output| path and
exits nonzero without one, which is deliberate: it prevents a re-run from
overwriting the delivered record.

\paragraph{\texttt{code/09-entire-hahn-holonomic-verification.py} --- 2{,}043
checks, 0 failures.}
It emphasises the boundary behaviour of the exclusion bounds and the exact
theta domain.
\begin{center}
\begin{tabularx}{\textwidth}{@{}X r@{}}
\toprule
Finite calculation & Checks\\
\midrule
Theta homogeneous equation, positive and negative valuation & 92\\
Theta first-order coefficient recurrence & 45\\
Dilation and differential forward recurrence conversions & 72\\
Unit orbit: parity, exceptional cancellation, residue polynomial & 504\\
Rank-two noncofinal affine minimum & 200\\
Theta excluded upper-scale argument & 200\\
Theta included coarse boundary & 170\\
Theta cofinal-scale growth & 120\\
Finite-coordinate entire witness growth & 480\\
Sparse recurrence exact identity & 50\\
Escape boundary: constant leading exponent, strict exterior descent & 50\\
Formal exponential recurrence & 60\\
\midrule
Total & 2{,}043\\
\bottomrule
\end{tabularx}
\end{center}
Two of these rows exist only in this suite and check the two boundary
phenomena that the report insists are different: ``theta included coarse
boundary'' checks that a zero coarse valuation coexists with positive growth
in the original lower-scale coordinate, which is
Section~\ref{hol:sub:coarseboundary}, and ``escape boundary'' checks the
constant-leading-exponent case that
Lemma~\ref{hol:lem:nonincreasing} rules out, which is what makes the excluded
region in Corollary~\ref{hol:cor:periodic} closed. Its \verb|--output|
argument is optional. Seeds are recorded in the shipped records.

\paragraph{\texttt{code/10-nonlinear-rigidity-verify.py} --- 2{,}113 checks,
0 failures.}
It checks the finite algebra of
Sections~\ref{hol:nl:sec:setting}--\ref{hol:nl:sec:consequences}.
\begin{center}
\begin{tabularx}{\textwidth}{@{}X r@{}}
\toprule
Finite calculation & Checks\\
\midrule
Finite Hahn initial layers and exclusion inequalities & 450\\
Automatic first-order nondegeneracy of the corner & 240\\
Univariate falling-factorial corner identities, orders $0$ to $3$ & 480\\
Lexicographic higher-rank ordered-group exclusion bounds & 500\\
Weighted Euler homogeneous-part identities & 360\\
Worked examples: Riccati, degree cutoff, resonance, Painlev\'e, vanishing
corner, Euler solution space & 83\\
\midrule
Total & 2{,}113\\
\bottomrule
\end{tabularx}
\end{center}
The group names in the record \path{data/10-nonlinear-rigidity-verification.json}
are \path{finite_hahn_certificate}, \path{first_order_nondegeneracy},
\path{higher_order_corner}, \path{ordered_group_radius},
\path{weighted_euler} and \path{worked_examples}, in the order of the
table; ``radius'' there is source~10's word for an exclusion bound. The
default seed is $20260922$. The \verb|--output| argument is required, and the
program refuses to overwrite an existing file. The vanishing-corner example
is tested for the monomial degrees $0,\ldots,20$ only. The counts are program
checks, not independently proved theorems.

\paragraph{What the checks do not do.}
No suite proves any infinite theorem. No finite prefix can prove strong
summability of an infinite family, cofinality of $(e_n)$, noncofinality of
$(n\epsilon)$, the existence of an infinite escape path, the eventual orbit
statements, the generic-line theorem, or the absence of every possible
annihilating equation. No infinite Hahn sum is evaluated, or numerically
approximated, by any of the programs. In particular a finite truncation of the
explicit series \eqref{hol:eq:explicitF} is a polynomial and is D-finite, so no
finite computation can witness its nonholonomicity; for the same reason no
finite computation can witness the two-jet independence of
Corollaries~\ref{hol:nl:cor:theta} and~\ref{hol:nl:cor:infiniterank}.
Source~10's suite in particular does not test arbitrary infinite supports,
cofinality, all Hahn elements, the nonexistence of all annihilating
equations, or historical novelty. What the suites can detect is sign,
indexing, conversion and finite valuation-comparison mistakes. Everything else
rests on the written arguments.

The mathematical dependency chain is short.
\begin{center}
\begin{tabularx}{\textwidth}{@{}l X@{}}
\toprule
Conclusion & Load-bearing ingredients\\
\midrule
Differential rigidity & Integer evaluation lemma; finite shift conversion;
escape chain and exclusion criterion.\\
Unit-dilation rigidity & Residue polynomial; finite binomial leading terms;
escape chain and exclusion criterion.\\
Noncofinal nonunit rigidity & Eventual order of finitely many affine values;
noncofinality witness; escape chain and exclusion criterion.\\
Coarsened exclusion & The same affine values, projected to
$\Gamma/\Delta_{\abs{v(q)}}$; strictness of the projected descent.\\
Positive dilation existence & Support monoid certificate and the eventual
strict increase of $h_n$; or, in the order-unit case, quadratic leading
valuations and the cofinal-valuation criterion.\\
Several-variable rigidity & Countable polynomial avoidance; finite-block
regrouping; rational differential system on a generic line.\\
First-order nonlinear rigidity & Hahn-family algebra and its residue
homomorphism; large active degrees; exact corner identity; uniqueness of a
corner term for each power of $Y'$.\\
Higher-order corner criterion & The same, with $\chi_P\ne0$ assumed.\\
Positive-weight Euler rigidity & The same algebra in $m$ variables; finite
bounded-weight sets; the eigenvalue identity of $\vartheta_q$.\\
\bottomrule
\end{tabularx}
\end{center}

All supports are sets. All operators are finite sums. All integer indices are
ordinary nonnegative integers. The coefficient valuation is trivial in
characteristic zero; the coefficient field is $\C$ in
Sections~\ref{hol:sec:intro}--\ref{hol:sec:certificates} and any
characteristic-zero field in the nonlinear part. The equations use the
external formal variable derivative. None of the proofs asserts convergence in the fine topology of
$\No[i]$, substitutes a real-valued norm for an arbitrary-rank valuation, or
transports a fixed-workspace evaluation to every surcomplex argument.

\section{Originality assessment and questions left open}
\label{hol:sec:status}

\subsection{Established background and proposed contribution}

The D-finite/P-recursive correspondence~\cite{stanley}, ordinary properties of
Hahn fields, the Hahn--Neumann support lemmas and Higman's finite-word
theorem~\cite{higman}, finite polynomial root bounds, the Conway normal-form
identification~\cite{bg,kks}, and partial theta series with their elementary
functional identity~\cite{andrewswarnaar} are established material. The
companion report's cofinality and scalar-extension work, including the group
$\Gamma_\infty$ and the distinction between order units and countable
cofinality, is prior project material~\cite{repoentire}. The rank-one
quadratic-exponent entire example is already in the
collection~\cite{reporankone}. We claim to originate none of these, and we do
not claim to have disproved any published theorem about a different analytic
function class.

The proposed research contribution is their combination into the sharp
arbitrary-rank dilation classification, Theorem~\ref{hol:main:q}: in
particular the necessity argument for all noncofinal scales, the unit case
with root-of-unity residue, the uniform exterior obstruction applying to every
solution of a fixed equation, the quantitative exclusion bounds with their two
different boundary behaviours, and the exact theta evaluation domain for an
arbitrary Hahn parameter. The mixed noncofinal extension
(Theorem~\ref{hol:thm:mixed}), generic-line rigidity
(Theorem~\ref{hol:main:multi}), and the explicit series excluded from every
nontorsion pure dilation equation (Theorem~\ref{hol:thm:explicit}) are further
consequences proved here.

The nonlinear part (source~10) inherits Gauss initial forms, Euler
eigenvalues and indicial-polynomial reasoning, and is related to
central-index, Newton-polygon, non-Archimedean value-distribution and tropical
differential methods (Section~\ref{hol:nl:sub:predecessors}). Its proposed
contribution is the precise arbitrary-rank strong-family formulation:
Theorem~\ref{hol:nl:main:first} with its finite exclusion certificate
(Theorem~\ref{hol:nl:thm:certificate}), the candidate degrees and the
entire-solution locus, the higher-order corner criterion, the positive-weight
Euler theorem, and the non-uniformity of the nonlinear exclusion bound. Its
priority is provisional in the sense recorded there.

\subsection{A result-by-result attribution ledger}

\begin{longtable}{@{}>{\raggedright\arraybackslash}p{0.40\linewidth}
                    >{\raggedright\arraybackslash}p{0.55\linewidth}@{}}
\toprule
Result or ingredient & Status in this report\\
\midrule
\endhead
Hahn arithmetic, strong support sums, surreal normal forms & Standard
background; references and support proofs supplied.\\[4pt]
Higman's finite-word theorem & External standard ingredient; alternatively
import the classical Hahn--Neumann support lemma directly.\\[4pt]
D-finite / polynomial recurrence conversion & Classical; credited to Stanley
and derived explicitly for Hahn coefficients.\\[4pt]
Partial theta series and functional identity & Classical; not a novelty
claim.\\[4pt]
Escape chain and exclusion criterion, with the local-finiteness boundary &
Elementary support-theoretic lemma and quantitative formulation developed
here.\\[4pt]
Eventual affine and periodic orbit valuations & Proved here for arbitrary Hahn
coefficients; the finite residue-polynomial treatment covers near-torsion
units.\\[4pt]
Exact fixed-$q$ order-unit classification & Principal proposed contribution:
arbitrary rank, arbitrary Hahn parameter, both signs of $v(q)$.\\[4pt]
Exact theta evaluation domain & Support-controlled sharpness theorem; related
in spirit to the companion report's scalar-extension results, proved
independently here.\\[4pt]
Uniform exterior obstruction and common degree bound & Proved here; stronger
than eventual termination of each solution separately.\\[4pt]
Universal single-dilation nonholonomicity on no-order-unit workspaces &
Consequence of the classification, with an explicit entire witness.\\[4pt]
Mixed differential--dilation rigidity at noncofinal nonzero scales & Proved
here; the unit-valuation mixed case is left open.\\[4pt]
Several-variable rigidity and multivariate algebraic rigidity & Proved here;
the generic line is an existence device.\\[4pt]
Gauss initial forms, Euler eigenvalues, indicial polynomials, central-index
and Newton-polygon ideas & Established (source~10); credited to Hayman, Cano
and Fortuny Ayuso, Hu--Luan, Hu--Yang and Giansiracusa--Mereta as
predecessors in different settings.\\[4pt]
First-order nonlinear rigidity (Theorem~\ref{hol:nl:main:first}), finite
exclusion certificate, candidate degrees and entire-solution locus &
Proposed contribution of source~10, over any characteristic-zero coefficient
field; priority provisional.\\[4pt]
Higher-order corner criterion and positive-weight Euler rigidity & Proposed
by source~10; sufficient conditions only; the vanishing-corner case is
Question~\ref{hol:q:nonlinear}.\\[4pt]
Solution-dependence of the nonlinear exclusion bound & Proved by source~10
with the explicit Riccati family of
Theorem~\ref{hol:nl:thm:nouniform}.\\[4pt]
Two-jet independence of $\sum t^{n^2}z^n$ and of
\eqref{hol:eq:explicitF} & New consequences (source~10); both series are
prior collection material.\\[4pt]
Named published open conjecture & None claimed solved; the nonlinear question
answered in order one is a question of this report, not a historical
conjecture.\\[4pt]
Priority, peer review, formal verification & Priority and independent
refereeing are not certified. The support word lemma, escape mechanism and
formal exponential domain have Lean coverage; the main classifications and
the whole nonlinear part remain unformalized.\\
\bottomrule
\end{longtable}

The targeted literature checks of sources~08 and~09 located the standard
D-finite framework,
classical complex $q$-difference growth theory, $q$-holonomic degree growth in
a rational-function setting, and ultrametric $q$-difference theory, but did
not establish whether an identical arbitrary-rank strong-Hahn formulation
appears elsewhere. Several broad keyword searches returned irrelevant material
and were not treated as evidence of absence; a repository content search for
\texttt{holonomic} returning no matches is likewise not proof that no sentence
anywhere in the repository can be reformulated into one of these results.
Accordingly this is a proposed original result with proofs, not a certified
claim of bibliographic priority. Independent mathematical review is still
needed; machine-checked scope is the partial coverage recorded in
Section~\ref{hol:sec:certificates} and the formalization ledger. The
corresponding record for the nonlinear part, with its own search limits, is
Section~\ref{hol:nl:sub:predecessors}.

\subsection{What the nonlinear part does not claim}
\label{hol:nl:sub:nonclaims}

Every limitation stated by source~10, in its manuscript or its two audit
files, is kept; they are collected here.
\begin{enumerate}[label=\textup{(N\arabic*)}]
\item Only the first-order case of Question~\ref{hol:q:nonlinear} is
answered, not the unrestricted higher-order question.
\item Theorem~\ref{hol:nl:main:first} does not say that a nonpolynomial
entire function is differentially transcendental in all orders; independence
of $f$ and $D_zf$ does not exclude a relation involving $D_z^2f$.
\item The criterion $\chi_P\ne0$ is sufficient, not a complete invariant of a
solution set, and the methods do not decide equations of order at least two
with vanishing residue corner. Proposition~\ref{hol:nl:prop:degenerate} is not
a nonpolynomial entire solution and does not refute higher-order rigidity.
\item The corner is a finite selection, not a full differential Newton
polygon; no invention of differential Newton polygons, no equivalence to
tropical differential algebra, and no new tropical fundamental theorem is
claimed.
\item The exclusion bound \eqref{hol:nl:eq:deltazero} is not claimed optimal;
the certificate has no converse; leading valuations alone do not decide
evaluation at a boundary argument.
\item Candidate degree sets are necessary, not realized; their finiteness is
not a uniform algorithm for arbitrary coefficient names, and the workflow of
Section~\ref{hol:nl:sub:workflow} is not an algorithm for arbitrary surreal
names, does not enumerate local formal solutions, does not decide
strong-support membership, and accepts no numerical solver failure as a
certificate.
\item Corollary~\ref{hol:nl:cor:locus} does not say that the number of
solutions is finite, is not a moduli statement over rings with nilpotents,
and covers only the extension class it states.
\item No theorem says that $G$, $F$ or any nonpolynomial series is entire on
the proper class $\No[i]$; entireness and the derivative are relative to one
set-sized workspace.
\item The independence statements concern formal functions; no statement
concerns algebraic independence of a value $f(x)$ from its own ambient field.
\item $D_z$ kills all Hahn coefficients and is not the Berarducci--Mantova
derivation; no scalar-derivation result is claimed.
\item The series $\sum t^{n^2}z^n$ and \eqref{hol:eq:explicitF}, and the
linear rigidity of the latter, are prior material; only their nonlinear
two-jet consequences are new. The tail-span report's all-order independence
results are neither weakened nor rediscovered.
\item Hahn fields, the positive-support summation lemma, Conway normal forms,
Gauss initial forms, Euler eigenvalues and indicial-polynomial reasoning are
established ideas, not relabelled as new; the positive-support lemma is
imported by source~10, and proved in this report by
Lemma~\ref{hol:lem:support}\,(2).
\item The precise all-rank statements are proposed contributions; historical
priority is not certified. The literature search of 22 September 2026 was
targeted and not exhaustive; Hayman's survey and the Hu--Yang monograph were
consulted through metadata only, and no claim is made that rank-one special
cases were previously unknown.
\item The repository comparison at \texttt{048b72c} was targeted: no
line-by-line audit of every report, Lean file, archive, branch or the Git
history, and no absence claim rests on a truncated preview.
\item The answered question is a question of this report, not a named
historical conjecture.
\item No Lean formalization and no independent refereeing of the nonlinear
part exists. The 2{,}113 finite checks are program checks, not theorems; they
do not test infinite supports, cofinality, all Hahn elements, nonexistence of
annihilating equations or novelty, and no infinite surreal sum is numerically
approximated. A successful PDF build is not a proof check.
\item The sequence arguments behind the support calculus use ordinary
set-theoretic choice.
\end{enumerate}

\subsection{Further mathematical questions}

\begin{question}[Nonlinear differential equations of order at least two]
\label{hol:q:nonlinear}
Can a nonpolynomial $f\in\E$ satisfy a nonzero algebraic differential
equation $P(z,f,D_zf,\ldots,D_z^sf)=0$ over $\K(z)$ of order $s\ge2$? By
Theorems~\ref{hol:nl:main:first} and~\ref{hol:nl:thm:corner}, such an $f$
satisfies no nonzero equation of order at most one and no equation with
$\chi_P\ne0$; equivalently, can a nonpolynomial strongly entire function
satisfy an equation whose residue corner polynomial vanishes identically,
without satisfying any other equation to which
Theorem~\ref{hol:nl:thm:certificate} applies? Under which hypotheses on
$\Gamma$ can linear rigidity be strengthened to differential algebraic
rigidity in all orders? Does the explicit series \eqref{hol:eq:explicitF}
avoid all equations of order at least two under a natural coefficient
restriction?
\end{question}

\paragraph{Status.}
As posed by sources~08 and~09, the question asked about nonlinear algebraic
differential equations of every order. Source~10 answers its first-order case
negatively, for every nonzero $\Gamma$ and every characteristic-zero
coefficient field (Theorem~\ref{hol:nl:main:first}): in order one no
hypothesis on $\Gamma$ beyond $\Gamma\ne0$ is needed, and
\eqref{hol:eq:explicitF} avoids every first-order equation with no
coefficient restriction (Corollary~\ref{hol:nl:cor:infiniterank}). It also
settles the higher-order equations with $\chi_P\ne0$
(Theorem~\ref{hol:nl:thm:corner}) and the relations between $f$ and a
positive-weight Euler derivative (Theorem~\ref{hol:nl:thm:euler}). What stays
open is order at least two with a vanishing residue corner.
Proposition~\ref{hol:nl:prop:degenerate} shows that such equations exist and
can have polynomial solutions of unbounded degree, so the study cannot rely on
a uniform degree bound for all polynomial solutions of an arbitrary
higher-order equation; it is not a nonpolynomial entire solution and gives no
negative answer. Source~10 suggests as a next step classifying the polynomial
solutions of a vanishing-corner residue equation beyond their leading monomial
and determining how those constraints propagate between scales.

The escape proof does not answer this, in any order. A nonlinear differential
relation produces convolution sums in the coefficients rather than a
bounded-step linear recurrence, and treating it as though it had the form
\eqref{hol:eq:generalrecurrence} would be an invalid extension. The
first-order answer uses the different route of
Section~\ref{hol:nl:sec:setting}.

\begin{question}[Beyond a positive Euler operator]\label{hol:nl:q:vectorfields}
Which polynomial vector fields $V$ have the property that every strongly
entire $f$ satisfying a nonzero equation $P(x,f,Vf)=0$ is a polynomial?
\end{question}

Positive integer Euler weights give one family, and zero and mixed weights
give explicit counterexamples (Section~\ref{hol:nl:sub:euler}). A useful
extension would need both a finite-dimensional filtration and a sufficiently
nondegenerate action on its highest layers; diagonalization alone is not
enough if infinitely many monomials share a bounded eigenvalue. (Source~10.)

\begin{question}[Effective representations and minimal certificates]
\label{hol:nl:q:effective}
For which exact Hahn coefficient representations can the candidate degree set
and an optimal excluded valuation cut be computed uniformly?
\end{question}

The bound \eqref{hol:nl:eq:deltazero} is always a valid choice from exact
finite valuation data but can be far from minimal, and
Theorem~\ref{hol:nl:thm:nouniform} shows that an optimal cut must in general
depend on solution data. An algorithmic refinement must state how it receives
those data and how it decides coefficient cancellation. (Source~10; the
linear analogue is Question~\ref{hol:q:optimal}.)

\begin{question}[Mixed operators at unit valuation]\label{hol:q:mixedunit}
For a nontorsion $q$ with $v(q)=0$, must every strongly entire solution of a
nonzero mixed equation \eqref{hol:eq:mixed} be a polynomial?
\end{question}

The unresolved point in the present method is uniform control of
$v(P(n,q^n))$ for a general bivariate $P$ in the unit case. The pure
$q$-difference proof controls $P(q^n)$ and the differential proof controls
$P(n)$; neither statement by itself handles their combination. This is a
question suggested by the argument, not asserted to be a previously
catalogued open problem.

\begin{question}[Optimal exterior regions and exact solution domains]
\label{hol:q:optimal}
For a fixed differential or rigid dilation operator, can one compute the
largest common strong evaluation domain of its formal solution space from
finite valued operator data, rather than only a sufficient exclusion bound?
Can the coarse exclusion bound of Theorem~\ref{hol:thm:coarse} be supplemented
by support hypotheses that give the exact domain?
\end{question}

The thresholds supplied here are uniform and explicit once the valuation
certificates are available, but need not be optimal, and distinct solutions may
have cancellations or terminating tails that change their individual domains.
Proposition~\ref{hol:prop:expdomain} shows that the uncoarsened bound can be
exactly attained, and Theorem~\ref{hol:thm:thetadomain} exhibits one
sufficient pattern for an exact answer: a common positive support monoid
together with a controlled sequence of leading exponents. A full answer should
retain arbitrary Hahn tails rather than use only a numerical radius.

\begin{question}[Several independent dilations]\label{hol:q:severaldilations}
What is the existence criterion for a nonpolynomial strongly entire solution
constrained by equations using several multiplicatively independent dilations?
Is an appropriate cofinality condition on their joint valuation data
sufficient, or do compatibility relations impose additional obstructions? For
a no-order-unit workspace, can one classify the finite linear relations among
$F(q_1z),\ldots,F(q_sz)$ when the parameters are not powers of a common
element?
\end{question}

The single-dilation theorem identifies the exact obstruction when only one
scale is repeatedly added. With several dilation directions, distinct
monomials can have equal valuation slopes and their leading coefficients can
produce nontrivial recurrence phenomena of their own. A meaningful formulation
must also quotient out identities caused by repeated parameters and any
genuine algebraic relations among them. No multidilation classification is
implied by Theorem~\ref{hol:main:q}.

These questions are directions suggested by the proved statements, not claims
that they were already posed in the published literature, and their
formulations are deliberately narrower than a general theory of nonlinear
surcomplex differential equations.

\section{Conclusion}

Strong entireness over a surreal Hahn workspace imposes a global scale
requirement that a finite linear differential equation cannot meet unless its
solution terminates. A dilation equation can meet that requirement, but
exactly when its valuation reaches every scale through integer multiples. The
decisive distinction is therefore not between Archimedean and non-Archimedean
fields, nor between rank one and higher rank; it is between \emph{a cofinal
dilation scale} and \emph{a scale trapped below another one}.

The obstruction is finite and robust. A linear recurrence forces a later
nonzero coefficient with a controlled valuation, and a large enough argument
turns an infinite chain of such coefficients into a nonincreasing sequence of
leading exponents. Descending supports and infinite reuse of one exponent are
two distinct ways in which the evaluation fails, and keeping them separate is
what makes one exclusion boundary closed and the other open. The mechanism
works with arbitrary Hahn coefficient tails, arbitrarily high valuation rank,
and sparse coefficient sequences, without replacing strong summability by
topological convergence and without any hidden completeness assumption.

The countable-rank example makes the separation concrete. The series
$\sum_{n\ge0}\omega^{-\omega^n}z^n$ is evaluable at every point of its full
workspace, yet no single nontorsion dilation in that workspace supplies a
finite linear equation for it. That gap between existence and finite linear
describability is the main mathematical contribution proposed here.

The nonlinear part pushes the same global requirement past linear
equations by a different mechanism. A finite algebraic relation between a
function and its first derivative is incompatible with nonpolynomial strong
entireness over any nontrivial characteristic-zero Hahn workspace: at one
large admissible scale the whole series leaves a finite initial polynomial,
whose top coefficient is controlled by a corner polynomial that is
automatically nonzero in order one. In higher order an exact cancellation of
that corner is a genuine obstruction to the method, and a fixed nonlinear
equation can then have polynomial solutions of unbounded degree; whether it
can have a nonpolynomial entire solution is the open part of
Question~\ref{hol:q:nonlinear}.

\appendix

\section{Notation}\label{hol:app:notation}

\begin{center}
\begin{tabularx}{\textwidth}{@{}l X@{}}
\toprule
Symbol & Meaning\\
\midrule
$\Gamma$ & A fixed nonzero set-sized ordered abelian group; divisible only
where stated.\\
$\K$ & The complex Hahn field $\C((t^\Gamma))$.\\
$v(a),\lc(a),\res(a)$ & Least support exponent, its coefficient, and the
residue when $v(a)=0$.\\
$\E$, $\Dom(f)$ & Series strongly evaluable at every point of $\K$; the exact
set where a series is strongly evaluable.\\
$D_z$ & Formal derivative in $z$, zero on every element of $\K$.\\
$\sigma_q$ & Argument dilation $f(z)\mapsto f(qz)$.\\
$M(S)$ & The additive monoid generated by a well-ordered $S\subseteq\Gamma_{>0}$.\\
$\Delta_\beta$ & The smallest convex subgroup containing $\beta>0$,
\eqref{hol:eq:delta}.\\
$B$, $B^{+}$, $B_{\mathrm{rec}}$ & Exclusion thresholds in $\Gamma$;
$B_{\mathrm{rec}}$ divides by the jump length and needs $\Gamma$ divisible.\\
$\overline B^{\flat}_{\mathrm{rec}}$,
$\overline B_{\mathrm{rec}}$ & Their coarsened analogues in
$\Gamma/\Delta_{\abs{v(q)}}$.\\
$\mu$, $\beta$ & Usually a positive dilation valuation; an order unit only
when stated.\\
$\fall Xr$ & The ordinary falling factorial polynomial of degree $r$.\\
$\Theta_p$ & The partial theta series \eqref{hol:eq:theta}.\\
$\Gamma_\infty$ & The increasing-rank direct sum \eqref{hol:eq:gammainfty}.\\
\midrule
\multicolumn{2}{@{}l}{\emph{Nonlinear part,
Sections~\ref{hol:nl:sec:setting}--\ref{hol:nl:sec:consequences}}}\\
$k$, $\KK$, $\EK$ & Any characteristic-zero field; $k((t^\Gamma))$; its
strongly entire series (Convention~\ref{hol:nl:conv:field}).\\
$\Sfam_m$, $\Sfam_m^+$, $\rho$ & Series with strongly summable coefficient
family; those with coefficients of nonnegative valuation; coefficientwise
residue.\\
$\gamma_\delta$, $\phi_\delta$, $N_\delta$ & Gauss value, initial polynomial
and top active degree at the scale $t^{-\delta}$,
\eqref{hol:nl:eq:gaussdata}.\\
$d_P$, $w_P$, $\mathcal C_P$, $\beta_P$ & Corner data of a differential
polynomial, \eqref{hol:nl:eq:dw}.\\
$\chi_P$, $I_P$ & Residue and full corner polynomials,
\eqref{hol:nl:eq:chiI}.\\
$\kappa_P$ & The rational degree cutoff \eqref{hol:nl:eq:kappa}.\\
$\delta_0$ & A solution-dependent exclusion bound,
\eqref{hol:nl:eq:deltazero}.\\
$\vartheta_q$, $\wt$ & Positive-weight Euler operator and weight.\\
\bottomrule
\end{tabularx}
\end{center}

\section{Proof-dependency and hypothesis audit}\label{hol:app:audit}

The logical dependencies are short. Differential rigidity uses the
integer-evaluation lemma, the coefficient conversion, and the escape chain
with its exclusion criterion. Dilation necessity uses the geometric-orbit
lemma or the unit-orbit theorem, the dilation conversion, and the same
criterion. Dilation sufficiency uses only the partial theta identity and the
support certificate. The mixed extension uses the bivariate orbit lemma. The
several-variable theorem uses countable polynomial avoidance, finite-block
regrouping and the one-variable theorem. The infinite-rank application uses a
direct strong-entireness proof and the completed classification. The
nonlinear part uses none of these linear arguments: first-order rigidity uses
the Hahn-family algebra (with the support calculus of
Section~\ref{hol:sec:support}), the large-active-degree lemma, the corner
identity and the order-one nondegeneracy of the corner; the higher-order
criterion replaces the last step by the hypothesis $\chi_P\ne0$; the Euler
theorem uses the same algebra in several variables with finite bounded-weight
sets.

No unproved ring-theoretic conclusion from another report in this collection
is used. The real-valued coarsening of~\cite[Lemma~9.1]{repoentire} is cited
and not reproved, and no proof here depends on it. Higman's theorem is an
external standard ingredient; alternatively one may import the classical
Hahn--Neumann support lemma directly. The surreal interpretation imports the
normal-form field embedding, while every substantive function statement is
already proved in the abstract Hahn field.

The potentially dangerous hypotheses, collected for review:
\begin{itemize}
\item Supports and the value group are sets; all operators are finite sums;
all integer indices are ordinary nonnegative integers.
\item $\Gamma\ne0$, set-sized, ordered abelian. \emph{Divisibility is not
assumed} except in Corollary~\ref{hol:cor:periodic}, the refined threshold
\eqref{hol:eq:refinedthreshold}, Theorem~\ref{hol:thm:coarse}\,(b),
and Example~\ref{hol:ex:sparse}, each of which states it.
\item Series are ordinary $\N$-indexed power series in one variable, except in
Section~\ref{hol:sec:multi}. Equation coefficients are rational functions of
that variable, not arbitrary entire functions; their scalar coefficients may
be arbitrary Hahn elements.
\item The residue field is characteristic zero and trivially valued.
Algebraic closedness is not used. It is $\C$ in
Sections~\ref{hol:sec:intro}--\ref{hol:sec:certificates} and an arbitrary
characteristic-zero field in the nonlinear part; the latter generality is not
transferred to the linear theorems (Convention~\ref{hol:nl:conv:field}).
Characteristic zero is essential (Example~\ref{hol:nl:ex:charp}).
\item Dilation parameters are nonzero and nontorsion; nontorsion is
essential, as Section~\ref{hol:sec:torsion} shows by counterexample.
\item The derivative is $d/dz$, not a scalar surreal derivation.
\item Entireness is always relative to one fixed workspace, never to
$\No[i]$. Strong summability requires both well-ordering and local finiteness.
\item The conclusions of
Sections~\ref{hol:sec:intro}--\ref{hol:sec:certificates} concern
\emph{linear} equations. Nonlinear differential equations are treated only
in order one, in higher order under $\chi_P\ne0$, and for positive-weight
Euler relations (the nonlinear part); no nonlinear differential
transcendence in all orders, no nonlinear difference transcendence and no
unrestricted mixed-dilation independence is claimed.
\end{itemize}

Two boundary behaviours deserve separate emphasis because they point in
opposite directions. The most delicate equality is in
Corollary~\ref{hol:cor:periodic}: weakly decreasing leading exponents already
obstruct strong summability, so the excluded region is closed, and
Proposition~\ref{hol:prop:expdomain} shows the bound is attained. The most
delicate strictness is in Theorem~\ref{hol:thm:coarse}: equal \emph{coarse}
values need not obstruct summability in $\Gamma$, so that excluded region must
be open, and Theorem~\ref{hol:thm:thetadomain} with
Section~\ref{hol:sub:coarseboundary} shows why.

Deliberate exclusions, recorded so that no reader supplies them by analogy:
finite-order $q$ is excluded from the main classification and has genuine
periodic-support counterexamples; mixed operators at unit valuation are not
classified; no nonlinear differential algebraic rigidity is claimed beyond
order one, the nondegenerate higher-order class and positive-weight Euler
relations, whose own non-claims are listed in
Section~\ref{hol:nl:sub:nonclaims}; no complete algorithm for arbitrary
surreal comparisons or cofinality decisions is claimed; the finite checks
prove no infinite theorem. The main theorem package, the nonlinear part
included, has no Lean verification or independent refereeing; the existing
support, escape and exponential-domain coverage is recorded in the
formalization ledger.

\section{Reproducing the document and the finite checks}
\label{hol:app:reproduce}

The directory contains this standalone \LaTeX{} source, its compiled PDF, a
README, the three source verification programs and source~10's build script
under \texttt{code/}, the three recorded outputs and the source build and
layout reports under \texttt{data/}, and five source audit documents. No
source document from the inspected repository is redistributed, and no source
manuscript is shipped. The bibliography is embedded: no BibTeX run, external
graphics, shell escape, or repository checkout is required, and a conventional
TeX Live or MiKTeX installation with the named packages suffices.

\begin{verbatim}
latexmk -pdf -interaction=nonstopmode -halt-on-error article.tex
python code/08-finite-recurrences-verify.py --output rerun-08.txt
python code/09-entire-hahn-holonomic-verification.py --output rerun-09.txt
python code/10-nonlinear-rigidity-verify.py --output rerun-10.json
\end{verbatim}

The \verb|--output| argument is required by the first and third scripts, so
that a rerun need not overwrite the record distributed with the report, and
optional for the second; the third refuses an existing output path. All use
Python 3.9 or later, exit with a nonzero status if any check fails, and use
exact integer and rational arithmetic throughout. The build script
\path{code/10-nonlinear-rigidity-build.py} was written for source~10's own
manuscript: run here it would typeset this report, leave auxiliary files in
the directory and write \texttt{data/build\_report.json}. The shipped
\path{data/10-nonlinear-rigidity-build_report.json} and
\path{data/10-nonlinear-rigidity-layout_audit.json} describe the 23-page
source manuscript, not this report; use \texttt{latexmk} as above, on a copy.

\clearpage
\begin{thebibliography}{99}
\addcontentsline{toc}{section}{References}
\raggedright

\bibitem{stanley}
R. P. Stanley,
\emph{Differentiably finite power series},
European Journal of Combinatorics \textbf{1} (1980), no.~2, 175--188.
Author-hosted copy:
\url{https://math.mit.edu/~rstan/pubs/pubfiles/45.pdf}.
Theorem~1.5 (printed page 176) gives the D-finite/P-recursive equivalence.
The finite coefficient conversion is reproved in this report for its stated
characteristic-zero Hahn coefficient field.

\bibitem{ramis}
J.-P. Ramis,
\emph{About the growth of entire functions solutions of linear algebraic
$q$-difference equations},
Annales de la Facult\'e des Sciences de Toulouse, Math\'ematiques,
S\'erie~6, \textbf{1} (1992), no.~1, 53--94.
\url{https://www.numdam.org/item/AFST_1992_6_1_1_53_0/}.

\bibitem{garoufalidis}
S. Garoufalidis,
\emph{The degree of a $q$-holonomic sequence is a quadratic quasi-polynomial},
Electronic Journal of Combinatorics \textbf{18} (2011), no.~2, Paper~P4.
Inspected manuscript: arXiv:1005.4580v4.
\url{https://arxiv.org/abs/1005.4580}.

\bibitem{divizio}
L. Di Vizio,
\emph{An ultrametric version of the Maillet--Malgrange theorem for nonlinear
$q$-difference equations},
Proceedings of the American Mathematical Society \textbf{136} (2008), no.~8,
2803--2814. Inspected manuscript: arXiv:0709.2464v2.
\url{https://arxiv.org/abs/0709.2464}.

\bibitem{andrewswarnaar}
G. E. Andrews and S. O. Warnaar,
\emph{The product of partial theta functions}.
Author-hosted manuscript; the definition on its first page is the series used
here, up to a sign convention.
\url{https://people.smp.uq.edu.au/OleWarnaar/pubs/Partial-thetas.pdf}.

\bibitem{higman}
G. Higman,
\emph{Ordering by divisibility in abstract algebras},
Proceedings of the London Mathematical Society (3) \textbf{2} (1952),
326--336.
\url{https://doi.org/10.1112/plms/s3-2.1.326}.

\bibitem{mueller}
J. M\"uller and A. Strohmaier,
\emph{The theory of Hahn meromorphic functions, a holomorphic Fredholm theorem
and its applications},
Analysis \& PDE \textbf{7} (2014), no.~3, 745--770.
\url{https://arxiv.org/abs/1205.0236}.
\url{https://doi.org/10.2140/apde.2014.7.745}.

\bibitem{bm}
A. Berarducci and V. Mantova,
\emph{Surreal numbers, derivations and transseries},
Journal of the European Mathematical Society \textbf{20} (2018), 339--390.
\url{https://arxiv.org/abs/1503.00315}.
\url{https://doi.org/10.4171/JEMS/769}.

\bibitem{bg}
O. Bournez and Q. Guilmant,
\emph{Surreal fields stable under exponential, logarithmic, derivative and
anti-derivative functions},
arXiv:2211.08396 (2022).
\url{https://arxiv.org/abs/2211.08396}.

\bibitem{kks}
E. Kaplan, L. S. Krapp and M. Serra,
\emph{Decomposing the automorphism group of the surreal numbers},
arXiv:2509.22374v3.
Used for the Hahn-field and surreal normal-form conventions, not for its
automorphism results.
\url{https://arxiv.org/html/2509.22374v3}.

\bibitem{repo}
\repo\ project,
\emph{Repository documentation and research-report catalogue},
snapshot \texttt{a3124af79f66b8b9c196d76b4cbc5ac3938907c4},
consulted 22 September 2026.
\url{https://github.com/VladimirReshetnikov/Surreal}.
The inspection used the root README and \texttt{docs/README.md} at this
commit; these are project documentation, not independent published
validation.

\bibitem{repoentire}
\repo\ project,
\emph{Cofinality, factorization and scalar extension for entire Hahn functions
at arbitrary rank},
\texttt{docs/surcomplex/entire-functions-at-arbitrary-rank/}, same pinned
snapshot.
\url{https://github.com/VladimirReshetnikov/Surreal/tree/a3124af79f66b8b9c196d76b4cbc5ac3938907c4/docs/surcomplex/entire-functions-at-arbitrary-rank}.
Its Lemma~9.1, labelled \texttt{ent:lem:coarsening}, is the real-valued
coarsening cited in Remark~\ref{hol:rem:coarsening}; its Theorem~9.3 is the
B\'ezout statement referred to in Section~\ref{hol:sub:position}.

\bibitem{repodifferential}
\repo\ project,
\emph{Differential algebra and differential equations over the surreal and
surcomplex numbers},
\texttt{docs/surcomplex/differential-equations/}, same pinned snapshot.
\url{https://github.com/VladimirReshetnikov/Surreal/tree/a3124af79f66b8b9c196d76b4cbc5ac3938907c4/docs/surcomplex/differential-equations}.
Cited to distinguish the scalar surreal derivation from $D_z$.

\bibitem{reporankone}
\repo\ project,
\emph{Rank-one non-Archimedean analytic geometry inside the surcomplex
numbers},
\texttt{docs/surcomplex/rank-one-berkovich/}, same pinned snapshot.
\url{https://github.com/VladimirReshetnikov/Surreal/tree/a3124af79f66b8b9c196d76b4cbc5ac3938907c4/docs/surcomplex/rank-one-berkovich}.
Cited for the pre-existing quadratic-exponent rank-one example.

\bibitem{repofoundations}
\repo\ project,
\emph{Foundations for Surreal and Surcomplex Mathematics},
\texttt{docs/foundations-and-computation/foundations/}.
Its Example \texttt{found:ex:internal} is the series $\sum t^{n^2}z^n$ over
$\C((t^\Q))$; cited by this merge, not by source~10.

\bibitem{repotail}
\repo\ project,
\emph{Tail-spans and differential transcendence in surreal and surcomplex Hahn
fields},
\texttt{docs/surreal/tail-spans-and-differential-transcendence/}, at
source~10's pin \texttt{048b72cf7cbfc8ab246e4f73788c10460cb3f6e0}.
\url{https://github.com/VladimirReshetnikov/Surreal/tree/048b72cf7cbfc8ab246e4f73788c10460cb3f6e0/docs/surreal/tail-spans-and-differential-transcendence}.
Its README was inspected by source~10 for function-class and derivative
boundaries.

\bibitem{mantovamatusinski}
V. Mantova and M. Matusinski,
\emph{Surreal numbers with derivation, Hardy fields and transseries: a
survey}, arXiv:1608.03413, inspected version~2, Section~2 on normal forms.
\url{https://arxiv.org/html/1608.03413v2}.

\bibitem{hayman}
W. K. Hayman,
\emph{The local growth of power series: a survey of the Wiman--Valiron
method},
Canadian Mathematical Bulletin \textbf{17} (1974), no.~3, 317--358.
\url{https://doi.org/10.4153/CMB-1974-064-0}.
Publisher metadata and abstract consulted by source~10.

\bibitem{cano}
J. Cano and P. Fortuny Ayuso,
\emph{Power series solutions of non-linear $q$-difference equations and the
Newton--Puiseux polygon},
Qualitative Theory of Dynamical Systems \textbf{21} (2022), article~123.
\url{https://doi.org/10.1007/s12346-022-00656-0}.
Introduction and Section~2 (initial and indicial polynomials) consulted by
source~10.

\bibitem{huluan}
P.-C. Hu and Y.-Z. Luan,
\emph{Non-Archimedean meromorphic solutions of functional equations},
arXiv:1311.5291v1 (2013).
\url{https://arxiv.org/abs/1311.5291}.
Field hypotheses and Theorems~1.1, 1.3 and~1.5 consulted by source~10.

\bibitem{huyang}
P.-C. Hu and C.-C. Yang,
\emph{Meromorphic Functions over Non-Archimedean Fields},
Mathematics and Its Applications, Springer, Dordrecht, 2000.
\url{https://doi.org/10.1007/978-94-015-9415-8}.
Publisher metadata and description only; the book was not audited.

\bibitem{giansiracusa}
J. Giansiracusa and S. Mereta,
\emph{A general framework for tropical differential equations},
manuscripta mathematica \textbf{173} (2024), 1273--1304.
\url{https://doi.org/10.1007/s00229-023-01492-5}.

\end{thebibliography}
\end{document}
